\input{preamble}

% OK, start here.
%
\begin{document}

\title{Propriétés des schémas}

\maketitle

\phantomsection
\label{section-phantom}

\tableofcontents

\section{Introduction}
\label{section-introduction}

\noindent
Dans ce chapitre, nous introduisons quelques propriétés absolues des schémas.
Une référence fondamentale est \cite{EGA}.

\section{Ensembles constructibles}
\label{section-constructible}

\noindent
Les ensembles constructibles et localement constructibles sont introduits dans
Topologie, section \ref{topology-section-constructible}.
Les sous-ensembles localement constructibles des schémas se caractérisent comme suit.

\begin{lemma}
\label{lemma-locally-constructible}
Soit $X$ un schéma.
Un sous-ensemble $E$ de $X$ est localement constructible dans $X$ si et seulement si
$E \cap U$ est constructible dans $U$ pour tout ouvert affine $U$ de $X$.
\end{lemma}

\begin{proof}
Supposons $E$ localement constructible. Il existe alors un recouvrement ouvert
$X = \bigcup U_i$ tel que $E \cap U_i$ soit constructible dans $U_i$
pour tout $i$. Soit $V \subset X$ un ouvert affine quelconque. On peut trouver un
recouvrement ouvert affine fini $V = V_1 \cup \ldots \cup V_m$ tel que, pour tout $j$,
on ait $V_j \subset U_i$ pour un certain $i = i(j)$. D'après
Topologie, lemme \ref{topology-lemma-open-immersion-constructible-inverse-image},
chaque $E \cap V_j$ est constructible dans $V_j$. Les inclusions
$V_j \to V$ étant quasi-compactes (voir
Schémas, lemme \ref{schemes-lemma-quasi-compact-affine}),
on en déduit que $E \cap V$ est constructible dans $V$ par
Topologie, lemme \ref{topology-lemma-collate-constructible}.
La réciproque est immédiate.
\end{proof}

\begin{lemma}
\label{lemma-generic-point-in-constructible}
Soient $X$ un schéma et $E \subset X$ un sous-ensemble localement constructible.
Soit $\xi \in X$ un point générique d'une composante irréductible de $X$.
\begin{enumerate}
\item Si $\xi \in E$, alors un voisinage ouvert de
$\xi$ est contenu dans $E$.
\item Si $\xi \not \in E$, alors un voisinage ouvert
de $\xi$ est disjoint de $E$.
\end{enumerate}
\end{lemma}

\begin{proof}
Le complémentaire d'un sous-ensemble localement constructible étant localement
constructible, il suffit de montrer (2). On peut supposer $X$
affine, et donc $E$ constructible (lemme \ref{lemma-locally-constructible}).
Dans ce cas, $X$ est un espace spectral
(Algèbre, lemme \ref{algebra-lemma-spec-spectral}).
Alors $\xi \not \in E$ entraîne $\xi \not \in \overline{E}$ d'après
Topologie, lemme \ref{topology-lemma-constructible-stable-specialization-closed},
car aucun point de $X$ distinct de $\xi$
ne se spécialise en $\xi$.
\end{proof}

\begin{lemma}
\label{lemma-quasi-separated-quasi-compact-open-retrocompact}
Soit $X$ un schéma quasi-séparé. L'intersection de deux
ouverts quasi-compacts de $X$ est un ouvert quasi-compact de $X$.
Tout ouvert quasi-compact de $X$ est rétrocompact dans $X$.
\end{lemma}

\begin{proof}
Si $U$ et $V$ sont des ouverts quasi-compacts, alors
$U \cap V = \Delta^{-1}(U \times V)$, où $\Delta : X \to X \times X$
est la diagonale. Comme $X$ est quasi-séparé, $\Delta$ est
quasi-compact. Ainsi $U \cap V$ est quasi-compact, puisque
$U \times V$ l'est (détails omis ; utiliser
Schémas, lemme \ref{schemes-lemma-affine-covering-fibre-product},
pour voir que $U \times V$ est une réunion finie d'ouverts affines).
Les autres assertions résultent de la première et de
Topologie, lemme \ref{topology-lemma-topology-quasi-separated-scheme}.
\end{proof}

\begin{lemma}
\label{lemma-quasi-compact-quasi-separated-spectral}
Soit $X$ un schéma quasi-compact et quasi-séparé.
Alors l'espace topologique sous-jacent à $X$ est spectral.
\end{lemma}

\begin{proof}
D'après Topologie, définition \ref{topology-definition-spectral-space},
il faut vérifier que $X$ est sobre, quasi-compact, admet une base
d'ouverts quasi-compacts, et que l'intersection de deux
ouverts quasi-compacts est quasi-compacte. Cela résulte de
Schémas, lemmes \ref{schemes-lemma-scheme-sober} et
\ref{schemes-lemma-basis-affine-opens},
ainsi que du
lemme \ref{lemma-quasi-separated-quasi-compact-open-retrocompact}
ci-dessus.
\end{proof}

\begin{lemma}
\label{lemma-constructible-quasi-compact-quasi-separated}
Soit $X$ un schéma quasi-compact et quasi-séparé.
Tout sous-ensemble localement constructible de $X$ est constructible.
\end{lemma}

\begin{proof}
Comme $X$ est quasi-compact, on peut choisir un recouvrement ouvert affine fini
$X = V_1 \cup \ldots \cup V_m$. Comme $X$ est quasi-séparé, chaque $V_i$ est
rétrocompact dans $X$ d'après le
lemme \ref{lemma-quasi-separated-quasi-compact-open-retrocompact}.
Ainsi, par
Topologie, lemme \ref{topology-lemma-collate-constructible},
un sous-ensemble $E \subset X$ est constructible dans $X$ si et seulement si
$E \cap V_j$ est constructible dans $V_j$. On conclut par le
lemme \ref{lemma-locally-constructible}.
\end{proof}

\begin{lemma}
\label{lemma-retrocompact}
Soit $X$ un schéma. Un sous-ensemble $E$ de $X$ est rétrocompact dans $X$ si et seulement si
$E \cap U$ est quasi-compact pour tout ouvert affine $U$ de $X$.
\end{lemma}

\begin{proof}
Cela résulte immédiatement du fait que tout ouvert quasi-compact de $X$ est une réunion finie
d'ouverts affines.
\end{proof}

\begin{lemma}
\label{lemma-stratification-locally-finite-constructible}
Une partition $X = \coprod_{i \in I} X_i$ d'un schéma $X$ dont les
parties sont rétrocompactes est localement finie si et seulement si ses parties
sont localement constructibles.
\end{lemma}

\begin{proof}
Voir Topologie, définitions
\ref{topology-definition-quasi-compact},
\ref{topology-definition-paritition} et
\ref{topology-definition-locally-finite},
pour les définitions de rétrocompact, de partition et de localement fini.

\medskip\noindent
Si la partition est localement finie et si $U \subset X$ est un
ouvert affine, alors $U = \coprod_{i \in I} U \cap X_i$
est une partition finie (plus précisément, toutes ses parties sauf un nombre
fini sont vides). Ainsi $U \cap X_i$ est quasi-compact
et son complémentaire est rétrocompact dans $U$, étant une réunion finie
de parties rétrocompactes. Donc $U \cap X_i$ est constructible
d'après Topologie, lemme \ref{topology-lemma-locally-closed-constructible-image}.
Il s'ensuit que $X_i$ est localement constructible par le
lemme \ref{lemma-locally-constructible}.

\medskip\noindent
Supposons les parties localement constructibles. Pour tout ouvert
affine $U \subset X$, on obtient alors un recouvrement $U = \coprod X_i \cap U$
par des sous-ensembles constructibles. La topologie constructible étant
quasi-compacte, voir
Topologie, lemme \ref{topology-lemma-constructible-hausdorff-quasi-compact},
ce recouvrement admet un raffinement fini ; autrement dit,
la partition est localement finie.
\end{proof}

\section{Schémas intègres, irréductibles et réduits}
\label{section-integral}

\begin{definition}
\label{definition-integral}
Soit $X$ un schéma. On dit que $X$ est {\it intègre} s'il est non vide et si,
pour tout ouvert affine non vide $\Spec(R) = U \subset X$, l'anneau $R$
est intègre.
\end{definition}

\begin{lemma}
\label{lemma-characterize-reduced}
Soit $X$ un schéma.
Les assertions suivantes sont équivalentes.
\begin{enumerate}
\item Le schéma $X$ est réduit, voir
Schémas, définition \ref{schemes-definition-reduced}.
\item Il existe un recouvrement ouvert affine $X = \bigcup U_i$
tel que chaque $\Gamma(U_i, \mathcal{O}_X)$ soit réduit.
\item Pour tout ouvert affine $U \subset X$, l'anneau
$\mathcal{O}_X(U)$ est réduit.
\item Pour tout ouvert $U \subset X$, l'anneau $\mathcal{O}_X(U)$ est réduit.
\end{enumerate}
\end{lemma}

\begin{proof}
Voir Schémas, lemmes \ref{schemes-lemma-reduced} et
\ref{schemes-lemma-affine-reduced}.
\end{proof}

\begin{lemma}
\label{lemma-characterize-irreducible}
Soit $X$ un schéma.
Les assertions suivantes sont équivalentes.
\begin{enumerate}
\item Le schéma $X$ est irréductible.
\item Il existe un recouvrement ouvert affine $X = \bigcup_{i \in I} U_i$
tel que $I$ soit non vide, que $U_i$ soit irréductible pour tout $i \in I$ et que
$U_i \cap U_j \not = \emptyset$ pour tous $i, j \in I$.
\item Le schéma $X$ est non vide et tout ouvert affine non vide
$U \subset X$ est irréductible.
\end{enumerate}
\end{lemma}

\begin{proof}
Supposons (1). D'après Schémas, lemme \ref{schemes-lemma-scheme-sober},
$X$ possède un unique point générique $\eta$. On a alors
$X = \overline{\{\eta\}}$. Par conséquent, $\eta$ appartient à
tout ouvert affine non vide $U \subset X$. Cela implique
que le point $\eta \in U$ est dense, donc que $U$ est irréductible.
Cela implique aussi que deux ouverts affines non vides se rencontrent.
Ainsi (1) entraîne à la fois (2) et (3).

\medskip\noindent
Supposons (2). Soit $X = Z_1 \cup Z_2$ une réunion de deux sous-ensembles fermés.
Pour tout $i$, on a soit $U_i \subset Z_1$, soit $U_i \subset Z_2$.
Choisissons un $i \in I$ et supposons $U_i \subset Z_1$ (quitte à permuter
$Z_1$, $Z_2$). Pour tout $j \in I$, l'ouvert $U_i \cap U_j$ est dense dans
$U_j$ et contenu dans le fermé $Z_1 \cap U_j$. On en déduit que
$U_j \subset Z_1$ également. Donc $X = Z_1$, comme voulu.

\medskip\noindent
Supposons (3). Choisissons un recouvrement ouvert affine $X = \bigcup_{i \in I} U_i$.
On peut supposer chaque $U_i$ non vide.
Puisque $X$ est non vide, $I$ est non vide.
Par hypothèse, chaque $U_i$ est irréductible.
Supposons $U_i \cap U_j = \emptyset$ pour un couple $i, j \in I$.
Alors l'ouvert $U_i \amalg  U_j = U_i \cup U_j$ est affine, voir
Schémas, lemme \ref{schemes-lemma-disjoint-union-affines}.
Il est donc irréductible par hypothèse, ce qui est absurde. On en conclut que (3)
entraîne (2). Le lemme est démontré.
\end{proof}

\begin{lemma}
\label{lemma-characterize-integral}
Un schéma $X$ est intègre si et seulement s'il est réduit et irréductible.
\end{lemma}

\begin{proof}
Si $X$ est irréductible, tout ouvert affine $\Spec(R) = U \subset X$
est irréductible. Si $X$ est réduit, $R$ est réduit d'après le
lemme \ref{lemma-characterize-reduced} ci-dessus. Ainsi $R$ est réduit
et $(0)$ est un idéal premier ; autrement dit, $R$ est un anneau intègre.

\medskip\noindent
Si $X$ est intègre, alors, pour tout ouvert affine non vide
$\Spec(R) = U \subset X$, l'anneau $R$ est réduit,
donc $X$ est réduit par le lemme \ref{lemma-characterize-reduced}.
De plus, tout ouvert affine non vide est irréductible.
Par conséquent, $X$ est irréductible, voir le lemme \ref{lemma-characterize-irreducible}.
\end{proof}

\noindent
Dans Exemples, section
\ref{examples-section-connected-locally-integral-not-integral},
on construit un schéma affine connexe dont tous les anneaux locaux sont intègres,
mais qui n'est pas intègre.

\section{Types de schémas définis par des propriétés d'anneaux}
\label{section-properties-rings}

\noindent
Dans cette section, nous étudions quelles propriétés d'anneaux
permettent de définir des propriétés locales des schémas.

\begin{definition}
\label{definition-property-local}
Soit $P$ une propriété des anneaux.
On dit que $P$ est {\it locale} si les conditions suivantes sont satisfaites :
\begin{enumerate}
\item Pour tout anneau $R$ et tout $f \in R$, on a
$P(R) \Rightarrow P(R_f)$.
\item Pour tout anneau $R$ et tous $f_i \in R$ tels que
$(f_1, \ldots, f_n) = R$, on a
$\forall i, P(R_{f_i}) \Rightarrow P(R)$.
\end{enumerate}
\end{definition}

\begin{definition}
\label{definition-locally-P}
Soient $P$ une propriété des anneaux et $X$ un schéma.
On dit que $X$ est {\it localement $P$} si, pour tout $x \in X$,
il existe un voisinage ouvert affine $U$ de $x$
dans $X$ tel que $\mathcal{O}_X(U)$ possède la propriété $P$.
\end{definition}

\noindent
Cette notion n'est satisfaisante que si la propriété est locale.
Même lorsque $P$ est une propriété locale, nous n'utiliserons pas
automatiquement cette définition pour dire qu'un schéma est
« localement $P$ », à moins d'énoncer aussi explicitement cette définition
ailleurs.

\begin{lemma}
\label{lemma-locally-P}
Soient $X$ un schéma et $P$ une propriété locale des anneaux.
Les assertions suivantes sont équivalentes :
\begin{enumerate}
\item Le schéma $X$ est localement $P$.
\item Pour tout ouvert affine $U \subset X$, la propriété
$P(\mathcal{O}_X(U))$ est satisfaite.
\item Il existe un recouvrement ouvert affine $X = \bigcup U_i$ tel que
chaque $\mathcal{O}_X(U_i)$ satisfasse $P$.
\item Il existe un recouvrement ouvert $X = \bigcup X_j$
tel que chaque sous-schéma ouvert $X_j$ soit localement $P$.
\end{enumerate}
De plus, si $X$ est localement $P$, tout sous-schéma ouvert
est localement $P$.
\end{lemma}

\begin{proof}
On a évidemment (1) $\Leftrightarrow$ (3) et (2) $\Rightarrow$ (1).
Si (3) $\Rightarrow$ (2), alors la dernière assertion du lemme
est vraie, et il en résulte aisément que (4) est aussi équivalente à (1).
Montrons donc (3) $\Rightarrow$ (2).

\medskip\noindent
Soit $X = \bigcup U_i$ un recouvrement ouvert affine, avec
$U_i = \Spec(R_i)$. Supposons $P(R_i)$.
Soit $\Spec(R) = U \subset X$ un ouvert affine quelconque.
D'après Schémas, lemme \ref{schemes-lemma-good-subcover},
il existe un recouvrement standard de $U = \Spec(R)$ par des
ouverts principaux $D(f_j)$ tel que chaque anneau $R_{f_j}$ soit une
localisation principale d'un des anneaux $R_i$. Par la
définition \ref{definition-property-local} (1), on obtient $P(R_{f_j})$.
D'où $P(R)$ par la définition \ref{definition-property-local} (2).
\end{proof}

\noindent
Voici un exemple d'application.

\begin{lemma}
\label{lemma-reduced-is-locally-reduced}
Soit $X$ un schéma. Alors $X$ est réduit si et seulement si $X$ est
« localement réduit » au sens de la définition \ref{definition-locally-P}.
\end{lemma}

\begin{proof}
Cela résulte immédiatement du lemme \ref{lemma-characterize-reduced}.
\end{proof}

\begin{lemma}
\label{lemma-properties-local}
Les propriétés suivantes d'un anneau $R$ sont locales.
\begin{enumerate}
\item (Cohen-Macaulay.)
L'anneau $R$ est noethérien et de Cohen--Macaulay, voir
Algèbre, définition \ref{algebra-definition-ring-CM}.
\item (Régulier.)
L'anneau $R$ est noethérien et régulier, voir
Algèbre, définition \ref{algebra-definition-regular}.
\item (Absolument noethérien.)
L'anneau $R$ est de type fini sur $Z$.
\item Ajouter ici d'autres propriétés au besoin.\footnote{Nous ne répertorions ici que les propriétés
qui n'ont pas déjà été traitées séparément ailleurs.}
\end{enumerate}
\end{lemma}

\begin{proof}
Omis.
\end{proof}

\section{Schémas noethériens}
\label{section-noetherian}

\noindent
Rappelons qu'un anneau $R$ est {\it noethérien} s'il satisfait à la condition
de chaîne ascendante sur les idéaux. De manière équivalente, tout idéal de $R$ est de type
fini.

\begin{definition}
\label{definition-noetherian}
Soit $X$ un schéma.
\begin{enumerate}
\item On dit que $X$ est {\it localement noethérien} si tout
$x \in X$ possède un voisinage ouvert affine
$\Spec(R) = U \subset X$ tel que l'anneau $R$ soit noethérien.
\item On dit que $X$ est {\it noethérien} si $X$ est localement noethérien
et quasi-compact.
\end{enumerate}
\end{definition}

\noindent
Voici la caractérisation usuelle des schémas localement noethériens.

\begin{lemma}
\label{lemma-locally-Noetherian}
Soit $X$ un schéma. Les assertions suivantes sont équivalentes :
\begin{enumerate}
\item Le schéma $X$ est localement noethérien.
\item Pour tout ouvert affine $U \subset X$, l'anneau $\mathcal{O}_X(U)$
est noethérien.
\item Il existe un recouvrement ouvert affine $X = \bigcup U_i$ tel que
chaque $\mathcal{O}_X(U_i)$ soit noethérien.
\item Il existe un recouvrement ouvert $X = \bigcup X_j$
tel que chaque sous-schéma ouvert $X_j$ soit localement noethérien.
\end{enumerate}
De plus, si $X$ est localement noethérien, tout sous-schéma ouvert
est localement noethérien.
\end{lemma}

\begin{proof}
Il suffit de montrer qu'être noethérien est une propriété locale
des anneaux, voir le lemme \ref{lemma-locally-P}.
Toute localisation d'un anneau noethérien est noethérienne, voir
Algèbre, lemme \ref{algebra-lemma-Noetherian-permanence}.
Le lemme \ref{algebra-lemma-cover} du chapitre Algèbre donne la seconde
condition de la définition \ref{definition-property-local}.
\end{proof}

\begin{lemma}
\label{lemma-immersion-into-noetherian}
Toute immersion $Z \to X$, avec $X$ localement noethérien, est quasi-compacte.
\end{lemma}

\begin{proof}
Une immersion fermée est évidemment quasi-compacte.
Un composé de morphismes quasi-compacts est quasi-compact,
voir Topologie, lemme \ref{topology-lemma-composition-quasi-compact}.
Il suffit donc de montrer qu'une immersion ouverte dans
un schéma localement noethérien est quasi-compacte.
À l'aide de Schémas, lemme \ref{schemes-lemma-quasi-compact-affine},
on se ramène au cas où $X$ est affine.
Tout ouvert du spectre d'un anneau noethérien
est quasi-compact (on peut, par exemple,
combiner Algèbre, lemme \ref{algebra-lemma-Noetherian-topology}, et
Topologie, lemmes \ref{topology-lemma-Noetherian} et
\ref{topology-lemma-Noetherian-quasi-compact}).
\end{proof}

\begin{lemma}
\label{lemma-locally-Noetherian-quasi-separated}
Un schéma localement noethérien est quasi-séparé.
\end{lemma}

\begin{proof}
D'après Schémas, lemme \ref{schemes-lemma-characterize-quasi-separated},
il faut montrer que l'intersection $U \cap V$ de deux
ouverts affines de $X$ est quasi-compacte. Cela résulte du
lemme \ref{lemma-immersion-into-noetherian} ci-dessus, en
considérant, par exemple, l'immersion ouverte $U \cap V \to U$.
(Au fond, cela vient simplement de ce que tout ouvert du spectre d'un
anneau noethérien est quasi-compact.)
\end{proof}

\begin{lemma}
\label{lemma-Noetherian-topology}
L'espace topologique sous-jacent à un schéma (localement) noethérien
est (localement) noethérien,
voir Topologie, définition \ref{topology-definition-noetherian}.
\end{lemma}

\begin{proof}
Cela résulte du fait qu'un schéma noethérien est une réunion finie de spectres
d'anneaux noethériens, ainsi que des résultats suivants :
Algèbre, lemme \ref{algebra-lemma-Noetherian-topology}, et
Topologie, lemme \ref{topology-lemma-finite-union-Noetherian}.
\end{proof}

\begin{lemma}
\label{lemma-locally-closed-in-Noetherian}
Tout sous-schéma localement fermé d'un schéma (localement) noethérien
est (localement) noethérien.
\end{lemma}

\begin{proof}
Omis. Indication : tout quotient et toute localisation d'un anneau
noethérien sont noethériens. Dans le cas noethérien, utiliser de nouveau
le fait que tout sous-ensemble d'un espace noethérien est noethérien
pour la topologie induite.
\end{proof}

\begin{lemma}
\label{lemma-Noetherian-irreducible-components}
Un schéma noethérien possède un nombre fini de composantes irréductibles.
\end{lemma}

\begin{proof}
L'espace topologique sous-jacent à un schéma noethérien est noethérien
(lemme \ref{lemma-Noetherian-topology}),
et l'on conclut parce qu'un espace topologique noethérien
n'a qu'un nombre fini de composantes irréductibles
(Topologie, lemme \ref{topology-lemma-Noetherian}).
\end{proof}

\begin{lemma}
\label{lemma-morphism-Noetherian-schemes-quasi-compact}
Tout morphisme de schémas $f : X \to Y$, avec $X$ noethérien,
est quasi-compact.
\end{lemma}

\begin{proof}
Utiliser le lemme \ref{lemma-Noetherian-topology}
et le fait que tout sous-ensemble d'un espace topologique
noethérien est quasi-compact (voir Topologie,
lemmes \ref{topology-lemma-Noetherian} et
\ref{topology-lemma-Noetherian-quasi-compact}).
\end{proof}

\noindent
Voici un lemme intéressant.
Il affirme que tout schéma localement noethérien possède beaucoup de
points fermés (au moins un dans tout fermé non vide).

\begin{lemma}
\label{lemma-locally-Noetherian-closed-point}
Tout schéma localement noethérien non vide possède un point fermé.
Tout fermé non vide d'un schéma localement noethérien possède un point fermé.
De manière équivalente, tout point d'un schéma localement noethérien se spécialise
en un point fermé.
\end{lemma}

\begin{proof}
La seconde assertion découle de la première (à l'aide de
Schémas, lemme \ref{schemes-lemma-reduced-closed-subscheme},
et du lemme \ref{lemma-locally-closed-in-Noetherian}).
Considérons un ouvert affine non vide $U \subset X$.
Soit $x \in U$ un point fermé. Si $x$ est fermé
dans $X$, la démonstration est terminée. Sinon, soit $X_0 \subset X$ le
sous-schéma fermé réduit induit sur $\overline{\{x\}}$.
Alors $U_0 = U \cap X_0$ est un ouvert affine de $X_0$ d'après
Schémas, lemme \ref{schemes-lemma-closed-subspace-scheme}, et
$U_0 = \{x\}$. Soit $y \in X_0$, $y \not = x$, une spécialisation de $x$.
Considérons l'anneau local $R = \mathcal{O}_{X_0, y}$.
C'est un anneau local noethérien, puisque $X_0$ est noethérien
d'après le lemme \ref{lemma-locally-closed-in-Noetherian}. Notons $V \subset \Spec(R)$
l'image inverse de $U_0$ dans $\Spec(R)$ par le morphisme canonique
$\Spec(R) \to X_0$ (voir Schémas, section \ref{schemes-section-points}).
Par construction, $V$ est un singleton dont l'unique point correspond à $x$ (utiliser
Schémas, lemme \ref{schemes-lemma-specialize-points}).
D'après
Algèbre, lemme \ref{algebra-lemma-Noetherian-local-domain-dim-2-infinite-opens},
on a $\dim(R) = 1$.
Autrement dit, $y$ est une spécialisation immédiate
de $x$ (voir Topologie, définition \ref{topology-definition-dimension-function}).
Ainsi, tout
point $y \not = x$ tel que $x \leadsto y$ est une spécialisation
immédiate de $x$. Chacun de ces points est évidemment
fermé, comme voulu.
\end{proof}

\begin{lemma}
\label{lemma-locally-Noetherian-specialization-dvr}
Soit $X$ un schéma localement noethérien.
Soit $x' \leadsto x$ une spécialisation de points de $X$. Alors
\begin{enumerate}
\item il existe un anneau de valuation discrète $R$ et un morphisme
$f : \Spec(R) \to X$ tels que le point générique $\eta$ de
$\Spec(R)$ ait pour image $x'$ et le point fermé ait pour image $x$ ;
\item si $x \not = x'$, étant donnée une extension de corps de type fini
$K/\kappa(x')$, on peut faire en sorte que l'extension
$\kappa(\eta)/\kappa(x')$
induite par $f$ soit isomorphe à l'extension donnée.
\end{enumerate}
\end{lemma}

\begin{proof}
Supposons d'abord que $x' \leadsto x$ soit une spécialisation dans $X$
avec $x \not = x'$, et soit $K/\kappa(x')$ une extension
de corps de type fini. D'après
Schémas, lemme \ref{schemes-lemma-specialize-points},
et la discussion qui suit
Schémas, lemme \ref{schemes-lemma-characterize-points},
on obtient des homomorphismes d'anneaux $\mathcal{O}_{X, x} \to \kappa(x') \to K$.
Puisque $x \not = x'$, l'image de $\mathcal{O}_{X, x}$ dans $K$
n'est pas un corps (détails omis).
Soit $R \subset K$ un anneau de valuation discrète dont le corps des fractions est
$K$ et qui domine l'image de $\mathcal{O}_{X, x} \to K$, voir
Algèbre, lemme \ref{algebra-lemma-exists-dvr}.
L'homomorphisme d'anneaux $\mathcal{O}_{X, x} \to R$ induit le morphisme
$f : \Spec(R) \to X$, voir
Schémas, lemme \ref{schemes-lemma-morphism-from-spec-local-ring}.
Par construction, ce morphisme possède toutes les propriétés voulues.
Si $x = x'$, on peut poser $R = \kappa(x)[t]_{(t)}$.
\end{proof}

\begin{lemma}
\label{lemma-thin-infinite-sequence}
Soient $S$ un schéma noethérien et $T \subset S$ un sous-ensemble infini.
Il existe alors un sous-ensemble infini $T' \subset T$
tel qu'il n'y ait aucune spécialisation non triviale entre les points de $T'$.
\end{lemma}

\begin{proof}
Soit $T_0 \subset T$ l'ensemble des $t \in T$ qui ne se spécialisent
en aucun autre point de $T$. Si $T_0$ est infini, $T' = T_0$ convient.
On peut donc supposer $T_0$ fini.
Pour $i > 0$, considérons, par récurrence, l'ensemble $T_i \subset T$
des $t \in T$ tels que
\begin{enumerate}
\item $t \not \in T_{i - 1} \cup T_{i - 2} \cup \ldots \cup T_0$ ;
\item il existe une spécialisation non triviale $t \leadsto t'$ avec
$t' \in T_{i - 1}$ ;
\item pour toute spécialisation non triviale
$t \leadsto t'$ avec $t' \in T$, on a
$t' \in T_{i - 1} \cup T_{i - 2} \cup \ldots \cup T_0$.
\end{enumerate}
Là encore, si $T_i$ est infini, $T' = T_i$ convient.
Soit $d$ le maximum des dimensions des anneaux locaux
$\mathcal{O}_{S, t}$ pour $t \in T_0$ ; alors $d$ est un entier,
car $T_0$ est fini et les dimensions des anneaux locaux
sont finies d'après Algèbre, proposition \ref{algebra-proposition-dimension}.
On a $T_i = \emptyset$ pour $i > d$.
En effet, si $t \in T_i$, on peut trouver une suite
de spécialisations non triviales
$t = t_i \leadsto t_{i - 1} \leadsto \ldots \leadsto t_0$
avec $t_0 \in T_0$. Comme
les points $t = t_i, t_{i - 1}, \ldots, t_0$ appartiennent à
$\Spec(\mathcal{O}_{S, t_0})$
(Schémas, lemme \ref{schemes-lemma-specialize-points}),
on obtient $i \leq d$.
Ainsi $\bigcup T_i = T_d \cup \ldots \cup T_0$ est un sous-ensemble fini de $T$.

\medskip\noindent
Supposons qu'un $t \in T$ n'appartienne pas à $\bigcup T_i$. Il doit alors exister une spécialisation
$t \leadsto t'$ avec $t' \in T$ et $t' \not \in \bigcup T_i$. (En effet, si
toutes les spécialisations de $t$ appartiennent à l'ensemble fini $T_d \cup \ldots \cup T_0$,
il existe un plus grand $i$ tel qu'il y ait une spécialisation
$t \leadsto t'$ avec $t' \in T_i$, et alors $t \in T_{i + 1}$ par construction.)
On obtient donc une suite infinie
$$
t \leadsto t' \leadsto t'' \leadsto \ldots
$$
de spécialisations non triviales entre points de $T \setminus \bigcup T_i$.
C'est impossible, car l'espace topologique sous-jacent à $S$
est noethérien d'après le lemme \ref{lemma-locally-Noetherian-quasi-separated}.
\end{proof}

\begin{lemma}
\label{lemma-maximal-points}
Soient $S$ un schéma noethérien et $T \subset S$ un sous-ensemble. Soit
$T_0 \subset T$ l'ensemble des $t \in T$ tels qu'il n'existe aucune spécialisation
non triviale $t' \leadsto t$ avec $t' \in T'$. Alors (a) il n'y a
aucune spécialisation entre les points de $T_0$, (b) tout point de
$T$ est une spécialisation d'un point de $T_0$, et (c) les adhérences
de $T$ et de $T_0$ sont égales.
\end{lemma}

\begin{proof}
Rappelons que $\dim(\mathcal{O}_{S, s}) < \infty$ pour tout $s \in S$, voir
Algèbre, proposition \ref{algebra-proposition-dimension}. Soit $t \in T$.
Si $t' \leadsto t$, la théorie de la dimension donne
$\dim(\mathcal{O}_{S, t'}) \leq \dim(\mathcal{O}_{S, t})$,
avec égalité si et seulement si $t' = t$. Ainsi, si l'on choisit $t' \leadsto t$
avec $\dim(\mathcal{O}_{T, t'})$ minimal, alors $t' \in T_0$.
Autrement dit,
tout $t \in T$ est une spécialisation d'un élément de $T_0$.
\end{proof}

\begin{lemma}
\label{lemma-countable-dense-subset}
Soient $S$ un schéma noethérien et $T \subset S$ un sous-ensemble infini dense.
Il existe alors un sous-ensemble dénombrable $E \subset T$ dense dans $S$.
\end{lemma}

\begin{proof}
Soit $T'$ l'ensemble des points $s \in S$ tels que $\overline{\{s\}} \cap T$
contienne un sous-ensemble dénombrable d'adhérence $\overline{\{s\}}$.
Tout ensemble fini étant dénombrable, on a $T \subset T'$.
Pour $s \in T'$, choisissons un tel sous-ensemble dénombrable
$E_s \subset \overline{\{s\}} \cap T$.
Soit $E' = \{s_1, s_2, s_3, \ldots\} \subset T'$
un sous-ensemble dénombrable. Alors l'adhérence de $E'$ dans $S$ est
l'adhérence du sous-ensemble dénombrable $\bigcup_n E_{s_n}$ de $T$.
Il s'ensuit que, si $Z$
est une composante irréductible de l'adhérence de $E'$, le point
générique de $Z$ appartient à $T'$.

\medskip\noindent
Notons $T'_0 \subset T'$ le sous-ensemble des $t \in T'$ tels
qu'il n'existe aucune spécialisation non triviale $t' \leadsto t$ avec $t' \in T'$,
comme dans le lemme \ref{lemma-maximal-points}, dont nous utiliserons les résultats
sans le mentionner à nouveau. Si $T'_0$ est infini, choisissons un
sous-ensemble dénombrable $E' \subset T'_0$. D'après l'argument du premier
paragraphe, les points génériques des composantes irréductibles de
l'adhérence de $E'$ appartiennent à $T'$. Mais l'un de ces points se spécialise en
une infinité d'éléments distincts de $E' \subset T'_0$,
ce qui est contradictoire. Ainsi $T'_0$ est fini ; écrivons
$T'_0 = \{s_1, \ldots, s_m\}$. Il en résulte que $S$, qui est
l'adhérence de $T$, est contenu dans l'adhérence de
$\{s_1, \ldots, s_m\}$, elle-même contenue dans l'adhérence
du sous-ensemble dénombrable $E_{s_1} \cup \ldots \cup E_{s_m} \subset T$,
comme voulu.
\end{proof}
\section{Schémas de Jacobson}
\label{section-jacobson}

\noindent
Rappelons qu'un espace est dit {\it de Jacobson} si ses points fermés sont
denses dans tout sous-ensemble fermé, voir
Topologie, section \ref{topology-section-space-jacobson}.

\begin{definition}
\label{definition-jacobson}
Un schéma $S$ est dit {\it de Jacobson} si son espace topologique
sous-jacent est de Jacobson.
\end{definition}

\noindent
Rappelons qu'un anneau $R$ est de Jacobson si tout idéal radical de $R$
est une intersection d'idéaux maximaux, voir
Algèbre, définition \ref{algebra-definition-ring-jacobson}.

\begin{lemma}
\label{lemma-affine-jacobson}
Un schéma affine $\Spec(R)$ est de Jacobson si et seulement si
l'anneau $R$ est de Jacobson.
\end{lemma}

\begin{proof}
C'est le lemme \ref{algebra-lemma-jacobson} du chapitre Algèbre.
\end{proof}

\noindent
Voici la caractérisation usuelle des schémas de Jacobson.
Intuitivement, elle affirme : de Jacobson $\Leftrightarrow$ localement de Jacobson.

\begin{lemma}
\label{lemma-locally-jacobson}
Soit $X$ un schéma. Les assertions suivantes sont équivalentes :
\begin{enumerate}
\item Le schéma $X$ est de Jacobson.
\item Le schéma $X$ est « localement de Jacobson » au sens de la
définition \ref{definition-locally-P}.
\item Pour tout ouvert affine $U \subset X$, l'anneau $\mathcal{O}_X(U)$
est de Jacobson.
\item Il existe un recouvrement ouvert affine $X = \bigcup U_i$ tel que
chaque $\mathcal{O}_X(U_i)$ soit de Jacobson.
\item Il existe un recouvrement ouvert $X = \bigcup X_j$
tel que chaque sous-schéma ouvert $X_j$ soit de Jacobson.
\end{enumerate}
De plus, si $X$ est de Jacobson, tout sous-schéma ouvert
est de Jacobson.
\end{lemma}

\begin{proof}
La dernière assertion du lemme résulte de
Topologie, lemme \ref{topology-lemma-jacobson-inherited}.
L'équivalence de (5) et de (1) est donnée par
Topologie, lemme \ref{topology-lemma-jacobson-local}.
Ainsi, à l'aide du lemme \ref{lemma-affine-jacobson},
on voit que (1) $\Leftrightarrow$ (2).
Pour achever la démonstration, il suffit de montrer qu'être de Jacobson
est une propriété locale des anneaux, voir le lemme \ref{lemma-locally-P}.
Toute localisation d'un anneau de Jacobson en un élément est de Jacobson, voir
Algèbre, lemme \ref{algebra-lemma-Jacobson-invert-element}.
Soit $R$ un anneau ; supposons que $f_1, \ldots, f_n \in R$ engendrent l'idéal
unité et que chaque $R_{f_i}$ soit de Jacobson. Alors
$\Spec(R) = \bigcup D(f_i)$ est une réunion d'ouverts
tous de Jacobson ; donc $\Spec(R)$ est de Jacobson,
de nouveau par Topologie, lemme \ref{topology-lemma-jacobson-local}.
Cela démontre la seconde condition de la définition \ref{definition-property-local}.
\end{proof}

\noindent
Beaucoup de schémas couramment utilisés en géométrie algébrique sont de Jacobson, voir
Morphismes, lemme \ref{morphisms-lemma-ubiquity-Jacobson-schemes}.
Mentionnons ici le cas intéressant suivant.

\begin{lemma}
\label{lemma-complement-closed-point-Jacobson}
Exemples de schémas noethériens de Jacobson.
\begin{enumerate}
\item Si $(R, \mathfrak m)$ est un anneau local noethérien, alors
le spectre épointé $\Spec(R) \setminus \{\mathfrak m\}$
est un schéma de Jacobson.
\item Si $R$ est un anneau noethérien de radical de Jacobson $\text{rad}(R)$,
alors $\Spec(R) \setminus V(\text{rad}(R))$ est un schéma de Jacobson.
\item Si $(R, I)$ est un couple de Zariski (Compléments sur l'algèbre, définition
\ref{more-algebra-definition-zariski-pair}),
avec $R$ noethérien, alors $\Spec(R) \setminus V(I)$ est un
schéma de Jacobson.
\end{enumerate}
\end{lemma}

\begin{proof}
Démonstration de (3). Remarquons que $\Spec(R) - V(I)$ est recouvert par les
ouverts affines $\Spec(R_f)$ pour $f \in I$. Les anneaux $R_f$ sont de Jacobson d'après
Compléments sur l'algèbre, lemme
\ref{more-algebra-lemma-noetherian-zariski-jacobson-complement}.
Donc $\Spec(R) \setminus V(I)$ est de Jacobson par le
lemme \ref{lemma-locally-jacobson}.
Les assertions (1) et (2) sont des cas particuliers de (3).

\medskip\noindent
Démonstration directe de (1).
Puisque $\Spec(R)$ est un schéma noethérien,
$S$ est un schéma noethérien (lemme \ref{lemma-locally-closed-in-Noetherian}).
Ainsi $S$ est un espace topologique sobre et noethérien (utiliser
Schémas, lemme \ref{schemes-lemma-scheme-sober}).
Supposons que $S$ ne soit pas de Jacobson et cherchons
une contradiction. D'après
Topologie, lemme \ref{topology-lemma-non-jacobson-Noetherian-characterize},
il existe un point non fermé $\xi \in S$
tel que $\{\xi\}$ soit localement fermé. Il correspond
à un idéal premier $\mathfrak p \subset R$ tel que (1) il existe
un idéal premier $\mathfrak q$, $\mathfrak p \subset \mathfrak q \subset \mathfrak m$,
les deux inclusions étant strictes, et (2) $\{\mathfrak p\}$ soit ouvert dans
$\Spec(R/\mathfrak p)$. C'est impossible d'après Algèbre,
lemme \ref{algebra-lemma-Noetherian-local-domain-dim-2-infinite-opens}.
\end{proof}

\section{Schémas normaux}
\label{section-normal}

\noindent
Rappelons qu'un anneau $R$ est dit normal si tous ses anneaux locaux
sont des anneaux intègres normaux,
voir Algèbre, définition \ref{algebra-definition-ring-normal}.
Un anneau intègre normal est un anneau intègre intégralement clos dans son corps
des fractions, voir
Algèbre, définition \ref{algebra-definition-domain-normal}.
Il est donc naturel de définir un schéma normal comme suit.

\begin{definition}
\label{definition-normal}
Un schéma $X$ est {\it normal} si et seulement si, pour tout $x \in X$, l'anneau local
$\mathcal{O}_{X, x}$ est un anneau intègre normal.
\end{definition}

\noindent
Cela semble être la définition utilisée dans EGA, voir \cite[0, 4.1.4]{EGA}.
Supposons $X = \Spec(A)$, avec $A$ réduit. Dire que $X$ est
normal n'équivaut pas à dire que $A$ est intégralement clos dans son
anneau total des fractions. Cette équivalence est toutefois vraie si $A$ est noethérien
(voir Algèbre, lemme \ref{algebra-lemma-characterize-reduced-ring-normal}).

\begin{lemma}
\label{lemma-locally-normal}
Soit $X$ un schéma. Les assertions suivantes sont équivalentes :
\begin{enumerate}
\item Le schéma $X$ est normal.
\item Pour tout ouvert affine $U \subset X$, l'anneau $\mathcal{O}_X(U)$
est normal.
\item Il existe un recouvrement ouvert affine $X = \bigcup U_i$ tel que
chaque $\mathcal{O}_X(U_i)$ soit normal.
\item Il existe un recouvrement ouvert $X = \bigcup X_j$
tel que chaque sous-schéma ouvert $X_j$ soit normal.
\end{enumerate}
De plus, si $X$ est normal, tout sous-schéma ouvert
est normal.
\end{lemma}

\begin{proof}
Cela résulte immédiatement des définitions.
\end{proof}

\begin{lemma}
\label{lemma-normal-reduced}
Un schéma normal est réduit.
\end{lemma}

\begin{proof}
Cela résulte immédiatement des définitions.
\end{proof}

\begin{lemma}
\label{lemma-integral-normal}
Soit $X$ un schéma intègre.
Alors $X$ est normal si et seulement si, pour tout ouvert affine non vide
$U \subset X$, l'anneau $\mathcal{O}_X(U)$ est un anneau intègre normal.
\end{lemma}

\begin{proof}
Cela résulte de
Algèbre, lemme \ref{algebra-lemma-normality-is-local}.
\end{proof}

\begin{lemma}
\label{lemma-normal-locally-finite-nr-irreducibles}
Soit $X$ un schéma dont tout ouvert quasi-compact possède un nombre fini
de composantes irréductibles. Les assertions suivantes sont équivalentes :
\begin{enumerate}
\item $X$ est normal ;
\item $X$ est une réunion disjointe de schémas intègres normaux.
\end{enumerate}
\end{lemma}

\begin{proof}
Il résulte immédiatement des définitions que (2) entraîne (1).
Soit $X$ un schéma normal dont tout ouvert quasi-compact
possède un nombre fini de composantes irréductibles.
Si $X$ est affine, $X$ satisfait (2) d'après
Algèbre, lemme \ref{algebra-lemma-characterize-reduced-ring-normal}.
Pour $X$ quelconque, soit $X = \bigcup X_i$ un
recouvrement ouvert affine. Remarquons que chaque $X_i$ n'a lui aussi
qu'un nombre fini de composantes irréductibles et que le lemme s'applique
à chaque $X_i$. Soit $T \subset X$ une composante irréductible.
D'après le cas affine, chaque intersection $T \cap X_i$ est ouverte dans $X_i$
et est un schéma intègre normal.
Ainsi $T \subset X$ est ouvert et est un schéma intègre normal.
Cela prouve que $X$ est la réunion disjointe de ses composantes irréductibles,
qui sont des schémas intègres normaux.
\end{proof}

\begin{lemma}
\label{lemma-normal-Noetherian}
Soit $X$ un schéma noethérien.
Les assertions suivantes sont équivalentes :
\begin{enumerate}
\item $X$ est normal ;
\item $X$ est une réunion disjointe finie de schémas intègres normaux.
\end{enumerate}
\end{lemma}

\begin{proof}
C'est un cas particulier du
lemme \ref{lemma-normal-locally-finite-nr-irreducibles}, car l'espace topologique
sous-jacent à un schéma noethérien est noethérien
(lemme \ref{lemma-Noetherian-topology}
et
Topologie, lemme \ref{topology-lemma-Noetherian}).
\end{proof}

\begin{lemma}
\label{lemma-normal-locally-Noetherian}
Soit $X$ un schéma localement noethérien.
Les assertions suivantes sont équivalentes :
\begin{enumerate}
\item $X$ est normal ;
\item $X$ est une réunion disjointe de schémas intègres normaux.
\end{enumerate}
\end{lemma}

\begin{proof}
Omis. Indication : cela se déduit du
lemme \ref{lemma-normal-Noetherian} par un raisonnement purement topologique.
\end{proof}

\begin{remark}
\label{remark-normal-connected-irreducible}
Soit $X$ un schéma normal. Si $X$ est localement noethérien,
$X$ est intègre si et seulement si $X$ est connexe, voir le
lemme \ref{lemma-normal-locally-Noetherian}.
Mais il existe un schéma affine connexe $X$ tel que
$\mathcal{O}_{X, x}$ soit intègre pour tout $x \in X$, sans que $X$ soit
irréductible, voir Exemples, section
\ref{examples-section-connected-locally-integral-not-integral}.
Cet exemple est même un schéma normal (démonstration omise) ; il faut donc y prendre garde !
\end{remark}

\begin{lemma}
\label{lemma-normal-integral-sections}
\begin{slogan}
L'anneau des fonctions sur un schéma normal est normal.
\end{slogan}
Soit $X$ un schéma intègre normal.
Alors $\Gamma(X, \mathcal{O}_X)$ est un anneau intègre normal.
\end{lemma}

\begin{proof}
Posons $R = \Gamma(X, \mathcal{O}_X)$.
Il est clair que $R$ est un anneau intègre.
Supposons que $f = a/b$ soit un élément de son corps des fractions
entier sur $R$. Écrivons
$f^d + \sum_{i = 0, \ldots, d - 1} a_i f^i = 0$, avec
$a_i \in R$. Soit $U \subset X$ un ouvert affine non vide.
Puisque $b \in R$ est non nul et que $X$ est intègre,
$b|_U \in \mathcal{O}_X(U)$ est également non nul.
Ainsi $a/b$ est un élément du corps des fractions de
$\mathcal{O}_X(U)$ entier sur $\mathcal{O}_X(U)$
(car on peut utiliser la même équation polynomiale
$f^d + \sum_{i = 0, \ldots, d - 1} a_i|_U f^i = 0$ sur $U$).
Comme $\mathcal{O}_X(U)$ est un anneau intègre normal
(lemme \ref{lemma-integral-normal}), on obtient
$f_U = (a|_U)/(b|_U) \in \mathcal{O}_X(U)$. Il est clair
que $f_U|_V = f_V$ dès que $V \subset U \subset X$ sont des
ouverts affines non vides. Les sections locales $f_U$ se recollent donc en un élément
$g \in R = \Gamma(X, \mathcal{O}_X)$. Alors $bg$ et $a$
se restreignent au même élément de $\mathcal{O}_X(U)$ pour
tout $U$ comme ci-dessus ; donc $bg = a$. Autrement dit, $g$ a pour image $f$
dans le corps des fractions de $R$.
\end{proof}

\section{Schémas de Cohen--Macaulay}
\label{section-Cohen-Macaulay}

\noindent
Rappelons, voir Algèbre, définition \ref{algebra-definition-local-ring-CM},
qu'un anneau local noethérien $(R, \mathfrak m)$ est
dit de Cohen--Macaulay si $\text{depth}_{\mathfrak m}(R) = \dim(R)$.
Rappelons qu'un anneau noethérien $R$ est dit de Cohen--Macaulay si
tout anneau local $R_{\mathfrak p}$ de $R$ est de Cohen--Macaulay,
voir Algèbre, définition \ref{algebra-definition-ring-CM}.

\begin{definition}
\label{definition-Cohen-Macaulay}
Soit $X$ un schéma. On dit que $X$ est {\it de Cohen--Macaulay} si,
pour tout $x \in X$, il existe un voisinage ouvert affine
$U \subset X$ de $x$ tel que l'anneau $\mathcal{O}_X(U)$ soit
noethérien et de Cohen--Macaulay.
\end{definition}

\begin{lemma}
\label{lemma-characterize-Cohen-Macaulay}
Soit $X$ un schéma. Les assertions suivantes sont équivalentes :
\begin{enumerate}
\item $X$ est de Cohen--Macaulay ;
\item $X$ est localement noethérien et tous ses anneaux locaux sont de Cohen--Macaulay ;

\item $X$ est localement noethérien et, pour tout point fermé $x \in X$,
l'anneau local $\mathcal{O}_{X, x}$ est de Cohen--Macaulay.
\end{enumerate}
\end{lemma}

\begin{proof}
Algèbre, lemme \ref{algebra-lemma-localize-CM}, affirme que toute localisation
d'un anneau local de Cohen--Macaulay est de Cohen--Macaulay. Le lemme s'en déduit
en combinant ce résultat avec le lemme \ref{lemma-locally-Noetherian},
l'existence de points
fermés sur les schémas localement noethériens
(lemme \ref{lemma-locally-Noetherian-closed-point}) et
les définitions.
\end{proof}

\begin{lemma}
\label{lemma-locally-Cohen-Macaulay}
Soit $X$ un schéma. Les assertions suivantes sont équivalentes :
\begin{enumerate}
\item Le schéma $X$ est de Cohen--Macaulay.
\item Pour tout ouvert affine $U \subset X$, l'anneau $\mathcal{O}_X(U)$
est noethérien et de Cohen--Macaulay.
\item Il existe un recouvrement ouvert affine $X = \bigcup U_i$ tel que
chaque $\mathcal{O}_X(U_i)$ soit noethérien et de Cohen--Macaulay.
\item Il existe un recouvrement ouvert $X = \bigcup X_j$
tel que chaque sous-schéma ouvert $X_j$ soit de Cohen--Macaulay.
\end{enumerate}
De plus, si $X$ est de Cohen--Macaulay, tout sous-schéma ouvert
est de Cohen--Macaulay.
\end{lemma}

\begin{proof}
Combiner les lemmes \ref{lemma-locally-Noetherian}
et \ref{lemma-characterize-Cohen-Macaulay}.
\end{proof}

\noindent
On trouvera davantage de résultats sur les schémas de Cohen--Macaulay et la profondeur dans
Cohomologie des schémas, section \ref{coherent-section-depth}.

\section{Schémas réguliers}
\label{section-regular}

\noindent
Rappelons, voir Algèbre, définition \ref{algebra-definition-regular-local},
qu'un anneau local noethérien $(R, \mathfrak m)$ est
dit {\it régulier} si $\mathfrak m$ peut être engendré
par $\dim(R)$ éléments.
Rappelons qu'un anneau noethérien $R$ est dit {\it régulier} si
tout anneau local $R_{\mathfrak p}$ de $R$ est régulier,
voir Algèbre, définition \ref{algebra-definition-regular}.

\begin{definition}
\label{definition-regular}
Soit $X$ un schéma. On dit que $X$ est {\it régulier}, ou {\it non singulier}, si,
pour tout $x \in X$, il existe un voisinage ouvert affine
$U \subset X$ de $x$ tel que l'anneau $\mathcal{O}_X(U)$ soit
noethérien et régulier.
\end{definition}

\begin{lemma}
\label{lemma-characterize-regular}
Soit $X$ un schéma. Les assertions suivantes sont équivalentes :
\begin{enumerate}
\item $X$ est régulier ;
\item $X$ est localement noethérien et tous ses anneaux locaux sont réguliers ;

\item $X$ est localement noethérien et, pour tout point fermé $x \in X$,
l'anneau local $\mathcal{O}_{X, x}$ est régulier.
\end{enumerate}
\end{lemma}

\begin{proof}
D'après la discussion précédant, dans le chapitre Algèbre, la définition
\ref{algebra-definition-regular}, toute localisation
d'un anneau local régulier est régulière. Le lemme s'en déduit
en combinant ce résultat avec le lemme \ref{lemma-locally-Noetherian},
l'existence de points
fermés sur les schémas localement noethériens
(lemme \ref{lemma-locally-Noetherian-closed-point}) et
les définitions.
\end{proof}

\begin{lemma}
\label{lemma-locally-regular}
Soit $X$ un schéma. Les assertions suivantes sont équivalentes :
\begin{enumerate}
\item Le schéma $X$ est régulier.
\item Pour tout ouvert affine $U \subset X$, l'anneau $\mathcal{O}_X(U)$
est noethérien et régulier.
\item Il existe un recouvrement ouvert affine $X = \bigcup U_i$ tel que
chaque $\mathcal{O}_X(U_i)$ soit noethérien et régulier.
\item Il existe un recouvrement ouvert $X = \bigcup X_j$
tel que chaque sous-schéma ouvert $X_j$ soit régulier.
\end{enumerate}
De plus, si $X$ est régulier, tout sous-schéma ouvert est régulier.
\end{lemma}

\begin{proof}
Combiner les lemmes \ref{lemma-locally-Noetherian}
et \ref{lemma-characterize-regular}.
\end{proof}

\begin{lemma}
\label{lemma-regular-normal}
Un schéma régulier est normal.
\end{lemma}

\begin{proof}
Voir
Algèbre, lemme \ref{algebra-lemma-regular-normal}.
\end{proof}

\section{Dimension}
\label{section-dimension}

\noindent
La dimension d'un schéma est simplement la dimension de son espace
topologique sous-jacent.

\begin{definition}
\label{definition-dimension}
Soit $X$ un schéma.
\begin{enumerate}
\item La {\it dimension} de $X$ est la dimension de $X$
en tant qu'espace topologique, voir
Topologie, définition \ref{topology-definition-Krull}.
\item Pour $x \in X$, on note $\dim_x(X)$ la dimension de l'espace topologique
sous-jacent à $X$ en $x$, au sens de
Topologie, définition \ref{topology-definition-Krull}.
On appelle $\dim_x(X)$ la {\it dimension de $X$ en $x$}.
\end{enumerate}
\end{definition}

\noindent
L'espace topologique sous-jacent à un schéma étant sobre
(Schémas, lemme \ref{schemes-lemma-scheme-sober}),
on peut calculer la dimension de $X$ comme la borne supérieure des longueurs $n$
des chaînes
$$
T_0 \subset T_1 \subset \ldots \subset T_n
$$
de fermés irréductibles de $X$, ou comme la borne supérieure des longueurs $n$
des chaînes de spécialisations
$$
\xi_n \leadsto \xi_{n - 1} \leadsto \ldots \leadsto \xi_0
$$
de points de $X$.

\begin{lemma}
\label{lemma-dimension}
Soit $X$ un schéma. Les quantités suivantes sont égales :
\begin{enumerate}
\item la dimension de $X$ ;
\item la borne supérieure des dimensions des anneaux locaux de $X$ ;
\item la borne supérieure des $\dim_x(X)$ pour $x \in X$.
\end{enumerate}
\end{lemma}

\begin{proof}
Remarquons qu'étant donnée une chaîne de spécialisations
$$
\xi_n \leadsto \xi_{n - 1} \leadsto \ldots \leadsto \xi_0
$$
de points de $X$, tous les points $\xi_i$ correspondent à des idéaux premiers
de l'anneau local de $X$ en $\xi_0$, d'après
Schémas, lemme \ref{schemes-lemma-specialize-points}.
On voit donc que la dimension de $X$ est la borne supérieure des dimensions
de ses anneaux locaux. En particulier, $\dim_x(X) \geq \dim(\mathcal{O}_{X, x})$,
puisque $\dim_x(X)$ est le minimum des dimensions des voisinages ouverts de
$x$. Ainsi $\sup_{x \in X} \dim_x(X) \geq \dim(X)$. D'autre part,
il est clair que $\sup_{x \in X} \dim_x(X) \leq \dim(X)$,
car $\dim(U) \leq \dim(X)$ pour tout ouvert de $X$.
\end{proof}

\begin{lemma}
\label{lemma-codimension-local-ring}
Soient $X$ un schéma et $Y \subset X$ un sous-ensemble fermé
irréductible. Soit $\xi \in Y$ son point générique. Alors
$$
\text{codim}(Y, X) = \dim(\mathcal{O}_{X, \xi})
$$
où la codimension est celle de
Topologie, définition \ref{topology-definition-codimension}.
\end{lemma}

\begin{proof}
D'après Topologie, lemme \ref{topology-lemma-codimension-at-generic-point},
on peut remplacer $X$ par un voisinage ouvert affine de $\xi$. Dans ce
cas, le résultat découle aisément de
Algèbre, lemme \ref{algebra-lemma-irreducible-components-containing-x}.
\end{proof}

\begin{lemma}
\label{lemma-generic-point}
Soient $X$ un schéma et $x \in X$. Alors $x$ est un point générique
d'une composante irréductible de $X$ si et seulement si $\dim(\mathcal{O}_{X, x}) = 0$.
\end{lemma}

\begin{proof}
Cela résulte, par exemple, du lemme \ref{lemma-codimension-local-ring}.
\end{proof}

\begin{lemma}
\label{lemma-locally-Noetherian-dimension-0}
Un schéma localement noethérien de dimension $0$ est une réunion
disjointe de spectres d'anneaux locaux artiniens.
\end{lemma}

\begin{proof}
Un anneau noethérien de dimension $0$ est un produit fini d'anneaux locaux
artiniens, voir
Algèbre, proposition \ref{algebra-proposition-dimension-zero-ring}.
L'espace topologique sous-jacent à un ouvert affine d'un schéma localement noethérien $X$ de dimension
$0$ est donc discret. Il en résulte que
la topologie de $X$ est discrète. Le lemme se déduit aisément de ces
remarques.
\end{proof}

\begin{lemma}
\label{lemma-dimension-zero}
\begin{reference}
Courriel d'Ofer Gabber du 4 juin 2016
\end{reference}
Soit $X$ un schéma de dimension zéro. Les assertions suivantes sont équivalentes :
\begin{enumerate}
\item $X$ est quasi-séparé ;
\item $X$ est séparé ;
\item $X$ est un espace topologique séparé ;
\item tout ouvert affine est fermé.
\end{enumerate}
Dans ce cas, les composantes connexes de $X$ sont des points et tout
ouvert quasi-compact de $X$ est affine. En particulier, si $X$
est quasi-compact, $X$ est affine.
\end{lemma}

\begin{proof}
Comme $X$ est de dimension zéro, pour tout ouvert affine
$U \subset X$, l'espace $U$ est profini et possède
diverses autres propriétés que nous utiliserons librement ci-dessous, voir
Algèbre, lemme \ref{algebra-lemma-ring-with-only-minimal-primes}.
Choisissons un recouvrement ouvert affine $X = \bigcup U_i$.

\medskip\noindent
Si (4) est satisfaite, $U_i \cap U_j$ est un fermé de
$U_i$, donc est quasi-compact. Par conséquent, $X$ est quasi-séparé,
d'après Schémas, lemme \ref{schemes-lemma-characterize-quasi-separated},
et (1) est satisfaite.

\medskip\noindent
Si (1) est satisfaite, $U_i \cap U_j$ est un ouvert quasi-compact
de $U_i$, donc est fermé dans $U_i$. Alors $U_i \cap U_j \to U_i$
est une immersion ouverte d'image fermée, et donc une
immersion fermée. En particulier, $U_i \cap U_j$ est affine
et $\mathcal{O}(U_i) \to \mathcal{O}_X(U_i \cap U_j)$ est surjectif.
Ainsi $X$ est séparé
d'après Schémas, lemme \ref{schemes-lemma-characterize-separated},
et (2) est satisfaite.

\medskip\noindent
Supposons (2) et soient $x, y \in X$. Supposons $x \in U_i$. Si $y \in U_i$
également, on peut trouver des voisinages ouverts disjoints de $x$ et de $y$,
car $U_i$ est un espace topologique séparé. Supposons $y \not \in U_i$ et $y \in U_j$.
Alors $y \not \in U_i \cap U_j$, qui est un ouvert affine de $U_j$,
donc fermé dans $U_j$. On peut ainsi trouver un voisinage ouvert
de $y$ ne rencontrant pas $U_i$, et l'on en déduit que $X$ est un espace topologique séparé ;
donc (3) est satisfaite.

\medskip\noindent
Supposons (3). Soit $U \subset X$ un ouvert affine.
Alors $U$ est fermé dans $X$ d'après Topologie, lemme
\ref{topology-lemma-quasi-compact-in-Hausdorff}.
Cela démontre (4).

\medskip\noindent
Supposons que $X$ satisfasse les conditions équivalentes (1) -- (4).
Démontrons les dernières assertions du lemme. Soient $x, y \in X$
avec $x \not = y$. Puisque $y$ ne se spécialise pas en $x$, on peut
choisir un ouvert affine $U \subset X$ avec $x \in U$ et $y \not \in U$.
Alors $X = U \amalg (X \setminus U)$ est une décomposition
en parties ouvertes et fermées, ce qui montre que $x$ et $y$ n'appartiennent pas
à la même composante connexe de $X$. Supposons maintenant que $U \subset X$
soit un ouvert quasi-compact. Écrivons $U = U_1 \cup \ldots \cup U_n$ comme réunion
d'ouverts affines. Montrons par récurrence sur $n$ que $U$ est affine.
On se ramène immédiatement au cas $n = 2$. Dans ce cas, on a
$U = (U_1 \setminus U_2) \amalg (U_1 \cap U_2) \amalg (U_2 \setminus U_1)$,
et les arguments précédents montrent que chacune de ces parties est affine.
\end{proof}

\begin{lemma}
\label{lemma-isolated-dim-0-locally-noetherian}
Soit $x$ un point d'un schéma localement noethérien $X$.
Alors $\dim_x(X) = 0$ si et seulement si $x$ est un point isolé de $X$.
\end{lemma}

\begin{proof}
Si $x$ est un point isolé, $\{x\}$ est un ouvert de $X$
(Topologie, définition \ref{topology-definition-isolated-point}),
et donc $\dim_x(X) = \dim(\{x\}) = 0$ par définition. Réciproquement, si
$\dim_x(X) = 0$, il existe un voisinage ouvert $U \subset X$
de $x$ tel que $\dim(U) = 0$
(Topologie, définition \ref{topology-definition-Krull}).
D'après le lemme \ref{lemma-locally-Noetherian-dimension-0},
la topologie de $U$ est discrète, et donc
$x$ est un point isolé.
\end{proof}

\section{Schémas caténaires}
\label{section-catenary}

\noindent
Rappelons qu'un espace topologique $X$ est dit {\it caténaire} si,
pour tout couple de fermés irréductibles $T \subset T'$,
il existe une chaîne maximale de fermés irréductibles
$$
T = T_0 \subset T_1 \subset \ldots \subset T_e = T'
$$
et si toutes ces chaînes ont la même longueur. Voir
Topologie, définition \ref{topology-definition-catenary}.

\begin{definition}
\label{definition-catenary}
Soit $S$ un schéma. On dit que $S$ est {\it caténaire} si
l'espace topologique sous-jacent à $S$ est caténaire.
\end{definition}

\noindent
Rappelons qu'un anneau $A$ est dit {\it caténaire} si,
pour tout couple d'idéaux premiers $\mathfrak p \subset \mathfrak q$,
il existe une chaîne maximale d'idéaux premiers
$$
\mathfrak p =
\mathfrak p_0 \subset \ldots \subset \mathfrak p_e
= \mathfrak q
$$
et si toutes ces chaînes ont la même longueur. Voir
Algèbre, définition \ref{algebra-definition-catenary}.

\begin{lemma}
\label{lemma-catenary-local}
Soit $S$ un schéma. Les assertions suivantes sont équivalentes :
\begin{enumerate}
\item $S$ est caténaire ;
\item il existe un recouvrement ouvert de $S$ dont tous les membres sont
des schémas caténaires ;
\item pour tout ouvert affine $\Spec(R) = U \subset S$, l'anneau
$R$ est caténaire ;
\item il existe un recouvrement ouvert affine $S = \bigcup U_i$ tel
que chaque $U_i$ soit le spectre d'un anneau caténaire.
\end{enumerate}
De plus, dans ce cas, tout sous-schéma localement fermé de $S$ est lui aussi
caténaire.
\end{lemma}

\begin{proof}
Combiner Topologie, lemme \ref{topology-lemma-catenary}, et
Algèbre, lemme \ref{algebra-lemma-catenary}.
\end{proof}

\begin{lemma}
\label{lemma-catenary-dimension-function}
Soit $S$ un schéma localement noethérien.
Les assertions suivantes sont équivalentes :
\begin{enumerate}
\item $S$ est caténaire ;
\item il existe, localement pour la topologie de Zariski, une fonction de dimension
sur $S$ (voir Topologie, définition \ref{topology-definition-dimension-function}).
\end{enumerate}
\end{lemma}

\begin{proof}
Cela résulte des résultats suivants :
Topologie, lemmes
\ref{topology-lemma-catenary},
\ref{topology-lemma-dimension-function-catenary} et
\ref{topology-lemma-locally-dimension-function},
Schémas, lemme \ref{schemes-lemma-scheme-sober},
et enfin le lemme \ref{lemma-Noetherian-topology}.
\end{proof}

\noindent
Un schéma est caténaire si et seulement si ses anneaux
locaux sont caténaires.

\begin{lemma}
\label{lemma-catenary-local-rings-catenary}
Soit $X$ un schéma. Les assertions suivantes sont équivalentes :
\begin{enumerate}
\item $X$ est caténaire ;
\item pour tout $x \in X$, l'anneau local $\mathcal{O}_{X, x}$ est
caténaire.
\end{enumerate}
\end{lemma}

\begin{proof}
Supposons $X$ caténaire. Soit $x \in X$. D'après le lemme \ref{lemma-catenary-local},
on peut remplacer $X$ par un voisinage ouvert affine de $x$ ; alors
$\Gamma(X, \mathcal{O}_X)$ est un anneau caténaire. D'après
Algèbre, lemme \ref{algebra-lemma-localization-catenary}, toute
localisation d'un anneau caténaire est
caténaire. Donc $\mathcal{O}_{X, x}$ est caténaire.

\medskip\noindent
Réciproquement, supposons tous les anneaux locaux de $X$ caténaires.
Soit $Y \subset Y'$ une inclusion de fermés
irréductibles de $X$. Soit $\xi \in Y$ le point générique.
Soit $\mathfrak p \subset \mathcal{O}_{X, \xi}$ l'idéal premier
correspondant au point générique de $Y'$, voir
Schémas, lemme \ref{schemes-lemma-specialize-points}. D'après ce même
lemme, les fermés irréductibles de $X$ compris entre $Y$ et $Y'$
correspondent aux idéaux premiers $\mathfrak q \subset \mathcal{O}_{X, \xi}$
tels que $\mathfrak p \subset \mathfrak q \subset \mathfrak m_{\xi}$.
Toutes leurs chaînes maximales sont donc finies et de même
longueur, puisque $\mathcal{O}_{X, \xi}$ est un anneau caténaire.
\end{proof}

\section{Conditions de Serre}
\label{section-Rk}

\noindent
Voici deux notions techniques souvent utiles.
Voir aussi Cohomologie des schémas, section \ref{coherent-section-depth}.

\begin{definition}
\label{definition-Rk}
Soit $X$ un schéma localement noethérien. Soit $k \geq 0$.
\begin{enumerate}
\item On dit que $X$ est {\it régulier en codimension $k$},
ou que $X$ possède la propriété {\it $(R_k)$}, si, pour tout $x \in X$,
on a
$$
\dim(\mathcal{O}_{X, x}) \leq k
\Rightarrow
\mathcal{O}_{X, x}\text{ est régulier}
$$
\item On dit que $X$ possède la propriété {\it $(S_k)$} si, pour tout $x \in X$, on a
$\text{depth}(\mathcal{O}_{X, x}) \geq \min(k, \dim(\mathcal{O}_{X, x}))$.
\end{enumerate}
\end{definition}

\noindent
L'expression « régulier en codimension $k$ » a un sens, car nous avons vu,
dans la section \ref{section-catenary}, que, si $Y \subset X$ est un fermé
irréductible de point générique $x$, alors
$\dim(\mathcal{O}_{X, x}) = \text{codim}(Y, X)$. Par exemple, la condition
$(R_0)$ signifie que, pour tout point générique $\eta \in X$ d'une composante
irréductible de $X$, l'anneau local $\mathcal{O}_{X, \eta}$ est un corps.
Mais, pour un schéma noethérien quelconque, le lieu régulier
de $X$ peut mal se comporter ; il faut donc être prudent.

\begin{lemma}
\label{lemma-scheme-regular-iff-all-Rk}
Soit $X$ un schéma localement noethérien.
Alors $X$ est régulier si et seulement si $X$ satisfait $(R_k)$ pour tout $k \geq 0$.
\end{lemma}

\begin{proof}
Cela résulte du lemme \ref{lemma-characterize-regular} et des définitions.
\end{proof}

\begin{lemma}
\label{lemma-scheme-CM-iff-all-Sk}
Soit $X$ un schéma localement noethérien.
Alors $X$ est de Cohen--Macaulay si et seulement si $X$ satisfait $(S_k)$ pour tout $k \geq 0$.
\end{lemma}

\begin{proof}
Le lemme \ref{lemma-characterize-Cohen-Macaulay}
permet de se ramener aux anneaux locaux.
Le résultat vient donc de ce qu'un anneau local noethérien est de Cohen--Macaulay
si et seulement si sa profondeur est égale à sa dimension.
\end{proof}

\begin{lemma}
\label{lemma-criterion-reduced}
Soit $X$ un schéma localement noethérien.
Alors $X$ est réduit si et seulement si $X$ possède les propriétés $(S_1)$ et $(R_0)$.
\end{lemma}

\begin{proof}
C'est le lemme \ref{algebra-lemma-criterion-reduced} du chapitre Algèbre.
\end{proof}

\begin{lemma}
\label{lemma-criterion-normal}
Soit $X$ un schéma localement noethérien.
Alors $X$ est normal si et seulement si $X$ possède les propriétés $(S_2)$ et $(R_1)$.
\end{lemma}

\begin{proof}
C'est le lemme \ref{algebra-lemma-criterion-normal} du chapitre Algèbre.
\end{proof}

\begin{lemma}
\label{lemma-normal-dimension-1-regular}
Soit $X$ un schéma localement noethérien, normal et
de dimension $\leq 1$. Alors $X$ est régulier.
\end{lemma}

\begin{proof}
Cela résulte du lemme \ref{lemma-criterion-normal} et des définitions.
\end{proof}

\begin{lemma}
\label{lemma-normal-dimension-2-Cohen-Macaulay}
Soit $X$ un schéma localement noethérien, normal et
de dimension $\leq 2$. Alors $X$ est de Cohen--Macaulay.
\end{lemma}

\begin{proof}
Cela résulte du lemme \ref{lemma-criterion-normal} et des définitions.
\end{proof}

\section{Schémas japonais et schémas de Nagata}
\label{section-nagata}

\noindent
Les notions considérées dans cette section ne font pas l'objet de définitions explicites bien mises en évidence dans EGA.
Un « schéma universellement japonais » est mentionné et défini dans
\cite[IV Corollary 5.11.4]{EGA}. Un « schéma japonais » est mentionné dans
\cite[IV Remark 10.4.14 (ii)]{EGA}, mais sans définition.
La notion de schéma de Nagata (définie ci-dessous) apparaît à quelques
endroits dans la littérature (voir, par exemple, \cite[Definition 8.2.30]{Liu} et
\cite[Page 142]{Greco}).

\medskip\noindent
Rappelons brièvement qu'un anneau intègre $R$ est dit {\it japonais} si la fermeture
intégrale de $R$ dans toute extension finie de son corps des fractions est finie sur
$R$. Un anneau $R$ est dit {\it universellement japonais} si, pour tout homomorphisme
d'anneaux de type fini $R \to S$ avec $S$ intègre, $S$ est japonais. Un anneau $R$ est dit
{\it de Nagata} s'il est noethérien et si $R/\mathfrak p$ est japonais pour tout
idéal premier $\mathfrak p$ de $R$.

\begin{definition}
\label{definition-nagata}
Soit $X$ un schéma.
\begin{enumerate}
\item Supposons $X$ intègre. On dit que $X$ est {\it japonais}
si, pour tout $x \in X$, il existe un
voisinage ouvert affine $x \in U \subset X$ tel que l'anneau
$\mathcal{O}_X(U)$ soit japonais (voir
Algèbre, définition \ref{algebra-definition-N}).
\item On dit que $X$ est {\it universellement japonais} si, pour tout $x \in X$,
il existe un voisinage ouvert affine $x \in U \subset X$ tel que
l'anneau $\mathcal{O}_X(U)$ soit universellement japonais (voir
Algèbre, définition \ref{algebra-definition-nagata}).
\item On dit que $X$ est {\it de Nagata} si, pour tout $x \in X$, il existe un
voisinage ouvert affine $x \in U \subset X$ tel que l'anneau
$\mathcal{O}_X(U)$ soit de Nagata (voir
Algèbre, définition \ref{algebra-definition-nagata}).
\end{enumerate}
\end{definition}

\noindent
Être de Nagata équivaut à être localement noethérien
et universellement japonais, voir le
lemme \ref{lemma-nagata-universally-Japanese}.

\begin{remark}
\label{remark-non-integral-Japanese}
Dans \cite{Hoobler-finite}, un schéma (localement noethérien) $X$ est dit
japonais si, pour tout $x \in X$ et tout idéal premier associé $\mathfrak p$
de $\mathcal{O}_{X, x}$, l'anneau $\mathcal{O}_{X, x}/\mathfrak p$ est
japonais. Nous n'utilisons pas cette définition, car il existe un anneau
intègre noethérien de dimension un dont les anneaux locaux sont excellents (en particulier
japonais), mais dont la normalisation n'est pas finie. Voir
\cite[Example 1]{Hochster-loci}, \cite{Heinzer-Levy} ou
\cite[Expos\'e XIX]{Traveaux}.
On pourrait toutefois contourner ce problème en disant qu'un schéma
$X$ est japonais si, pour tout ouvert affine $\Spec(A) \subset X$, l'anneau
$A/\mathfrak p$ est japonais pour tout idéal premier associé $\mathfrak p$ de $A$.
\end{remark}

\begin{lemma}
\label{lemma-nagata-locally-Noetherian}
Un schéma de Nagata est localement noethérien.
\end{lemma}

\begin{proof}
En effet, un anneau de Nagata est noethérien par définition.
\end{proof}

\begin{lemma}
\label{lemma-locally-Japanese}
Soit $X$ un schéma intègre. Les assertions suivantes sont équivalentes :
\begin{enumerate}
\item Le schéma $X$ est japonais.
\item Pour tout ouvert affine $U \subset X$, l'anneau intègre $\mathcal{O}_X(U)$
est japonais.
\item Il existe un recouvrement ouvert affine $X = \bigcup U_i$
tel que chaque $\mathcal{O}_X(U_i)$ soit japonais.
\item Il existe un recouvrement ouvert $X = \bigcup X_j$
tel que chaque sous-schéma ouvert $X_j$ soit japonais.
\end{enumerate}
De plus, si $X$ est japonais, tout sous-schéma ouvert
est japonais.
\end{lemma}

\begin{proof}
Cela résulte du lemme \ref{lemma-locally-P} et des lemmes suivants du chapitre Algèbre : \ref{algebra-lemma-localize-N} et
\ref{algebra-lemma-Japanese-local}.
\end{proof}

\begin{lemma}
\label{lemma-locally-universally-Japanese}
Soit $X$ un schéma. Les assertions suivantes sont équivalentes :
\begin{enumerate}
\item Le schéma $X$ est universellement japonais.
\item Pour tout ouvert affine $U \subset X$, l'anneau $\mathcal{O}_X(U)$
est universellement japonais.
\item Il existe un recouvrement ouvert affine $X = \bigcup U_i$
tel que chaque $\mathcal{O}_X(U_i)$ soit universellement japonais.
\item Il existe un recouvrement ouvert $X = \bigcup X_j$
tel que chaque sous-schéma ouvert $X_j$ soit universellement japonais.
\end{enumerate}
De plus, si $X$ est universellement japonais, tout sous-schéma ouvert
est universellement japonais.
\end{lemma}

\begin{proof}
Cela résulte du lemme \ref{lemma-locally-P} et des lemmes suivants du chapitre Algèbre : \ref{algebra-lemma-universally-japanese} et
\ref{algebra-lemma-nagata-local}.
\end{proof}

\begin{lemma}
\label{lemma-locally-nagata}
Soit $X$ un schéma. Les assertions suivantes sont équivalentes :
\begin{enumerate}
\item Le schéma $X$ est de Nagata.
\item Pour tout ouvert affine $U \subset X$, l'anneau $\mathcal{O}_X(U)$
est de Nagata.
\item Il existe un recouvrement ouvert affine $X = \bigcup U_i$
tel que chaque $\mathcal{O}_X(U_i)$ soit de Nagata.
\item Il existe un recouvrement ouvert $X = \bigcup X_j$
tel que chaque sous-schéma ouvert $X_j$ soit de Nagata.
\end{enumerate}
De plus, si $X$ est de Nagata, tout sous-schéma ouvert est de Nagata.
\end{lemma}

\begin{proof}
Cela résulte du lemme \ref{lemma-locally-P} et des lemmes suivants du chapitre Algèbre : \ref{algebra-lemma-nagata-localize} et
\ref{algebra-lemma-nagata-local}.
\end{proof}

\begin{lemma}
\label{lemma-characterize-nagata}
Soit $X$ un schéma localement noethérien.
Alors $X$ est de Nagata si et seulement si tout sous-schéma fermé intègre
$Z \subset X$ est japonais.
\end{lemma}

\begin{proof}
Supposons $X$ de Nagata. Soit $Z \subset X$ un sous-schéma fermé intègre.
Soit $z \in Z$.
Soit $\Spec(A) = U \subset X$ un ouvert affine contenant $z$
tel que $A$ soit de Nagata. Alors
$Z \cap U \cong \Spec(A/\mathfrak p)$ pour un certain idéal premier $\mathfrak p$,
voir Schémas, lemme \ref{schemes-lemma-closed-subspace-scheme} (et la
définition \ref{definition-integral}). D'après
Algèbre, définition \ref{algebra-definition-nagata},
$A/\mathfrak p$ est japonais. Donc $Z$ est japonais par définition.

\medskip\noindent
Supposons tout sous-schéma fermé intègre de $X$ japonais.
Soit $\Spec(A) = U \subset X$ un ouvert affine quelconque.
Comme $X$ est localement noethérien, $A$ est noethérien
(lemme \ref{lemma-locally-Noetherian}). Soit $\mathfrak p \subset A$
un idéal premier. Il faut montrer que $A/\mathfrak p$ est japonais.
Soit $T \subset U$ le fermé $V(\mathfrak p) \subset \Spec(A)$.
Soit $\overline{T} \subset X$ son adhérence. Alors $\overline{T}$ est
irréductible, étant l'adhérence d'un sous-ensemble irréductible. Le sous-schéma
fermé réduit défini par $\overline{T}$ est donc un sous-schéma fermé intègre
(encore noté $\overline{T}$), voir
Schémas, lemme \ref{schemes-lemma-reduced-closed-subscheme}.
Autrement dit, $\Spec(A/\mathfrak p)$ est un ouvert
affine d'un sous-schéma fermé intègre de $X$. Ce sous-schéma est japonais
par hypothèse ; le lemme \ref{lemma-locally-Japanese} montre donc que
$A/\mathfrak p$ est japonais.
\end{proof}

\begin{lemma}
\label{lemma-nagata-universally-Japanese}
Soit $X$ un schéma.
Les assertions suivantes sont équivalentes :
\begin{enumerate}
\item $X$ est de Nagata ;
\item $X$ est localement noethérien et universellement japonais.
\end{enumerate}
\end{lemma}

\begin{proof}
C'est
Algèbre, proposition \ref{algebra-proposition-nagata-universally-japanese}.
\end{proof}

\noindent
Cette discussion sera poursuivie dans
Morphismes, section \ref{morphisms-section-nagata}.

\section{G-schémas}
\label{section-G-scheme}

\noindent
Rappelons qu'un anneau $R$ est un G-anneau si $R$ est noethérien et si,
pour tout idéal premier $\mathfrak p$ de $R$, l'homomorphisme d'anneaux
$R_\mathfrak p \to (R_\mathfrak p)^\wedge$ est régulier.

\begin{definition}
\label{definition-G-scheme}
Soit $X$ un schéma. On dit que $X$ est un {\it G-schéma}\footnote{Cette terminologie
n'est peut-être pas usuelle. Si $G$ est un groupe fini ou un schéma en groupes,
un $G$-schéma désigne parfois un schéma muni d'une
action de $G$.} si, pour tout $x \in X$, il existe un voisinage
ouvert affine $x \in U \subset X$ tel que l'anneau
$\mathcal{O}_X(U)$ soit un G-anneau (voir
Compléments sur l'algèbre, définition \ref{more-algebra-definition-G-ring}).
\end{definition}

\begin{lemma}
\label{lemma-characterize-G-scheme}
Soit $X$ un schéma. Les assertions suivantes sont équivalentes :
\begin{enumerate}
\item $X$ est un G-schéma ;
\item $X$ est localement noethérien et, pour tout $x \in X$, l'homomorphisme
d'anneaux $\mathcal{O}_{X, x} \to \mathcal{O}_{X, x}^\wedge$ est régulier ;
\item $X$ est localement noethérien et, pour tout point fermé $x \in X$,
l'homomorphisme d'anneaux $\mathcal{O}_{X, x} \to \mathcal{O}_{X, x}^\wedge$ est régulier.
\end{enumerate}
\end{lemma}

\begin{proof}
Il est clair que (1) entraîne (2) et que (2) entraîne (3).
Supposons (3). Soit $x \in X$ un point quelconque. D'après le
lemme \ref{lemma-locally-Noetherian-closed-point},
il existe une spécialisation $x \leadsto x'$ avec $x'$ fermé
dans $X$. Alors $\mathcal{O}_{X, x'}$ est un G-anneau d'après
Compléments sur l'algèbre, lemme \ref{more-algebra-lemma-check-G-ring-maximal-ideals}.
Comme $\mathcal{O}_{X, x}$ est la localisation de $\mathcal{O}_{X, x'}$
en un idéal premier (Schémas, lemme \ref{schemes-lemma-specialize-points}),
l'homomorphisme d'anneaux $\mathcal{O}_{X, x} \to \mathcal{O}_{X, x}^\wedge$
est régulier. Le point $x \in X$ étant quelconque, on en déduit
que, pour tout ouvert affine $U \subset X$, l'anneau noethérien
(lemme \ref{lemma-locally-Noetherian}) $\mathcal{O}_X(U)$ est un G-anneau.
Cela démontre (1).
\end{proof}

\begin{lemma}
\label{lemma-locally-G-scheme}
Soit $X$ un schéma. Les assertions suivantes sont équivalentes :
\begin{enumerate}
\item Le schéma $X$ est un G-schéma.
\item Pour tout ouvert affine $U \subset X$, l'anneau $\mathcal{O}_X(U)$
est un G-anneau.
\item Il existe un recouvrement ouvert affine $X = \bigcup U_i$ tel que
chaque $\mathcal{O}_X(U_i)$ soit un G-anneau.
\item Il existe un recouvrement ouvert $X = \bigcup X_j$
tel que chaque sous-schéma ouvert $X_j$ soit un G-schéma.
\end{enumerate}
De plus, si $X$ est un G-schéma, tout sous-schéma ouvert est un
G-schéma.
\end{lemma}

\begin{proof}
Combiner les lemmes \ref{lemma-locally-Noetherian}
et \ref{lemma-characterize-G-scheme}.
\end{proof}

\section{Le lieu singulier}
\label{section-singular-locus}

\noindent
Voici la définition.

\begin{definition}
\label{definition-singular-locus}
Soit $X$ un schéma localement noethérien. Le {\it lieu régulier}
$\text{Reg}(X)$ de $X$ est l'ensemble des $x \in X$ tels que $\mathcal{O}_{X, x}$
soit un anneau local régulier. Le {\it lieu singulier} $\text{Sing}(X)$ est le
complémentaire $X \setminus \text{Reg}(X)$, c'est-à-dire l'ensemble des points $x \in X$
tels que $\mathcal{O}_{X, x}$ ne soit pas un anneau local régulier.
\end{definition}

\noindent
Le lieu régulier d'un schéma localement noethérien est stable par
générisation, voir la discussion précédant
Algèbre, définition \ref{algebra-definition-regular}.
Toutefois, pour un schéma localement noethérien quelconque, le lieu régulier
n'est pas nécessairement ouvert. Dans
Compléments sur l'algèbre, section \ref{more-algebra-section-singular-locus},
on trouvera quelques critères assurant qu'il l'est.
Nous reviendrons sur cette question dans
Morphismes, section \ref{morphisms-section-singular-locus}.

\section{Irréductibilité locale}
\label{section-unibranch}

\noindent
Rappelons que, dans Compléments sur l'algèbre, section \ref{more-algebra-section-unibranch},
nous avons introduit la notion d'anneau local (géométriquement) unibranche.

\begin{definition}
\label{definition-unibranch}
\begin{reference}
\cite[Chapter IV (6.15.1)]{EGA4}
\end{reference}
Soient $X$ un schéma et $x \in X$. On dit que $X$ est {\it unibranche en $x$}
si l'anneau local $\mathcal{O}_{X, x}$ est unibranche. On dit que $X$ est
{\it géométriquement unibranche en $x$}
si l'anneau local $\mathcal{O}_{X, x}$ est géométriquement unibranche.
On dit que $X$ est {\it unibranche} si $X$ est unibranche en chacun de ses points.
On dit que $X$ est {\it géométriquement unibranche} si $X$ est
géométriquement unibranche en chacun de ses points.
\end{definition}

\noindent
Il peut effectivement arriver qu'un anneau local $A$ soit géométriquement unibranche
(au sens de
Compléments sur l'algèbre, définition \ref{more-algebra-definition-unibranch}),
sans que le schéma $\Spec(A)$ soit géométriquement unibranche au sens
de la définition \ref{definition-unibranch}. C'est le cas, par exemple,
si $A$ est l'anneau local au sommet du cône sur une courbe plane
irréductible ayant un point double ordinaire (un nœud).

\begin{lemma}
\label{lemma-normal-geometrically-unibranch}
Un schéma normal est géométriquement unibranche.
\end{lemma}

\begin{proof}
Cela résulte des définitions. En effet, un schéma
est normal si ses anneaux locaux sont des anneaux intègres normaux. Il résulte immédiatement
de Compléments sur l'algèbre, définition \ref{more-algebra-definition-unibranch},
qu'un anneau local intègre normal est géométriquement unibranche.
\end{proof}

\begin{lemma}
\label{lemma-geometrically-unibranch}
\begin{reference}
Comparer avec \cite[Proposition 2.3]{Etale-coverings}
\end{reference}
Soit $X$ un schéma noethérien. Les assertions suivantes sont équivalentes :
\begin{enumerate}
\item $X$ est géométriquement unibranche (définition \ref{definition-unibranch}) ;
\item pour tout point $x \in X$ qui n'est pas le point générique d'une
composante irréductible de $X$, le spectre épointé de
l'hensélisé strict $\mathcal{O}_{X, x}^{sh}$ est connexe.
\end{enumerate}
\end{lemma}

\begin{proof}
Compléments sur l'algèbre, lemme \ref{more-algebra-lemma-geometrically-unibranch},
montre que (1) entraîne que les spectres épointés de (2) sont
irréductibles, donc connexes.

\medskip\noindent
Supposons (2). Soit $x \in X$. Il faut montrer que $\mathcal{O}_{X, x}$
est géométriquement unibranche. Par récurrence sur $\dim(\mathcal{O}_{X, x})$,
on peut supposer le résultat vrai pour toute générisation non triviale de $x$.
On peut remplacer $X$ par $\Spec(\mathcal{O}_{X, x})$. Autrement dit,
on peut supposer que $X = \Spec(A)$, avec $A$ local, et que
$A_\mathfrak p$ est géométriquement unibranche pour tout idéal premier
non maximal $\mathfrak p \subset A$.

\medskip\noindent
Soit $A^{sh}$ l'hensélisé strict de $A$. Si
$\mathfrak q \subset A^{sh}$ est un idéal premier au-dessus de $\mathfrak p \subset A$,
alors $A_\mathfrak p \to A^{sh}_\mathfrak q$ est une
colimite filtrante d'algèbres \'etales. Les hensélisés stricts de
$A_\mathfrak p$ et de $A^{sh}_\mathfrak q$ sont donc isomorphes.
Ainsi, d'après Compléments sur l'algèbre, lemme \ref{more-algebra-lemma-geometrically-unibranch},
$A^{sh}_\mathfrak q$
possède un unique idéal premier minimal pour tout idéal premier non maximal $\mathfrak q$ de
$A^{sh}$.

\medskip\noindent
Soient $\mathfrak q_1, \ldots, \mathfrak q_r$ les idéaux premiers minimaux
de $A^{sh}$. Il faut montrer que $r = 1$. D'après ce qui précède,
on a $V(\mathfrak q_1) \cap V(\mathfrak q_j) = \{\mathfrak m^{sh}\}$
pour $j = 2, \ldots, r$. Ainsi $V(\mathfrak q_1) \setminus \{\mathfrak m^{sh}\}$
est un ouvert fermé du spectre épointé de $A^{sh}$,
ce qui contredit l'hypothèse selon laquelle ce spectre épointé
est connexe, à moins que $r = 1$.
\end{proof}

\begin{definition}
\label{definition-number-of-branches}
Soient $X$ un schéma et $x \in X$. Le {\it nombre de branches de $X$
en $x$} est le nombre de branches de l'anneau local $\mathcal{O}_{X, x}$
au sens de
Compléments sur l'algèbre, définition \ref{more-algebra-definition-number-of-branches}.
Le {\it nombre de branches géométriques de $X$ en $x$} est le nombre de
branches géométriques de l'anneau local $\mathcal{O}_{X, x}$ au sens de
Compléments sur l'algèbre, définition \ref{more-algebra-definition-number-of-branches}.
\end{definition}

\noindent
On souhaite souvent comparer ce nombre à celui des branches de l'anneau local
complété, mais cette comparaison n'est généralement pas immédiate ; on trouvera quelques résultats
à ce sujet dans Compléments sur l'algèbre, section
\ref{more-algebra-section-branches-completion}.

\begin{lemma}
\label{lemma-number-of-branches-irreducible-components}
Soient $X$ un schéma et $x \in X$. Soient $X_i$, $i \in I$, les
composantes irréductibles de $X$ passant par $x$.
Alors le nombre de branches (géométriques) de $X$ en $x$
est la somme, sur $i \in I$, des nombres de branches (géométriques)
de $X_i$ en $x$.
\end{lemma}

\begin{proof}
On considère les $X_i$ comme des sous-schémas fermés intègres de $X$, voir
Schémas, définition \ref{schemes-definition-reduced-induced-scheme}, et le
lemme \ref{lemma-characterize-integral}.
Remarquons que le nombre de branches (géométriques) de $X_i$ en $x$
est au moins $1$ pour tout $i$ (essentiellement par définition).
Rappelons que les $X_i$ correspondent $1$-à-$1$ aux idéaux premiers
minimaux $\mathfrak p_i \subset \mathcal{O}_{X, x}$, voir
Algèbre, lemme \ref{algebra-lemma-irreducible-components-containing-x}.
Ainsi, si $I$ est infini, $\mathcal{O}_{X, x}$ possède une infinité
d'idéaux premiers minimaux, et donc $\mathcal{O}_{X, x}^h$
et $\mathcal{O}_{X, x}^{sh}$ possèdent tous deux une infinité d'idéaux premiers
minimaux (combiner Algèbre, lemmes
\ref{algebra-lemma-injective-minimal-primes-in-image} et
\ref{algebra-lemma-minimal-prime-image-minimal-prime}, avec
l'injectivité des homomorphismes
$\mathcal{O}_{X, x} \to  \mathcal{O}_{X, x}^h \to \mathcal{O}_{X, x}^{sh}$).
Dans ce cas, le nombre de branches (géométriques) de $X$ en $x$
est, par définition, $\infty$, tout comme la somme considérée.
On peut donc supposer $I$ fini.
Soit $A'$ la fermeture intégrale de $\mathcal{O}_{X, x}$
dans l'anneau total des fractions $Q$ de $(\mathcal{O}_{X, x})_{red}$.
Soit $A'_i$ la fermeture intégrale de $\mathcal{O}_{X, x}/\mathfrak p_i$
dans l'anneau total des fractions $Q_i$ de $\mathcal{O}_{X, x}/\mathfrak p_i$.
D'après Algèbre, lemme \ref{algebra-lemma-total-ring-fractions-no-embedded-points},
on a $Q = \prod_{i \in I} Q_i$. Ainsi $A' = \prod A'_i$.
L'égalité du lemme résulte alors de
Compléments sur l'algèbre, lemme \ref{more-algebra-lemma-number-of-branches-1},
qui exprime le nombre de branches (géométriques) en fonction
des idéaux maximaux de $A'$.
\end{proof}

\begin{lemma}
\label{lemma-number-of-branches-1}
Soient $X$ un schéma et $x \in X$.
\begin{enumerate}
\item Le nombre de branches de $X$ en $x$ est $1$ si et seulement si
$X$ est unibranche en $x$.
\item Le nombre de branches géométriques de $X$ en $x$ est $1$ si et seulement si
$X$ est géométriquement unibranche en $x$.
\end{enumerate}
\end{lemma}

\begin{proof}
Ce lemme résulte immédiatement des définitions et du résultat correspondant
pour les anneaux, voir Compléments sur l'algèbre, lemme
\ref{more-algebra-lemma-number-of-branches-1}.
\end{proof}
\section{Caractérisation des modules de type fini et de présentation finie}
\label{section-characterizing-finite-type-presentation}

\noindent
Soit $X$ un schéma.
Soit $\mathcal{F}$ un $\mathcal{O}_X$-module quasi-cohérent.
Le lemme suivant entraîne que $\mathcal{F}$ est de type fini
(voir Modules, définition \ref{modules-definition-finite-type})
si et seulement si $\mathcal{F}$ est,
sur tout ouvert affine $\Spec(A) = U \subset X$,
de la forme $\widetilde M$ pour un $A$-module de type fini $M$.
De même, $\mathcal{F}$ est de présentation finie
(voir Modules, définition \ref{modules-definition-finite-presentation})
si et seulement si $\mathcal{F}$ est,
sur tout ouvert affine $\Spec(A) = U \subset X$,
de la forme $\widetilde M$ pour un $A$-module de présentation finie $M$.

\begin{lemma}
\label{lemma-finite-type-module}
Soit $X = \Spec(R)$ un schéma affine.
Le faisceau quasi-cohérent de $\mathcal{O}_X$-modules
$\widetilde M$ est un $\mathcal{O}_X$-module de type fini
si et seulement si $M$ est un $R$-module de type fini.
\end{lemma}

\begin{proof}
Supposons que $\widetilde M$ soit un $\mathcal{O}_X$-module de type fini.
Cela signifie qu'il existe un recouvrement ouvert de $X$ tel que
la restriction de $\widetilde M$ à chacun des membres de ce recouvrement soit
engendrée par un nombre fini de sections globales.
Il existe donc aussi un recouvrement ouvert standard
$X = \bigcup_{i = 1, \ldots, n} D(f_i)$ tel que $\widetilde M|_{D(f_i)}$
soit engendré par un nombre fini de sections. Ainsi $M_{f_i}$ est de type
fini pour tout $i$. On conclut par
Algèbre, lemme \ref{algebra-lemma-cover}.
\end{proof}

\begin{lemma}
\label{lemma-finite-presentation-module}
Soit $X = \Spec(R)$ un schéma affine. Le faisceau quasi-cohérent
de $\mathcal{O}_X$-modules $\widetilde M$ est un $\mathcal{O}_X$-module de
présentation finie si et seulement si $M$ est un $R$-module de présentation finie.
\end{lemma}

\begin{proof}
Supposons que $\widetilde M$ soit un $\mathcal{O}_X$-module de présentation finie.
Le lemme \ref{lemma-finite-type-module} montre que $M$ est un $R$-module de type fini.
Choisissons une surjection $R^n \to M$ de noyau $K$. D'après
Schémas, lemme \ref{schemes-lemma-spec-sheaves},
on a une suite exacte courte
$$
0 \to \widetilde{K} \to
\bigoplus \mathcal{O}_X^{\oplus n} \to
\widetilde{M} \to 0
$$
D'après
Modules, lemme
\ref{modules-lemma-kernel-surjection-finite-free-onto-finite-presentation},
$\widetilde{K}$ est un $\mathcal{O}_X$-module de type fini.
Donc, de nouveau par le lemme \ref{lemma-finite-type-module},
$K$ est un $R$-module de type fini.
Ainsi $M$ est un $R$-module de présentation finie.
\end{proof}

\section{Sections sur les ouverts principaux}
\label{section-principal-opens}

\noindent
Voici un résultat typique de cette nature. Dans les lemmes suivants, nous utiliserons une
démonstration plus élémentaire, mais plus directe.

\begin{lemma}
\label{lemma-invert-f-sections}
\begin{slogan}
Les sections des faisceaux quasi-cohérents n'ont que des singularités méromorphes
à l'infini.
\end{slogan}
Soient $X$ un schéma et $f \in \Gamma(X, \mathcal{O}_X)$.
Notons $X_f \subset X$ l'ouvert où $f$ est inversible, voir
Schémas, lemme \ref{schemes-lemma-f-open}.
Si $X$ est quasi-compact et quasi-séparé, l'homomorphisme canonique
$$
\Gamma(X, \mathcal{O}_X)_f \longrightarrow \Gamma(X_f, \mathcal{O}_X)
$$
est un isomorphisme. De plus, si $\mathcal{F}$ est un faisceau
quasi-cohérent de $\mathcal{O}_X$-modules, l'homomorphisme
$$
\Gamma(X, \mathcal{F})_f \longrightarrow \Gamma(X_f, \mathcal{F})
$$
est un isomorphisme.
\end{lemma}

\begin{proof}
Écrivons $R = \Gamma(X, \mathcal{O}_X)$. Considérons le morphisme canonique
$$
\varphi : X \longrightarrow \Spec(R)
$$
de schémas, voir
Schémas, lemme
\ref{schemes-lemma-morphism-into-affine}.
L'image inverse de l'ouvert principal $D(f)$ du membre de droite
est alors l'ouvert $X_f$ du membre de gauche.
De plus, comme $X$ est supposé quasi-compact et quasi-séparé,
le morphisme $\varphi$ est quasi-compact et quasi-séparé,
voir Schémas, lemmes \ref{schemes-lemma-quasi-compact-affine} et
\ref{schemes-lemma-compose-after-separated}. Ainsi, d'après
Schémas, lemme \ref{schemes-lemma-push-forward-quasi-coherent},
$\varphi_*\mathcal{F}$ est quasi-cohérent.
On a donc $\varphi_*\mathcal{F} = \widetilde M$,
où $M = \Gamma(X, \mathcal{F})$, considéré comme $R$-module.
Il en résulte que
$$
\Gamma(X_f, \mathcal{F}) =
\Gamma(D(f), \varphi_*\mathcal{F}) =
\Gamma(D(f), \widetilde M) = M_f
$$
ce qui est précisément l'assertion du lemme. Le premier isomorphisme affiché
s'obtient en prenant $\mathcal{F} = \mathcal{O}_X$.
\end{proof}

\noindent
Rappelons qu'étant donnés un schéma $X$, un faisceau inversible $\mathcal{L}$
sur $X$ et un faisceau de $\mathcal{O}_X$-modules $\mathcal{F}$,
on obtient un anneau gradué
$\Gamma_*(X, \mathcal{L}) =
\bigoplus\nolimits_{n \geq 0} \Gamma(X, \mathcal{L}^{\otimes n})$
et un $\Gamma_*(X, \mathcal{L})$-module gradué
$\Gamma_*(X, \mathcal{L}, \mathcal{F}) =
\bigoplus\nolimits_{n \in \mathbf{Z}}
\Gamma(X, \mathcal{F} \otimes_{\mathcal{O}_X} \mathcal{L}^{\otimes n})$,
voir Modules, définition \ref{modules-definition-gamma-star}.
Si l'on se donne de plus une section $s \in \Gamma(X, \mathcal{L})$,
on obtient une application
\begin{equation}
\label{equation-module-invert-s}
\Gamma_*(X, \mathcal{L}, \mathcal{F})_{(s)}
\longrightarrow
\Gamma(X_s, \mathcal{F}|_{X_s})
\end{equation}
qui envoie $t/s^n$, où
$t \in \Gamma(X, \mathcal{F} \otimes_{\mathcal{O}_X} \mathcal{L}^{\otimes n})$,
sur $t|_{X_s} \otimes s|_{X_s}^{-n}$. Cela a un sens,
car $X_s \subset X$ est, par définition, l'ouvert sur lequel
$s$ possède un inverse, voir Modules, lemme \ref{modules-lemma-s-open}.

\begin{lemma}
\label{lemma-invert-s-sections}
\begin{reference}
\cite[Section 6.8]{EGA1-second}
\end{reference}
Soit $X$ un schéma. Soit $\mathcal{L}$ un faisceau inversible sur $X$.
Soit $s \in \Gamma(X, \mathcal{L})$. Soit $\mathcal{F}$ un
$\mathcal{O}_X$-module quasi-cohérent.
\begin{enumerate}
\item Si $X$ est quasi-compact, l'application (\ref{equation-module-invert-s})
est injective ;
\item si $X$ est quasi-compact et quasi-séparé,
l'application (\ref{equation-module-invert-s}) est un isomorphisme.
\end{enumerate}
En particulier, l'homomorphisme canonique
$$
\Gamma_*(X, \mathcal{L})_{(s)}
\longrightarrow
\Gamma(X_s, \mathcal{O}_X),\quad
a/s^n \longmapsto a \otimes s^{-n}
$$
est un isomorphisme si $X$ est quasi-compact et quasi-séparé.
\end{lemma}

\begin{proof}
Supposons $X$ quasi-compact. Choisissons un recouvrement ouvert affine fini
$X = U_1 \cup \ldots \cup U_m$, avec $U_j$ affine et
$\mathcal{L}|_{U_j} \cong \mathcal{O}_{U_j}$. Par cet isomorphisme,
la section $s|_{U_j}$ correspond à un élément
$f_j \in \Gamma(U_j, \mathcal{O}_{U_j})$. Alors
$X_s \cap U_j = D(f_j)$.

\medskip\noindent
Démonstration de (1). Soit $t/s^n$ un élément du noyau de
(\ref{equation-module-invert-s}). Alors $t|_{X_s} = 0$.
Donc $(t|_{U_j})|_{D(f_j)} = 0$. Par le
lemme \ref{lemma-invert-f-sections}, on obtient
$f_j^{e_j} t|_{U_j} = 0$ pour un certain
$e_j \geq 0$. Soit $e = \max(e_j)$. Alors $t \otimes s^e$
se restreint à zéro sur $U_j$ pour tout $j$, et est donc nul. Comme $t/s^n$
est égal à $t \otimes s^e/s^{n + e}$ dans
$\Gamma_*(X, \mathcal{L}, \mathcal{F})_{(s)}$, on en déduit que $t/s^n = 0$,
comme voulu.

\medskip\noindent
Démonstration de (2). Supposons $X$ quasi-compact et quasi-séparé.
Alors $U_j \cap U_{j'}$ est quasi-compact pour tout couple $j, j'$, voir
Schémas, lemme \ref{schemes-lemma-characterize-quasi-separated}.
D'après (1), on sait que (\ref{equation-module-invert-s}) est injective.
Soit $t' \in \Gamma(X_s, \mathcal{F}|_{X_s})$. Pour tout $j$, il existe un
entier $e_j \geq 0$ et une section $t'_j \in \Gamma(U_j, \mathcal{F}|_{U_j})$ tels que
$t'|_{D(f_j)}$ corresponde à $t'_j/f_j^{e_j}$
par l'isomorphisme du lemme \ref{lemma-invert-f-sections}.
Posons $e = \max(e_j)$ et
$$
t_j = f_j^{e - e_j} t'_j \otimes q_j^e \in
\Gamma(U_j,
(\mathcal{F} \otimes_{\mathcal{O}_X} \mathcal{L}^{\otimes e})|_{U_j})
$$
où $q_j \in \Gamma(U_j, \mathcal{L}|_{U_j})$ est la section trivialisante
provenant de l'isomorphisme
$\mathcal{L}|_{U_j} \cong \mathcal{O}_{U_j}$. En particulier, on a
$s|_{U_j} = f_j q_j$. Un calcul utilisant cette égalité montre que
$t_j|_{U_j \cap U_{j'}}$ et $t_{j'}|_{U_j \cap U_{j'}}$
ont pour image la même section de $\mathcal{F}$ sur $U_j \cap U_{j'} \cap X_s$.
Par quasi-compacité de $U_j \cap U_{j'}$ et d'après (1), il existe un
entier $e' \geq 0$ tel que
$$
t_j|_{U_j \cap U_{j'}} \otimes s^{e'}|_{U_j \cap U_{j'}} =
t_{j'}|_{U_j \cap U_{j'}} \otimes s^{e'}|_{U_j \cap U_{j'}}
$$
comme sections de $\mathcal{F} \otimes \mathcal{L}^{\otimes e + e'}$ sur
$U_j \cap U_{j'}$. On peut choisir un même $e'$ convenant à tous les couples
$j, j'$. La condition de faisceau fournit alors une section
$t \in \Gamma(X, \mathcal{F} \otimes \mathcal{L}^{\otimes e + e'})$
dont la restriction à $U_j$ est $t_j \otimes s^{e'}|_{U_j}$.
Un calcul immédiat montre que $t/s^{e + e'}$ a pour image $t'$,
comme voulu.
\end{proof}

\noindent
Soit $X$ un schéma. Soit $\mathcal{L}$ un $\mathcal{O}_X$-module inversible.
Soient $\mathcal{F}$ et $\mathcal{G}$ des $\mathcal{O}_X$-modules quasi-cohérents.
Considérons le $\Gamma_*(X, \mathcal{L})$-module gradué
$$
M = \bigoplus\nolimits_{n \in \mathbf{Z}} \Hom_{\mathcal{O}_X}(\mathcal{F},
\mathcal{G} \otimes_{\mathcal{O}_X} \mathcal{L}^{\otimes n})
$$
Soit maintenant $s \in \Gamma(X, \mathcal{L})$ une section. On dispose d'une
application canonique
\begin{equation}
\label{equation-hom-invert-s}
M_{(s)} \longrightarrow
\Hom_{\mathcal{O}_{X_s}}(\mathcal{F}|_{X_s}, \mathcal{G}|_{X_s})
\end{equation}
qui envoie $\alpha/s^n$ sur le morphisme $\alpha|_{X_s} \otimes s|_{X_s}^{-n}$.
Le lemme suivant, combiné au
lemme \ref{lemma-extend-finite-presentation},
affirme en substance que, si $X$ est quasi-compact et quasi-séparé,
la catégorie des $\mathcal{O}_{X_s}$-modules de présentation finie
s'obtient à partir de la catégorie des $\mathcal{O}_X$-modules de présentation finie
en inversant le système multiplicatif de morphismes
$s^n: \mathcal{F} \to
\mathcal{F} \otimes_{\mathcal{O}_X} \mathcal{L}^{\otimes n}$.

\begin{lemma}
\label{lemma-section-maps-backwards}
Soit $X$ un schéma. Soit $\mathcal{L}$ un $\mathcal{O}_X$-module inversible.
Soit $s \in \Gamma(X, \mathcal{L})$ une section.
Soient $\mathcal{F}$, $\mathcal{G}$ des $\mathcal{O}_X$-modules quasi-cohérents.
\begin{enumerate}
\item Si $X$ est quasi-compact et si $\mathcal{F}$ est de type fini,
l'application (\ref{equation-hom-invert-s}) est injective ;
\item si $X$ est quasi-compact et quasi-séparé et si $\mathcal{F}$
est de présentation finie, l'application
(\ref{equation-hom-invert-s})
est bijective.
\end{enumerate}
\end{lemma}

\begin{proof}
Démontrons d'abord le lemme lorsque $X = \Spec(A)$ est affine
et $\mathcal{L} = \mathcal{O}_X$. Dans ce cas, $s$ correspond
à un élément $f \in A$. Écrivons
$\mathcal{F} = \widetilde{M}$ et $\mathcal{G} = \widetilde{N}$
pour des $A$-modules $M$ et $N$. Le lemme se traduit alors,
à l'aide des lemmes \ref{lemma-finite-type-module} et
\ref{lemma-finite-presentation-module}, par
les assertions algébriques suivantes :
\begin{enumerate}
\item Si $M$ est un $A$-module de type fini et si $\varphi : M \to N$ est
un morphisme de $A$-modules tel que le morphisme induit $M_f \to N_f$ soit nul,
alors $f^n\varphi = 0$ pour un certain $n$.
\item Si $M$ est un $A$-module de présentation finie, alors
$\Hom_A(M, N)_f = \Hom_{A_f}(M_f, N_f)$.
\end{enumerate}
La seconde assertion est
Algèbre, lemme \ref{algebra-lemma-hom-from-finitely-presented} ; nous omettons
la démonstration de la première.

\medskip\noindent
Démontrons maintenant (1) pour $X$ quelconque.
Supposons $X$ quasi-compact et choisissons un recouvrement ouvert affine fini
$X = U_1 \cup \ldots \cup U_m$, avec $U_j$ affine et
$\mathcal{L}|_{U_j} \cong \mathcal{O}_{U_j}$. Par cet isomorphisme,
la section $s|_{U_j}$ correspond à un élément
$f_j \in \Gamma(U_j, \mathcal{O}_{U_j})$. Alors
$X_s \cap U_j = D(f_j)$.
Soit $\alpha/s^n$ un élément du noyau de
(\ref{equation-hom-invert-s}). Alors $\alpha|_{X_s} = 0$.
Donc $(\alpha|_{U_j})|_{D(f_j)} = 0$. D'après le cas affine traité ci-dessus,
on obtient $f_j^{e_j} \alpha|_{U_j} = 0$ pour un certain
$e_j \geq 0$. Soit $e = \max(e_j)$. Alors $\alpha \otimes s^e$
se restreint à zéro sur $U_j$ pour tout $j$, et est donc nul. Comme $\alpha/s^n$
est égal à $\alpha \otimes s^e/s^{n + e}$ dans $M_{(s)}$, on en déduit
que $\alpha/s^n = 0$, comme voulu.

\medskip\noindent
Démonstration de (2). Comme $\mathcal{F}$ est de présentation finie, le
faisceau $\SheafHom_{\mathcal{O}_X}(\mathcal{F}, \mathcal{G})$ est
quasi-cohérent, voir Schémas, section \ref{schemes-section-quasi-coherent}.
De plus, il est clair que
$$
\SheafHom_{\mathcal{O}_X}(\mathcal{F},
\mathcal{G} \otimes_{\mathcal{O}_X} \mathcal{L}^{\otimes n}) =
\SheafHom_{\mathcal{O}_X}(\mathcal{F}, \mathcal{G})
\otimes_{\mathcal{O}_X} \mathcal{L}^{\otimes n}
$$
pour tout $n$. Dans ce cas, l'assertion résulte donc du
lemme \ref{lemma-invert-s-sections} appliqué à
$\SheafHom_{\mathcal{O}_X}(\mathcal{F}, \mathcal{G})$.
\end{proof}

\section{Schémas quasi-affines}
\label{section-quasi-affine}

\begin{definition}
\label{definition-quasi-affine}
Un schéma $X$ est dit {\it quasi-affine} s'il est quasi-compact
et isomorphe à un sous-schéma ouvert d'un schéma affine.
\end{definition}

\begin{lemma}
\label{lemma-quasi-coherent-quasi-affine}
Soient $A$ un anneau et $U \subset \Spec(A)$ un sous-schéma ouvert
quasi-compact. Pour $\mathcal{F}$ quasi-cohérent sur $U$, le morphisme canonique
$$
\widetilde{H^0(U, \mathcal{F})}|_U \to \mathcal{F}
$$
est un isomorphisme.
\end{lemma}

\begin{proof}
Notons $j : U \to \Spec(A)$ le morphisme d'inclusion. Alors
$H^0(U, \mathcal{F}) = H^0(\Spec(A), j_*\mathcal{F})$, et
$j_*\mathcal{F}$ est quasi-cohérent d'après
Schémas, lemme \ref{schemes-lemma-push-forward-quasi-coherent}.
Ainsi $j_*\mathcal{F} = \widetilde{H^0(U, \mathcal{F})}$ par
Schémas, lemme \ref{schemes-lemma-equivalence-quasi-coherent}.
En restreignant à $U$, on obtient le lemme.
\end{proof}

\begin{lemma}
\label{lemma-invert-f-affine}
Soient $X$ un schéma et $f \in \Gamma(X, \mathcal{O}_X)$.
Supposons $X$ quasi-compact et quasi-séparé, et supposons que
$X_f$ soit affine. Alors le morphisme canonique
$$
j : X \longrightarrow \Spec(\Gamma(X, \mathcal{O}_X))
$$
de Schémas, lemme \ref{schemes-lemma-morphism-into-affine},
induit un isomorphisme de $X_f = j^{-1}(D(f))$ sur l'ouvert affine principal
$D(f) \subset \Spec(\Gamma(X, \mathcal{O}_X))$.
\end{lemma}

\begin{proof}
C'est immédiat, puisque $j$ induit un isomorphisme d'anneaux
$\Gamma(X, \mathcal{O}_X)_f \to \mathcal{O}_X(X_f)$ d'après le
lemme \ref{lemma-invert-f-sections} ci-dessus.
\end{proof}

\begin{lemma}
\label{lemma-quasi-affine}
Soit $X$ un schéma. Alors $X$ est quasi-affine si et seulement si
le morphisme canonique
$$
X \longrightarrow \Spec(\Gamma(X, \mathcal{O}_X))
$$
de Schémas, lemme \ref{schemes-lemma-morphism-into-affine}, est
une immersion ouverte quasi-compacte.
\end{lemma}

\begin{proof}
Si le morphisme affiché est une immersion ouverte quasi-compacte,
$X$ est isomorphe à un sous-schéma ouvert quasi-compact de
$\Spec(\Gamma(X, \mathcal{O}_X))$, et $X$ est évidemment quasi-affine.

\medskip\noindent
Supposons $X$ quasi-affine ; écrivons $X \subset \Spec(R)$ comme un
ouvert quasi-compact. Cela entraîne en particulier que $X$ est
séparé, voir
Schémas, lemme \ref{schemes-lemma-subscheme-of-separated-scheme}.
Soit $A = \Gamma(X, \mathcal{O}_X)$.
Considérons l'homomorphisme d'anneaux $R \to A$ provenant de
$R = \Gamma(\Spec(R), \mathcal{O}_{\Spec(R)})$
et de l'application de restriction du faisceau $\mathcal{O}_{\Spec(R)}$.
D'après Schémas, lemme \ref{schemes-lemma-morphism-into-affine},
on obtient une factorisation
$$
X \longrightarrow
\Spec(A) \longrightarrow
\Spec(R)
$$
du morphisme d'inclusion. Soit $x \in X$. Choisissons $r \in R$ tel que
$x \in D(r)$ et $D(r) \subset X$. Notons $f \in A$ l'image de $r$
dans $A$. L'ouvert $X_f$ du lemme \ref{lemma-invert-f-sections}
ci-dessus est égal à $D(r) \subset X$ ; donc $A_f \cong R_r$ d'après
ce lemme.
Ainsi $D(r) \to \Spec(A)$ est un isomorphisme sur
l'ouvert affine principal $D(f)$ de $\Spec(A)$. Comme $X$
peut être recouvert par de tels ouverts affines $D(f)$, on conclut.
\end{proof}

\begin{lemma}
\label{lemma-cartesian-diagram-quasi-affine}
Soit $U \to V$ une immersion ouverte de schémas quasi-affines. Alors
$$
\xymatrix{
U \ar[d] \ar[rr]_-j & & \Spec(\Gamma(U, \mathcal{O}_U)) \ar[d] \\
U \ar[r] & V \ar[r]^-{j'} & \Spec(\Gamma(V, \mathcal{O}_V))
}
$$
est cartésien.
\end{lemma}

\begin{proof}
Le diagramme est commutatif d'après Schémas, lemme
\ref{schemes-lemma-morphism-into-affine}.
Écrivons $A = \Gamma(U, \mathcal{O}_U)$ et $B = \Gamma(V, \mathcal{O}_V)$. Soit
$g \in B$ tel que $V_g$ soit affine et contenu dans $U$. Cela
signifie que, si $f$ est l'image de $g$ dans $A$, alors $U_f = V_g$. D'après le lemme
\ref{lemma-invert-f-affine}, $j'$ induit un isomorphisme de
$V_g$ sur l'ouvert principal $D(g)$ de $\Spec(B)$. Ainsi
$V_g \times_{\Spec(B)} \Spec(A) \to \Spec(A)$ est un
isomorphisme sur $D(f) \subset \Spec(A)$. De nouveau par le lemme \ref{lemma-invert-f-affine},
$j$ envoie $U_f$ isomorphiquement sur $D(f)$. On obtient donc
$U_f = U_f \times_{\Spec(B)} \Spec(A)$. Puisque le
lemme \ref{lemma-quasi-affine} permet de recouvrir $U$ par des $V_g = U_f$ comme ci-dessus,
on voit que $U \to U \times_{\Spec(B)} \Spec(A)$ est un isomorphisme.
\end{proof}

\begin{lemma}
\label{lemma-quasi-affine-presentation}
Soit $X$ un schéma quasi-affine. Il existe un entier $n \geq 0$,
un schéma affine $T$ et un morphisme $T \to X$ tels que, pour tout
morphisme $X' \to X$ avec $X'$ affine, le produit fibré $X' \times_X T$
soit isomorphe à $\mathbf{A}^n_{X'}$ au-dessus de $X'$.
\end{lemma}

\begin{proof}
Par définition, il existe un anneau $A$ tel que $X$ soit isomorphe à un
sous-schéma ouvert quasi-compact $U \subset \Spec(A)$. Rappelons que les ouverts
principaux $D(f) \subset \Spec(A)$ forment une base de la topologie, voir
Algèbre, section \ref{algebra-section-spectrum-ring}. Comme $U$ est
quasi-compact, on peut choisir $f_1, \ldots, f_n \in A$ tels que
$U = D(f_1) \cup \ldots \cup D(f_n)$. On peut donc supposer
$X = \Spec(A) \setminus V(I)$, où $I = (f_1, \ldots, f_n)$. Posons
$$
T = \Spec(A[t, x_1, \ldots, x_n]/(f_1 x_1 + \ldots + f_n x_n - 1))
$$
Le morphisme structural $T \to \Spec(A)$ se factorise par l'ouvert $X$
et fournit le morphisme $T \to X$. Si $X' = \Spec(A')$ et si le morphisme
$X' \to X$ correspond à l'homomorphisme d'anneaux $A \to A'$, les images
$f'_1, \ldots, f'_n \in A'$ de $f_1, \ldots, f_n$
engendrent l'idéal unité de $A'$.
Écrivons $1 = f'_1 a'_1 + \ldots + f'_n a'_n$.
Le changement de base $X' \times_X T$ est le spectre de
$A'[t, x_1, \ldots, x_n]/(f'_1 x_1 + \ldots + f'_n x_n - 1)$.
Nous affirmons que l'homomorphisme de $A'$-algèbres
$$
\varphi :
A'[y_1, \ldots, y_n]
\longrightarrow
A'[t, x_1, \ldots, x_n, x_{n + 1}]/(f'_1 x_1 + \ldots + f'_n x_n - 1)
$$
envoyant $y_i$ sur $a'_i t + x_i$ est un isomorphisme. Cette assertion achève
la démonstration du lemme. L'inverse de $\varphi$ est donné par l'homomorphisme de $A'$-algèbres

$$
\psi :
A'[t, x_1, \ldots, x_n, x_{n + 1}]/(f'_1 x_1 + \ldots + f'_n x_n - 1)
\longrightarrow
A'[y_1, \ldots, y_n]
$$
envoyant $t$ sur $-1 + f'_1 y_1 + \ldots + f'_n y_n$ et $x_i$ sur
$y_i + a'_i - a'_i(f'_1 y_1 + \ldots + f'_n y_n)$ pour $i = 1, \ldots, n$.
Cela a un sens, puisque $\sum f'_ix_i$ est envoyé sur
$$
\begin{matrix}
\sum f'_i(y_i + a'_i - a'_i(\sum f'_j y_j)) =
(\sum f'_iy_i) + 1 - (\sum f'_j y_j) = 1
\end{matrix}
$$
Pour vérifier que ces homomorphismes sont inverses l'un de l'autre, on calcule :
$$
\begin{matrix}
\varphi(\psi(t) = \varphi(-1 +  \sum f'_i y_i) =
-1 + \sum f'_i (a'_i t + x_i) = t \\
\varphi(\psi(x_i)) = \varphi(y_i + a'_i - a'_i(\sum f'_j y_j)) =
a'_i t + x_i + a'_i - a'_i(\sum f'_ja'_jt + f'_jx_j) = x_i \\
\psi(\varphi(y_i)) = \psi(a'_i t + x_i) =
a'_i(-1 + \sum f'_j y_j) + y_i + a'_i - a'_i(\sum f'_j y_j) = y_i
\end{matrix}
$$
Cela achève la démonstration.
\end{proof}

\section{Modules plats}
\label{section-flat}

\noindent
Sur tout espace annelé $(X, \mathcal{O}_X)$,
on sait ce que signifie, pour un $\mathcal{O}_X$-module,
être plat (en un point), voir
Modules, définition \ref{modules-definition-flat}
(définition \ref{modules-definition-flat-at-point}).
Pour les faisceaux quasi-cohérents sur un schéma affine, cette notion coïncide
avec celle définie dans le chapitre d'algèbre.

\begin{lemma}
\label{lemma-flat-module}
\begin{slogan}
La platitude est la même pour les modules et les faisceaux.
\end{slogan}
Soit $X = \Spec(R)$ un schéma affine.
Soit $\mathcal{F} = \widetilde{M}$ pour un $R$-module $M$.
Le faisceau quasi-cohérent $\mathcal{F}$ est un
$\mathcal{O}_X$-module plat si et seulement si $M$ est un $R$-module plat.
\end{lemma}

\begin{proof}
La platitude de $\mathcal{F}$ se vérifie sur les germes, voir
Modules, lemme \ref{modules-lemma-flat-stalks-flat}.
Il en va de même pour les modules sur un anneau, voir
Algèbre, lemme \ref{algebra-lemma-flat-localization}.
Puisque $\mathcal{F}_x = M_{\mathfrak p}$ lorsque $x$ correspond
à $\mathfrak p$, le lemme en résulte.
\end{proof}

\section{Modules localement libres}
\label{section-finite-locally-free}

\noindent
Sur tout espace annelé, on sait ce que signifie, pour un $\mathcal{O}_X$-module,
être localement libre (de rang fini). Sur un schéma affine, cette notion coïncide
avec celle définie dans le chapitre d'algèbre.

\begin{lemma}
\label{lemma-locally-free-module}
Soit $X = \Spec(R)$ un schéma affine.
Soit $\mathcal{F} = \widetilde{M}$ pour un $R$-module $M$.
Le faisceau quasi-cohérent $\mathcal{F}$ est un
$\mathcal{O}_X$-module localement libre (de rang fini) si et seulement si $M$ est un
$R$-module localement libre (de rang fini).
\end{lemma}

\begin{proof}
Cela résulte des définitions, voir
Modules, définition \ref{modules-definition-locally-free},
et
Algèbre, définition \ref{algebra-definition-locally-free}.
\end{proof}

\noindent
On peut caractériser les modules localement libres de rang fini de plusieurs façons.

\begin{lemma}
\label{lemma-finite-locally-free}
Soit $X$ un schéma.
Soit $\mathcal{F}$ un $\mathcal{O}_X$-module quasi-cohérent.
Les assertions suivantes sont équivalentes :
\begin{enumerate}
\item $\mathcal{F}$ est un $\mathcal{O}_X$-module plat de présentation finie ;
\item $\mathcal{F}$ est un $\mathcal{O}_X$-module de présentation finie et,
pour tout $x \in X$, le germe $\mathcal{F}_x$ est un
$\mathcal{O}_{X, x}$-module libre ;
\item $\mathcal{F}$ est un $\mathcal{O}_X$-module localement libre et de type fini ;
\item $\mathcal{F}$ est un $\mathcal{O}_X$-module localement libre de rang fini ;
\item $\mathcal{F}$ est un $\mathcal{O}_X$-module de type fini,
pour tout $x \in X$, le germe $\mathcal{F}_x$ est un
$\mathcal{O}_{X, x}$-module libre, et la fonction
$$
\rho_\mathcal{F} : X \to \mathbf{Z}, \quad
x \longmapsto
\dim_{\kappa(x)} \mathcal{F}_x \otimes_{\mathcal{O}_{X, x}} \kappa(x)
$$
est localement constante pour la topologie de Zariski de $X$.
\end{enumerate}
\end{lemma}

\begin{proof}
On se ramène immédiatement au cas affine.
Dans ce cas, le lemme est une reformulation de
Algèbre, lemme \ref{algebra-lemma-finite-projective}.
Le passage d'une formulation à l'autre utilise les
lemmes \ref{lemma-finite-type-module},
\ref{lemma-finite-presentation-module},
\ref{lemma-flat-module} et
\ref{lemma-locally-free-module}.
\end{proof}

\begin{lemma}
\label{lemma-finite-locally-free-reduced}
Soit $X$ un schéma réduit. Soit $\mathcal{F}$ un
$\mathcal{O}_X$-module quasi-cohérent. Les conditions équivalentes du
lemme \ref{lemma-finite-locally-free} sont alors aussi équivalentes à la suivante :
\begin{enumerate}
\item[(6)] $\mathcal{F}$ est un $\mathcal{O}_X$-module de type fini et
la fonction
$$
\rho_\mathcal{F} : X \to \mathbf{Z}, \quad
x \longmapsto
\dim_{\kappa(x)} \mathcal{F}_x \otimes_{\mathcal{O}_{X, x}} \kappa(x)
$$
est localement constante pour la topologie de Zariski de $X$.
\end{enumerate}
\end{lemma}

\begin{proof}
On se ramène immédiatement au cas affine.
Dans ce cas, le lemme est une reformulation de
Algèbre, lemme \ref{algebra-lemma-finite-projective-reduced}.
\end{proof}

\section{Modules localement projectifs}
\label{section-locally-projective}

\noindent
Les résultats établis dans le chapitre d'algèbre permettent
de définir les modules localement projectifs comme suit.

\begin{definition}
\label{definition-locally-projective}
Soit $X$ un schéma. Soit $\mathcal{F}$ un
$\mathcal{O}_X$-module quasi-cohérent. On dit que $\mathcal{F}$ est {\it localement projectif}
si, pour tout ouvert affine $U \subset X$, le $\mathcal{O}_X(U)$-module
$\mathcal{F}(U)$ est projectif.
\end{definition}

\begin{lemma}
\label{lemma-locally-projective}
Soit $X$ un schéma.
Soit $\mathcal{F}$ un $\mathcal{O}_X$-module quasi-cohérent.
Les assertions suivantes sont équivalentes :
\begin{enumerate}
\item $\mathcal{F}$ est localement projectif ;
\item il existe un recouvrement ouvert affine $X = \bigcup U_i$
tel que le $\mathcal{O}_X(U_i)$-module
$\mathcal{F}(U_i)$ soit projectif pour tout $i$.
\end{enumerate}
En particulier, si $X = \Spec(A)$ et $\mathcal{F} = \widetilde{M}$,
alors $\mathcal{F}$ est localement projectif si et seulement si $M$ est un
$A$-module projectif.
\end{lemma}

\begin{proof}
Remarquons d'abord que, si $M$ est un $A$-module projectif et si $A \to B$ est un
homomorphisme d'anneaux, alors $M \otimes_A B$ est un $B$-module projectif, voir
Algèbre, lemme \ref{algebra-lemma-ascend-properties-modules}.
Ainsi, si $U$ est un ouvert affine tel que $\mathcal{F}(U)$ soit un
$\mathcal{O}_X(U)$-module projectif, l'ouvert principal $D(f)$ est un
ouvert affine tel que $\mathcal{F}(D(f))$ soit un
$\mathcal{O}_X(D(f))$-module projectif, pour tout $f \in \mathcal{O}_X(U)$.
Supposons (2). Soit $U \subset X$ un ouvert affine quelconque.
On peut trouver un recouvrement ouvert $U = \bigcup_{j = 1, \ldots, m} D(f_j)$
par un nombre fini d'ouverts principaux $D(f_j)$ tels que, pour tout
$j$, l'ouvert $D(f_j)$ soit un ouvert principal d'un certain $U_i$, voir
Schémas, lemme \ref{schemes-lemma-standard-open-two-affines}.
Ainsi, si l'on pose $A = \mathcal{O}_X(U)$ et si $M$ est un $A$-module
tel que $\mathcal{F}|_U$ corresponde à $M$, alors
$M_{f_j}$ est un $A_{f_j}$-module projectif. Il en résulte que
$A \to B = \prod A_{f_j}$ est un homomorphisme d'anneaux fidèlement plat
tel que $M \otimes_A B$ soit un $B$-module projectif.
Donc $M$ est projectif d'après
Algèbre, théorème \ref{algebra-theorem-ffdescent-projectivity}.
\end{proof}

\begin{lemma}
\label{lemma-locally-projective-pullback}
Soit $f : X \to Y$ un morphisme de schémas.
Soit $\mathcal{G}$ un $\mathcal{O}_Y$-module quasi-cohérent.
Si $\mathcal{G}$ est localement projectif sur $Y$, alors $f^*\mathcal{G}$
est localement projectif sur $X$.
\end{lemma}

\begin{proof}
Cela résulte de
Algèbre, lemme \ref{algebra-lemma-ascend-properties-modules},
et du lemme \ref{lemma-locally-projective}.
\end{proof}

\section{Prolongement des faisceaux quasi-cohérents}
\label{section-extending-quasi-coherent-sheaves}

\noindent
Il est parfois utile de montrer qu'un faisceau quasi-cohérent donné
sur un sous-schéma ouvert se prolonge au schéma tout entier.

\begin{lemma}
\label{lemma-extend-trivial}
Soit $j : U \to X$ une immersion ouverte quasi-compacte de schémas.
\begin{enumerate}
\item Tout faisceau quasi-cohérent sur $U$ se prolonge en un faisceau
quasi-cohérent sur $X$.
\item Soit $\mathcal{F}$ un faisceau quasi-cohérent sur $X$.
Soit $\mathcal{G} \subset \mathcal{F}|_U$ un sous-faisceau
quasi-cohérent. Il existe un sous-faisceau quasi-cohérent $\mathcal{H}$ de
$\mathcal{F}$ tel que $\mathcal{H}|_U = \mathcal{G}$
comme sous-faisceaux de $\mathcal{F}|_U$.
\item Soit $\mathcal{F}$ un faisceau quasi-cohérent sur $X$.
Soit $\mathcal{G}$ un faisceau quasi-cohérent sur $U$.
Soit $\varphi : \mathcal{G} \to \mathcal{F}|_U$ un morphisme
de $\mathcal{O}_U$-modules. Il existe un faisceau quasi-cohérent $\mathcal{H}$
de $\mathcal{O}_X$-modules et un morphisme $\psi : \mathcal{H} \to \mathcal{F}$
tels que $\mathcal{H}|_U = \mathcal{G}$ et que
$\psi|_U = \varphi$.
\end{enumerate}
\end{lemma}

\begin{proof}
Une immersion est séparée
(voir Schémas, lemme \ref{schemes-lemma-immersions-monomorphisms}),
et $j$ est quasi-compact par hypothèse.
Ainsi, pour tout faisceau quasi-cohérent $\mathcal{G}$ sur $U$, le faisceau
$j_*\mathcal{G}$ est un prolongement à $X$. Voir
Schémas, lemme \ref{schemes-lemma-push-forward-quasi-coherent}, et
Faisceaux, section \ref{sheaves-section-open-immersions}.

\medskip\noindent
Supposons $\mathcal{F}$ et $\mathcal{G}$ comme dans (2).
Alors $j_*\mathcal{G}$ est un faisceau quasi-cohérent sur $X$ (voir ci-dessus).
C'est un sous-faisceau de $j_*j^*\mathcal{F}$.
Le noyau
$$
\mathcal{H} =
\Ker(\mathcal{F} \oplus j_* \mathcal{G}
\longrightarrow j_*j^*\mathcal{F})
$$
est donc lui aussi quasi-cohérent, voir
Schémas, section \ref{schemes-section-quasi-coherent}.
On vérifie formellement que $\mathcal{H} \subset \mathcal{F}$ et que
$\mathcal{H}|_U = \mathcal{G}$ (en utilisant de nouveau les résultats de
Faisceaux, section \ref{sheaves-section-open-immersions}).

\medskip\noindent
L'assertion (3) se démontre de la même manière que (2). Il suffit de prendre
$\mathcal{H} = \Ker(\mathcal{F} \oplus j_* \mathcal{G}
\to j_*j^*\mathcal{F})$, muni de son morphisme évident vers $\mathcal{F}$
et de son identification évidente avec $\mathcal{G}$ sur $U$.
\end{proof}

\begin{lemma}
\label{lemma-extend}
Soit $X$ un schéma quasi-compact et quasi-séparé.
Soit $U \subset X$ un ouvert quasi-compact.
Soit $\mathcal{F}$ un $\mathcal{O}_X$-module quasi-cohérent.
Soit $\mathcal{G} \subset \mathcal{F}|_U$ un
$\mathcal{O}_U$-sous-module quasi-cohérent de type fini. Alors
il existe un sous-module quasi-cohérent $\mathcal{G}' \subset \mathcal{F}$
de type fini tel que $\mathcal{G}'|_U = \mathcal{G}$.
\end{lemma}

\begin{proof}
Soit $n$ le nombre minimal d'ouverts affines $U_i \subset X$,
$i = 1, \ldots , n$, tels que $X = U \cup \bigcup U_i$.
(On utilise ici la quasi-compacité de $X$.) Supposons que
l'on sache démontrer le lemme dans le cas $n = 1$. On peut alors prolonger successivement
$\mathcal{G}$
en un faisceau $\mathcal{G}_1$ sur $U \cup U_1$,
puis en un faisceau $\mathcal{G}_2$ sur $U \cup U_1 \cup U_2$,
puis en un faisceau $\mathcal{G}_3$ sur $U \cup U_1 \cup U_2 \cup U_3$,
et ainsi de suite.
On se ramène donc au cas $n = 1$.

\medskip\noindent
On peut ainsi supposer que $X = U \cup V$, avec $V$ affine.
Comme $X$ est quasi-séparé et que $U$, $V$ sont des ouverts quasi-compacts,
$U \cap V$ est un ouvert quasi-compact. Il suffit de démontrer le
lemme pour le système $(V, U \cap V, \mathcal{F}|_V, \mathcal{G}|_{U \cap V})$,
car on peut recoller le faisceau $\mathcal{G}'$ ainsi obtenu sur $V$
au faisceau donné $\mathcal{G}$ sur $U$, le long de leur restriction commune
à $U \cap V$.
On se ramène donc au cas où $X$ est affine.

\medskip\noindent
Supposons $X = \Spec(R)$. Écrivons $\mathcal{F} = \widetilde M$
pour un $R$-module $M$. Le lemme \ref{lemma-extend-trivial} ci-dessus permet
de trouver un sous-faisceau quasi-cohérent $\mathcal{H} \subset \mathcal{F}$
qui se restreigne à $\mathcal{G}$ sur $U$.
Écrivons $\mathcal{H} = \widetilde N$ pour un $R$-module $N$.
Pour tout $u \in U$, il existe un $f \in R$ tel que
$u \in D(f) \subset U$ et que $N_f$ soit de type fini,
voir le lemme \ref{lemma-finite-type-module}.
Comme $U$ est quasi-compact, on peut le recouvrir par un nombre fini
d'ouverts $D(f_i)$ tels que $N_{f_i}$ soit engendré par
un nombre fini d'éléments, notés $x_{i, 1}/f_i^N, \ldots, x_{i, r_i}/f_i^N$.
Soit $N' \subset N$ le sous-module engendré par les éléments
$x_{i, j}$. Le sous-faisceau
$\mathcal{G}' = \widetilde{N'} \subset \mathcal{H} \subset \mathcal{F}$
convient.
\end{proof}

\begin{lemma}
\label{lemma-quasi-coherent-colimit-finite-type}
Soit $X$ un schéma quasi-compact et quasi-séparé.
Tout faisceau quasi-cohérent de $\mathcal{O}_X$-modules
est la colimite filtrante de ses
$\mathcal{O}_X$-sous-modules quasi-cohérents de type fini.
\end{lemma}

\begin{proof}
La colimite est filtrante, car, si $\mathcal{G}_1$ et $\mathcal{G}_2$
sont des sous-faisceaux quasi-cohérents de type fini, l'image de
$\mathcal{G}_1 \oplus \mathcal{G}_2 \to \mathcal{F}$ est
un sous-module quasi-cohérent de type fini.
Soit $U \subset X$ un ouvert affine quelconque, et soit
$s \in \Gamma(U, \mathcal{F})$ une section quelconque.
Soit $\mathcal{G} \subset \mathcal{F}|_U$ le
sous-faisceau engendré par $s$. Alors $\mathcal{G}$
est évidemment quasi-cohérent et de type fini comme $\mathcal{O}_U$-module.
Le lemme \ref{lemma-extend} montre que $\mathcal{G}$ est la restriction
d'un sous-faisceau quasi-cohérent $\mathcal{G}' \subset \mathcal{F}$
de type fini. Comme $X$ possède une base de la topologie formée
d'ouverts affines, on en déduit que toute section locale de
$\mathcal{F}$ appartient localement à un sous-module quasi-cohérent
de type fini. Cela conclut la démonstration.
\end{proof}

\begin{lemma}
\label{lemma-extend-finite-presentation}
Soit $X$ un schéma quasi-compact et quasi-séparé.
Soit $\mathcal{F}$ un $\mathcal{O}_X$-module quasi-cohérent.
Soit $U \subset X$ un ouvert quasi-compact.
Soit $\mathcal{G}$ un $\mathcal{O}_U$-module de présentation finie.
Soit $\varphi : \mathcal{G} \to \mathcal{F}|_U$ un morphisme de
$\mathcal{O}_U$-modules.
Il existe alors un $\mathcal{O}_X$-module
$\mathcal{G}'$ de présentation finie et un morphisme
de $\mathcal{O}_X$-modules $\varphi' : \mathcal{G}' \to \mathcal{F}$
tels que $\mathcal{G}'|_U = \mathcal{G}$ et que
$\varphi'|_U = \varphi$.
\end{lemma}

\begin{proof}
Le début de la démonstration reprend celui de la démonstration du
lemme \ref{lemma-extend}. Nous le détaillons néanmoins.

\medskip\noindent
Soit $n$ le nombre minimal d'ouverts affines $U_i \subset X$,
$i = 1, \ldots , n$, tels que $X = U \cup \bigcup U_i$.
(On utilise ici la quasi-compacité de $X$.) Supposons que
l'on sache démontrer le lemme dans le cas $n = 1$. On peut alors prolonger successivement
le couple $(\mathcal{G}, \varphi)$
en un couple $(\mathcal{G}_1, \varphi_1)$ sur $U \cup U_1$,
puis en un couple $(\mathcal{G}_2, \varphi_2)$ sur $U \cup U_1 \cup U_2$,
puis en un couple $(\mathcal{G}_3, \varphi_3)$ sur $U \cup U_1 \cup U_2 \cup U_3$,
et ainsi de suite.
On se ramène donc au cas $n = 1$.

\medskip\noindent
On peut ainsi supposer que $X = U \cup V$, avec $V$ affine.
Comme $X$ est quasi-séparé et que $U$ est quasi-compact,
$U \cap V \subset V$ est quasi-compact.
Supposons que l'on démontre le lemme pour le système
$(V, U \cap V, \mathcal{F}|_V, \mathcal{G}|_{U \cap V}, \varphi|_{U \cap V})$,
en obtenant un couple $(\mathcal{G}', \varphi')$ sur $V$.
On peut alors recoller $\mathcal{G}'$ sur $V$ au faisceau donné $\mathcal{G}$
sur $U$, le long de leur restriction commune à $U \cap V$ ; de même, on peut recoller
le morphisme $\varphi'$ au morphisme $\varphi$ le long de leur restriction commune à
$U \cap V$. On se ramène donc au cas où $X$ est affine.

\medskip\noindent
Supposons $X = \Spec(R)$.
Le lemme \ref{lemma-extend-trivial} ci-dessus permet
de trouver un faisceau quasi-cohérent $\mathcal{H}$ muni
d'un morphisme $\psi : \mathcal{H} \to \mathcal{F}$ sur $X$,
qui se restreignent à $\mathcal{G}$ et à $\varphi$ sur $U$.
Le lemme \ref{lemma-extend} permet de trouver un
$\mathcal{O}_X$-sous-module quasi-cohérent de type fini
$\mathcal{H}' \subset \mathcal{H}$
tel que $\mathcal{H}'|_U = \mathcal{G}$. Ainsi, en
remplaçant $\mathcal{H}$ par $\mathcal{H}'$
et $\psi$ par la restriction de $\psi$ à $\mathcal{H}'$,
on peut supposer $\mathcal{H}$ de type fini.
D'après le lemme \ref{lemma-finite-presentation-module},
on a $\mathcal{H} = \widetilde{N}$, avec
$N$ un $R$-module de type fini. Il existe donc une surjection
figurant dans la suite exacte courte suivante de
$\mathcal{O}_X$-modules quasi-cohérents :
$$
0 \to \mathcal{K} \to \mathcal{O}_X^{\oplus n} \to \mathcal{H} \to 0
$$
où $\mathcal{K}$ est défini comme le noyau.
Comme $\mathcal{G}$ est de présentation finie et que
$\mathcal{H}|_U = \mathcal{G}$, d'après
Modules, lemme
\ref{modules-lemma-kernel-surjection-finite-free-onto-finite-presentation},
la restriction $\mathcal{K}|_U$ est
un $\mathcal{O}_U$-module de type fini. Ainsi, de nouveau par le lemme \ref{lemma-extend},
il existe un
$\mathcal{O}_X$-sous-module quasi-cohérent de type fini $\mathcal{K}' \subset \mathcal{K}$ tel
que $\mathcal{K}'|_U = \mathcal{K}|_U$. Pour résoudre le problème
posé dans le lemme, on prend
$$
\mathcal{G}' = \mathcal{O}_X^{\oplus n}/\mathcal{K}'
$$
qui est évidemment de présentation finie et dont la restriction est $\mathcal{G}$
sur $U$, avec $\varphi'$ égal au composé
$$
\mathcal{G}' = \mathcal{O}_X^{\oplus n}/\mathcal{K}'
\to \mathcal{O}_X^{\oplus n}/\mathcal{K} = \mathcal{H} \xrightarrow{\psi}
\mathcal{F}.
$$
Cela achève la démonstration du lemme.
\end{proof}

\begin{lemma}
\label{lemma-lift-finite-presentation}
Soit $X$ un schéma quasi-compact et quasi-séparé. Soit $U \subset X$
un ouvert quasi-compact. Soit $\mathcal{G}$ un $\mathcal{O}_U$-module.
\begin{enumerate}
\item Si $\mathcal{G}$ est quasi-cohérent et de type fini,
il existe un $\mathcal{O}_X$-module quasi-cohérent $\mathcal{G}'$
de type fini tel que $\mathcal{G}'|_U = \mathcal{G}$.
\item Si $\mathcal{G}$ est de présentation finie,
il existe un $\mathcal{O}_X$-module $\mathcal{G}'$
de présentation finie tel que $\mathcal{G}'|_U = \mathcal{G}$.
\end{enumerate}
\end{lemma}

\begin{proof}
L'assertion (2) est le cas particulier du lemme \ref{lemma-extend-finite-presentation}
où $\mathcal{F} = 0$. Pour (1), on écrit d'abord
$\mathcal{G} = \mathcal{F}|_U$ pour un $\mathcal{O}_X$-module quasi-cohérent,
d'après le lemme \ref{lemma-extend-trivial},
puis on applique le lemme \ref{lemma-extend} avec $\mathcal{G} = \mathcal{F}|_U$.
\end{proof}
\noindent
Le lemme suivant affirme que tout faisceau quasi-cohérent sur un schéma quasi-compact
et quasi-séparé est une colimite filtrante de $\mathcal{O}$-modules
de présentation finie. Dans le lemme qui suit, nous reformulerons ce résultat en termes
de colimites indexées par des ensembles ordonnés filtrants, formulation peut-être plus familière.

\begin{lemma}
\label{lemma-directed-colimit-diagram-finite-presentation}
\begin{slogan}
Les modules quasi-cohérents sur les schémas quasi-compacts et quasi-séparés
sont des colimites filtrantes de modules de présentation finie.
\end{slogan}
Soit $X$ un schéma. Supposons $X$ quasi-compact et quasi-séparé.
Soit $\mathcal{F}$ un $\mathcal{O}_X$-module quasi-cohérent.
Il existe
\begin{enumerate}
\item une catégorie d'indices filtrante $\mathcal{I}$ (voir
Catégories, définition \ref{categories-definition-directed}) ;
\item un diagramme $\mathcal{I} \to \textit{Mod}(\mathcal{O}_X)$ (voir
Catégories, section \ref{categories-section-limits}),
$i \mapsto \mathcal{F}_i$ ;
\item des morphismes de $\mathcal{O}_X$-modules
$\varphi_i : \mathcal{F}_i \to \mathcal{F}$
\end{enumerate}
tels que chaque $\mathcal{F}_i$ soit de présentation finie
et que les morphismes $\varphi_i$ induisent un isomorphisme
$$
\colim_i \mathcal{F}_i
=
\mathcal{F}.
$$
\end{lemma}

\begin{proof}
Choisissons un ensemble $I$ et, pour tout $i \in I$, un $\mathcal{O}_X$-module
de présentation finie muni d'un homomorphisme de $\mathcal{O}_X$-modules
$\varphi_i : \mathcal{F}_i \to \mathcal{F}$, de façon à satisfaire la
propriété suivante : pour tout $\psi : \mathcal{G} \to \mathcal{F}$ avec $\mathcal{G}$
de présentation finie, il existe un $i \in I$ et
un isomorphisme $\alpha : \mathcal{F}_i \to \mathcal{G}$ tels que
$\varphi_i = \psi \circ \alpha$. Il résulte clairement de
Modules, lemme \ref{modules-lemma-set-isomorphism-classes-finite-type-modules},
qu'un tel ensemble existe (voir aussi sa démonstration).
Notons $\mathcal{I}$ la catégorie
dont $\Ob(\mathcal{I}) = I$ et dont, pour $i, i' \in I$,
les ensembles de morphismes sont définis par
$$
\Mor_\mathcal{I}(i, i') =
\{\alpha : \mathcal{F}_i \to \mathcal{F}_{i'} \mid
\alpha \circ \varphi_{i'} = \varphi_i
\}.
$$
Nous affirmons que $\mathcal{I}$ est une catégorie filtrante et que
$\mathcal{F} = \colim_i \mathcal{F}_i$.

\medskip\noindent
Soient $i, i' \in I$. Considérons le morphisme
$$
\mathcal{F}_i \oplus \mathcal{F}_{i'} \longrightarrow \mathcal{F}
$$
qui est la somme directe de $\varphi_i$ et de $\varphi_{i'}$.
Comme une somme directe de $\mathcal{O}_X$-modules de présentation finie
est de présentation finie, il existe un $i'' \in I$
tel que $\varphi_{i''} : \mathcal{F}_{i''} \to \mathcal{F}$
soit isomorphe à la flèche affichée ci-dessus vers $\mathcal{F}$.
Puisque l'on dispose des diagrammes commutatifs
$$
\xymatrix{
\mathcal{F}_i \ar[r] \ar[d] & \mathcal{F} \ar@{=}[d] \\
\mathcal{F}_i \oplus \mathcal{F}_{i'} \ar[r] & \mathcal{F}
}
\quad
\text{et}
\quad
\xymatrix{
\mathcal{F}_{i'} \ar[r] \ar[d] & \mathcal{F} \ar@{=}[d] \\
\mathcal{F}_i \oplus \mathcal{F}_{i'} \ar[r] & \mathcal{F}
}
$$
il existe des morphismes $i \to i''$ et $i' \to i''$
dans $\mathcal{I}$. Supposons maintenant donnés $i, i' \in I$ et
des morphismes $\alpha, \beta : i \to i'$ (correspondant à des morphismes de $\mathcal{O}_X$-modules
$\alpha, \beta : \mathcal{F}_i \to \mathcal{F}_{i'}$).
Considérons leur coégalisateur
$$
\mathcal{G} =
\Coker(
\mathcal{F}_i \xrightarrow{\alpha - \beta} \mathcal{F}_{i'}
)
$$
Remarquons que $\mathcal{G}$ est un $\mathcal{O}_X$-module de présentation finie.
Puisque, par définition des morphismes de la catégorie $\mathcal{I}$,
on a $\varphi_{i'} \circ \alpha = \varphi_{i'} \circ \beta$,
on obtient un morphisme induit $\psi : \mathcal{G} \to \mathcal{F}$.
Le couple $(\mathcal{G}, \psi)$ est donc encore isomorphe
à un couple $(\mathcal{F}_{i''}, \varphi_{i''})$ pour un certain $i''$.
Il existe donc un morphisme $i' \to i''$ dans
$\mathcal{I}$ qui égalise $\alpha$ et $\beta$. Nous avons ainsi
montré que la catégorie $\mathcal{I}$ est filtrante.

\medskip\noindent
Il reste à montrer que la colimite du diagramme est $\mathcal{F}$.
Par définition de la colimite et par notre définition de la catégorie
$\mathcal{I}$, on dispose d'un morphisme canonique
$$
\varphi :
\colim_i \mathcal{F}_i
\longrightarrow
\mathcal{F}.
$$
Soit $x \in X$. Montrons que $\varphi_x$ est un isomorphisme.
Rappelons que
$$
(\colim_i \mathcal{F}_i)_x
=
\colim_i \mathcal{F}_{i, x},
$$
voir
Faisceaux, section \ref{sheaves-section-limits-sheaves}.
Montrons d'abord que $\varphi_x$ est injectif.
Soit $s \in \mathcal{F}_{i, x}$ un élément
tel que $s$ ait une image nulle dans $\mathcal{F}_x$. Il existe alors
un ouvert quasi-compact $U$ tel que $s$ provienne d'une section $s \in \mathcal{F}_i(U)$
et que $\varphi_i(s) = 0$ dans $\mathcal{F}(U)$.
D'après le lemme \ref{lemma-extend},
on peut trouver un sous-faisceau quasi-cohérent de type fini
$\mathcal{K} \subset \Ker(\varphi_i)$ dont la restriction soit
le $\mathcal{O}_U$-sous-module quasi-cohérent de $\mathcal{F}_i$
engendré par $s$ :
$\mathcal{K}|_U = \mathcal{O}_U\cdot s \subset \mathcal{F}_i|_U$.
Il est clair que $\mathcal{F}_i/\mathcal{K}$ est de présentation finie et que
le morphisme $\varphi_i$ se factorise par le morphisme quotient
$\mathcal{F}_i \to \mathcal{F}_i/\mathcal{K}$. On peut donc trouver
un $i' \in I$ et un morphisme $\alpha : \mathcal{F}_i \to \mathcal{F}_{i'}$
dans $\mathcal{I}$ qui s'identifie au morphisme quotient
$\mathcal{F}_i \to \mathcal{F}_i/\mathcal{K}$. Il en résulte
que la section $s$ a une image nulle dans $\mathcal{F}_{i'}(U)$ et,
en particulier, dans
$(\colim_i \mathcal{F}_i)_x =
\colim_i \mathcal{F}_{i, x}$.
L'injectivité en découle.
Montrons enfin que $\varphi_x$ est surjectif.
Soit $s \in \mathcal{F}_x$. Choisissons un voisinage ouvert quasi-compact
$U \subset X$ de $x$ tel que $s$ corresponde à une section
$s \in \mathcal{F}(U)$. Considérons le morphisme
$s : \mathcal{O}_U \to \mathcal{F}$ (multiplication par $s$).
D'après le lemme \ref{lemma-extend-finite-presentation},
il existe un $\mathcal{O}_X$-module $\mathcal{G}$
de présentation finie et un morphisme de $\mathcal{O}_X$-modules
$\mathcal{G} \to \mathcal{F}$ tels que $\mathcal{G}|_U \to \mathcal{F}|_U$
s'identifie à
$s : \mathcal{O}_U \to \mathcal{F}$.
De nouveau, par définition de $\mathcal{I}$, il existe un $i \in I$
tel que $\mathcal{G} \to \mathcal{F}$ soit isomorphe à
$\varphi_i : \mathcal{F}_i \to \mathcal{F}$. Il existe évidemment
une section $s' \in \mathcal{F}_i(U)$ d'image $s \in \mathcal{F}(U)$.
Cela prouve la surjectivité et achève la démonstration du lemme.
\end{proof}

\begin{lemma}
\label{lemma-directed-colimit-finite-presentation}
Soit $X$ un schéma. Supposons $X$ quasi-compact et quasi-séparé.
Soit $\mathcal{F}$ un $\mathcal{O}_X$-module quasi-cohérent.
Il existe
\begin{enumerate}
\item un ensemble ordonné filtrant $I$ (voir
Catégories, définition \ref{categories-definition-directed-set}) ;
\item un système $(\mathcal{F}_i, \varphi_{ii'})$
sur $I$ dans $\textit{Mod}(\mathcal{O}_X)$ (voir
Catégories, définition \ref{categories-definition-system-over-poset}) ;
\item des morphismes de $\mathcal{O}_X$-modules
$\varphi_i : \mathcal{F}_i \to \mathcal{F}$
\end{enumerate}
tels que chaque $\mathcal{F}_i$ soit de présentation finie
et que les morphismes $\varphi_i$ induisent un isomorphisme
$$
\colim_i \mathcal{F}_i
=
\mathcal{F}.
$$
\end{lemma}

\begin{proof}
C'est une conséquence directe du
lemme \ref{lemma-directed-colimit-diagram-finite-presentation} et de
Catégories, lemme \ref{categories-lemma-directed-category-system},
combinés au fait que
les colimites existent dans la catégorie des faisceaux de $\mathcal{O}_X$-modules, voir
Faisceaux, section \ref{sheaves-section-limits-sheaves}.
\end{proof}

\begin{lemma}
\label{lemma-finite-directed-colimit-surjective-maps}
Soit $X$ un schéma. Supposons $X$ quasi-compact et quasi-séparé.
Soit $\mathcal{F}$ un $\mathcal{O}_X$-module quasi-cohérent de type fini.
On peut alors écrire $\mathcal{F} = \colim \mathcal{F}_i$, avec $\mathcal{F}_i$
de présentation finie et tous les morphismes de transition
$\mathcal{F}_i \to \mathcal{F}_{i'}$ surjectifs.
\end{lemma}

\begin{proof}
Écrivons $\mathcal{F} = \colim \mathcal{G}_i$ comme une colimite filtrante de
$\mathcal{O}_X$-modules de présentation finie
(lemme \ref{lemma-directed-colimit-finite-presentation}).
Nous affirmons que $\mathcal{G}_i \to \mathcal{F}$ est surjectif pour un certain $i$.
En effet, choisissons un recouvrement ouvert affine fini $X = U_1 \cup \ldots \cup U_m$.
Choisissons des sections $s_{jl} \in \mathcal{F}(U_j)$ engendrant
$\mathcal{F}|_{U_j}$, voir le lemme \ref{lemma-finite-type-module}.
D'après Faisceaux, lemme \ref{sheaves-lemma-directed-colimits-sections},
$s_{jl}$ appartient à l'image de $\mathcal{G}_i \to \mathcal{F}$
pour $i$ assez grand. Ainsi $\mathcal{G}_i \to \mathcal{F}$ est surjectif
pour $i$ assez grand. Choisissons un tel $i$ et soit
$\mathcal{K} \subset \mathcal{G}_i$ le noyau du morphisme
$\mathcal{G}_i \to \mathcal{F}$. Écrivons $\mathcal{K} = \colim \mathcal{K}_a$
comme la colimite filtrante de ses sous-modules quasi-cohérents de type fini
(lemme \ref{lemma-quasi-coherent-colimit-finite-type}). Alors
$\mathcal{F} = \colim \mathcal{G}_i/\mathcal{K}_a$ résout
le problème posé dans le lemme.
\end{proof}

\begin{lemma}
\label{lemma-application-directed-colimit}
Soit $X$ un schéma quasi-compact et quasi-séparé.
Soit $\mathcal{F}$ un $\mathcal{O}_X$-module quasi-cohérent de type fini.
Soit $U \subset X$ un ouvert quasi-compact tel que $\mathcal{F}|_U$
soit de présentation finie. Il existe alors un morphisme de $\mathcal{O}_X$-modules
$\varphi : \mathcal{G} \to \mathcal{F}$ tel que
(a) $\mathcal{G}$ soit de présentation finie,
(b) $\varphi$ soit surjectif et
(c) $\varphi|_U$ soit un isomorphisme.
\end{lemma}

\begin{proof}
Écrivons $\mathcal{F} = \colim \mathcal{F}_i$ comme une colimite filtrante,
avec chaque $\mathcal{F}_i$ de présentation finie,
voir le lemme \ref{lemma-directed-colimit-finite-presentation}.
Choisissons un recouvrement ouvert affine fini $X = \bigcup V_j$ et choisissons
un nombre fini de sections $s_{jl} \in \mathcal{F}(V_j)$ engendrant
$\mathcal{F}|_{V_j}$, voir le lemme \ref{lemma-finite-type-module}.
D'après Faisceaux, lemme \ref{sheaves-lemma-directed-colimits-sections},
$s_{jl}$ appartient à l'image de $\mathcal{F}_i \to \mathcal{F}$
pour $i$ assez grand. Ainsi $\mathcal{F}_i \to \mathcal{F}$ est surjectif
pour $i$ assez grand. Choisissons un tel $i$ et soit
$\mathcal{K} \subset \mathcal{F}_i$ le noyau du morphisme
$\mathcal{F}_i \to \mathcal{F}$. Comme $\mathcal{F}_U$ est de présentation
finie, $\mathcal{K}|_U$ est de type fini, voir
Modules, lemme
\ref{modules-lemma-kernel-surjection-finite-free-onto-finite-presentation}.
On peut donc trouver un sous-module quasi-cohérent de type fini
$\mathcal{K}' \subset \mathcal{K}$ avec $\mathcal{K}'|_U = \mathcal{K}|_U$,
voir le lemme \ref{lemma-extend}. Alors
$\mathcal{G} = \mathcal{F}_i/\mathcal{K}'$,
muni du morphisme donné $\mathcal{G} \to \mathcal{F}$, convient.
\end{proof}

\noindent
Soit $X$ un schéma. Dans le lemme suivant, nous utilisons la notion
de {\it $\mathcal{O}_X$-algèbre quasi-cohérente $\mathcal{A}$
de présentation finie}. Cela signifie que, pour tout ouvert affine
$\Spec(R) \subset X$, on a $\mathcal{A} = \widetilde{A}$,
où $A$ est une $R$-algèbre (commutative) de présentation finie
comme $R$-algèbre.

\begin{lemma}
\label{lemma-algebra-directed-colimit-finite-presentation}
Soit $X$ un schéma. Supposons $X$ quasi-compact et quasi-séparé.
Soit $\mathcal{A}$ une $\mathcal{O}_X$-algèbre quasi-cohérente.
Il existe
\begin{enumerate}
\item un ensemble ordonné filtrant $I$ (voir
Catégories, définition \ref{categories-definition-directed-set}) ;
\item un système $(\mathcal{A}_i, \varphi_{ii'})$
sur $I$ dans la catégorie des $\mathcal{O}_X$-algèbres ;
\item des morphismes de $\mathcal{O}_X$-algèbres
$\varphi_i : \mathcal{A}_i \to \mathcal{A}$
\end{enumerate}
tels que chaque $\mathcal{A}_i$ soit une $\mathcal{O}_X$-algèbre quasi-cohérente
de présentation finie et que les morphismes $\varphi_i$
induisent un isomorphisme
$$
\colim_i \mathcal{A}_i
=
\mathcal{A}.
$$
\end{lemma}

\begin{proof}
Écrivons d'abord $\mathcal{A} = \colim_i \mathcal{F}_i$ comme une colimite
filtrante de faisceaux quasi-cohérents de présentation finie, comme dans le
lemme \ref{lemma-directed-colimit-finite-presentation}.
Pour tout $i$, soit $\mathcal{B}_i = \text{Sym}(\mathcal{F}_i)$ l'algèbre
symétrique de $\mathcal{F}_i$ sur $\mathcal{O}_X$. Écrivons
$\mathcal{I}_i = \Ker(\mathcal{B}_i \to \mathcal{A})$. Écrivons
$\mathcal{I}_i = \colim_j \mathcal{F}_{i, j}$, où
$\mathcal{F}_{i, j}$ est un sous-module quasi-cohérent de type fini de
$\mathcal{I}_i$, voir le
lemme \ref{lemma-quasi-coherent-colimit-finite-type}.
Définissons $\mathcal{I}_{i, j} \subset \mathcal{I}_i$
comme le $\mathcal{B}_i$-idéal engendré par $\mathcal{F}_{i, j}$.
Posons $\mathcal{A}_{i, j} = \mathcal{B}_i/\mathcal{I}_{i, j}$.
Alors $\mathcal{A}_{i, j}$ est une
$\mathcal{O}_X$-algèbre quasi-cohérente de présentation finie. Définissons $(i, j) \leq (i', j')$ par
les conditions $i \leq i'$ et le fait que le morphisme $\mathcal{B}_i \to \mathcal{B}_{i'}$
envoie l'idéal $\mathcal{I}_{i, j}$ dans l'idéal $\mathcal{I}_{i', j'}$.
Il est alors clair que $\mathcal{A} = \colim_{i, j} \mathcal{A}_{i, j}$.
\end{proof}

\noindent
Soit $X$ un schéma. Dans le lemme suivant, nous utilisons la notion
de {\it $\mathcal{O}_X$-algèbre quasi-cohérente $\mathcal{A}$
de type fini}. Cela signifie que, pour tout ouvert affine
$\Spec(R) \subset X$, on a $\mathcal{A} = \widetilde{A}$,
où $A$ est une $R$-algèbre (commutative) de type fini
comme $R$-algèbre.

\begin{lemma}
\label{lemma-algebra-directed-colimit-finite-type}
Soit $X$ un schéma. Supposons $X$ quasi-compact et quasi-séparé.
Soit $\mathcal{A}$ une $\mathcal{O}_X$-algèbre quasi-cohérente.
Alors $\mathcal{A}$ est la colimite filtrante de ses
$\mathcal{O}_X$-sous-algèbres quasi-cohérentes de type fini.
\end{lemma}

\begin{proof}
Si $\mathcal{A}_1, \mathcal{A}_2 \subset \mathcal{A}$ sont des
$\mathcal{O}_X$-sous-algèbres quasi-cohérentes de type fini, l'image de
$\mathcal{A}_1 \otimes_{\mathcal{O}_X} \mathcal{A}_2 \to \mathcal{A}$
est encore une $\mathcal{O}_X$-sous-algèbre quasi-cohérente de type fini
(certains détails sont omis), qui contient à la fois $\mathcal{A}_1$ et $\mathcal{A}_2$.
Cela montre que le système est filtrant.
Pour montrer que $\mathcal{A}$ en est la colimite, écrivons
$\mathcal{A} = \colim_i \mathcal{A}_i$ comme une colimite
filtrante de $\mathcal{O}_X$-algèbres quasi-cohérentes de présentation finie, comme dans le
lemme \ref{lemma-algebra-directed-colimit-finite-presentation}.
Les images $\mathcal{A}'_i = \Im(\mathcal{A}_i \to \mathcal{A})$ sont alors
des sous-algèbres quasi-cohérentes de type fini de $\mathcal{A}$. Puisque
$\mathcal{A}$ en est la colimite, le résultat en découle.
\end{proof}

\noindent
Soit $X$ un schéma. Dans le lemme suivant, nous utilisons la notion
d'{\it algèbre finie (resp.\ entière) quasi-cohérente
sur $\mathcal{O}_X$, notée $\mathcal{A}$}. Cela signifie que, pour tout
ouvert affine $\Spec(R) \subset X$, on a
$\mathcal{A} = \widetilde{A}$, où $A$ est une $R$-algèbre (commutative)
finie (resp.\ entière) comme $R$-algèbre.

\begin{lemma}
\label{lemma-finite-algebra-directed-colimit-finite-finitely-presented}
Soit $X$ un schéma. Supposons $X$ quasi-compact et quasi-séparé.
Soit $\mathcal{A}$ une $\mathcal{O}_X$-algèbre quasi-cohérente finie.
Alors $\mathcal{A} = \colim \mathcal{A}_i$ est une colimite filtrante de
$\mathcal{O}_X$-algèbres quasi-cohérentes finies et de présentation finie,
telle que tous les morphismes de transition $\mathcal{A}_{i'} \to \mathcal{A}_i$
soient surjectifs.
\end{lemma}

\begin{proof}
D'après le lemme \ref{lemma-finite-directed-colimit-surjective-maps},
il existe un $\mathcal{O}_X$-module de présentation finie
$\mathcal{F}$ et une surjection $\mathcal{F} \to \mathcal{A}$.
La structure d'algèbre fournit une surjection
$$
\text{Sym}^*_{\mathcal{O}_X}(\mathcal{F}) \longrightarrow \mathcal{A}
$$
Notons $\mathcal{J}$ son noyau. Écrivons $\mathcal{J} = \colim \mathcal{E}_i$
comme une colimite filtrante de $\mathcal{O}_X$-sous-modules de type fini
$\mathcal{E}_i$ (lemme \ref{lemma-quasi-coherent-colimit-finite-type}). Posons
$$
\mathcal{A}_i = \text{Sym}^*_{\mathcal{O}_X}(\mathcal{F})/(\mathcal{E}_i)
$$
où $(\mathcal{E}_i)$ désigne le faisceau d'idéaux engendré par
l'image de $\mathcal{E}_i \to \text{Sym}^*_{\mathcal{O}_X}(\mathcal{F})$.
Chaque $\mathcal{A}_i$ est alors une $\mathcal{O}_X$-algèbre de présentation finie,
les morphismes de transition sont surjectifs
et $\mathcal{A} = \colim \mathcal{A}_i$. Pour achever la démonstration, il
reste à montrer que $\mathcal{A}_i$ est une $\mathcal{O}_X$-algèbre finie
pour $i$ assez grand. Pour cela, choisissons un recouvrement ouvert
affine $X = U_1 \cup \ldots \cup U_m$. Prenons des générateurs
$f_{j, 1}, \ldots, f_{j, N_j} \in \Gamma(U_i, \mathcal{F})$.
Comme $\mathcal{A}(U_j)$ est une $\mathcal{O}_X(U_j)$-algèbre finie,
pour tout $k$, il existe un polynôme unitaire
$P_{j, k} \in \mathcal{O}(U_j)[T]$ tel que $P_{j, k}(f_{j, k})$
soit nul dans $\mathcal{A}(U_j)$. Puisque
$\mathcal{A} = \colim \mathcal{A}_i$ par construction,
on a $P_{j, k}(f_{j, k}) = 0$ dans $\mathcal{A}_i(U_j)$
pour tout $i$ assez grand. Pour de tels $i$, les algèbres
$\mathcal{A}_i$ sont finies.
\end{proof}

\begin{lemma}
\label{lemma-integral-algebra-directed-colimit-finite}
Soit $X$ un schéma. Supposons $X$ quasi-compact et quasi-séparé.
Soit $\mathcal{A}$ une $\mathcal{O}_X$-algèbre quasi-cohérente entière.
Alors
\begin{enumerate}
\item $\mathcal{A}$ est la colimite filtrante de ses
$\mathcal{O}_X$-sous-algèbres quasi-cohérentes finies ;
\item $\mathcal{A}$ est une colimite filtrante de
$\mathcal{O}_X$-algèbres quasi-cohérentes finies et de présentation finie.
\end{enumerate}
\end{lemma}

\begin{proof}
D'après le lemme \ref{lemma-algebra-directed-colimit-finite-type}, on a
$\mathcal{A} = \colim \mathcal{A}_i$, où
$\mathcal{A}_i \subset \mathcal{A}$ parcourt les
$\mathcal{O}_X$-algèbres quasi-cohérentes de type fini.
Toute $\mathcal{O}_X$-sous-algèbre quasi-cohérente de type fini
de $\mathcal{A}$ est finie (appliquer Algèbre, lemme
\ref{algebra-lemma-characterize-finite-in-terms-of-integral},
à $\mathcal{A}_i(U) \subset \mathcal{A}(U)$ pour les ouverts affines $U$
de $X$). Cela démontre (1).

\medskip\noindent
Pour démontrer (2), écrivons $\mathcal{A} = \colim \mathcal{F}_i$
comme une colimite de $\mathcal{O}_X$-modules de présentation finie à l'aide du
lemme \ref{lemma-directed-colimit-finite-presentation}.
Pour tout $i$, soit $\mathcal{J}_i$ le noyau du morphisme
$$
\text{Sym}^*_{\mathcal{O}_X}(\mathcal{F}_i) \longrightarrow \mathcal{A}
$$
Pour $i' \geq i$, on a un morphisme induit $\mathcal{J}_i \to \mathcal{J}_{i'}$,
et l'on a $\mathcal{A} =
\colim \text{Sym}^*_{\mathcal{O}_X}(\mathcal{F}_i)/\mathcal{J}_i$.
De plus, les $\mathcal{O}_X$-algèbres quasi-cohérentes
$\text{Sym}^*_{\mathcal{O}_X}(\mathcal{F}_i)/\mathcal{J}_i$
sont finies (voir ci-dessus). Écrivons $\mathcal{J}_i = \colim \mathcal{E}_{ik}$
comme une colimite de $\mathcal{O}_X$-modules de présentation finie.
Étant donnés $i' \geq i$ et $k$, il existe un $k'$ tel que l'on
dispose d'un morphisme $\mathcal{E}_{ik} \to \mathcal{E}_{i'k'}$
rendant
$$
\xymatrix{
\mathcal{J}_i \ar[r] & \mathcal{J}_{i'} \\
\mathcal{E}_{ik} \ar[u] \ar[r] & \mathcal{E}_{i'k'} \ar[u]
}
$$
commutatif. Cela résulte de
Modules, lemme \ref{modules-lemma-finite-presentation-quasi-compact-colimit}.
On en déduit un morphisme
$$
\mathcal{A}_{ik} =
\text{Sym}^*_{\mathcal{O}_X}(\mathcal{F}_i)/(\mathcal{E}_{ik})
\longrightarrow
\text{Sym}^*_{\mathcal{O}_X}(\mathcal{F}_{i'})/(\mathcal{E}_{i'k'}) =
\mathcal{A}_{i'k'}
$$
où $(\mathcal{E}_{ik})$ désigne l'idéal engendré par $\mathcal{E}_{ik}$.
Les $\mathcal{O}_X$-algèbres quasi-cohérentes $\mathcal{A}_{ki}$
sont de présentation finie et sont finies pour $k$ assez grand
(voir la démonstration du
lemme \ref{lemma-finite-algebra-directed-colimit-finite-finitely-presented}).
Enfin, on a
$$
\colim \mathcal{A}_{ik} = \colim \mathcal{A}_i = \mathcal{A}
$$
En effet, la première égalité a été montrée dans la démonstration du
lemme \ref{lemma-finite-algebra-directed-colimit-finite-finitely-presented},
et la seconde résulte de ce que $\mathcal{A}$ est la colimite des
modules $\mathcal{F}_i$.
\end{proof}
\section{Le résultat de Gabber}
\label{section-gabber}

\noindent
Dans cette section, nous démontrons un résultat de Gabber assurant que, sur tout
schéma, il existe un cardinal $\kappa$ tel que tout module quasi-cohérent
$\mathcal{F}$ soit la réunion de ses sous-faisceaux quasi-cohérents
$\kappa$-engendrés. Il en résulte que la catégorie des faisceaux quasi-cohérents
sur un schéma est une catégorie abélienne de Grothendieck, qui possède
des limites et suffisamment d'injectifs\footnote{Ce résultat est bien expliqué dans un
\href{https://amathew.wordpress.com/2011/07/30/quasi-coherent-sheaves-presentable-categories-and-a-result-of-gabber/}{billet de blog}
d'Akhil Mathew.}.

\begin{definition}
\label{definition-kappa-generated}
Soit $(X, \mathcal{O}_X)$ un espace annelé. Soit $\kappa$ un cardinal
infini. On dit qu'un faisceau de $\mathcal{O}_X$-modules $\mathcal{F}$ est
{\it $\kappa$-engendré} s'il existe un recouvrement ouvert
$X = \bigcup U_i$ tel que $\mathcal{F}|_{U_i}$ soit engendré par
un sous-ensemble $R_i \subset \mathcal{F}(U_i)$ de cardinal
au plus $\kappa$.
\end{definition}

\noindent
Remarquons qu'une somme directe d'au plus $\kappa$ modules $\kappa$-engendrés est
encore $\kappa$-engendrée, car $\kappa \otimes \kappa = \kappa$, voir
Ensembles, section \ref{sets-section-cardinals}.
C'est en particulier le cas de la somme directe de deux modules $\kappa$-engendrés.
De plus, un quotient d'un faisceau $\kappa$-engendré est $\kappa$-engendré.
(Il n'en va pas nécessairement de même pour les sous-modules.)

\begin{lemma}
\label{lemma-set-of-iso-classes}
Soit $(X, \mathcal{O}_X)$ un espace annelé. Soit $\kappa$ un cardinal.
Il existe un ensemble $T$ et une famille $(\mathcal{F}_t)_{t \in T}$ de modules
$\kappa$-engendrés sur $\mathcal{O}_X$ tels que tout module $\kappa$-engendré
sur $\mathcal{O}_X$ soit isomorphe à l'un des $\mathcal{F}_t$.
\end{lemma}

\begin{proof}
Les recouvrements de $X$ forment un ensemble si l'on interdit les répétitions.
Supposons que $X = \bigcup U_i$ soit un recouvrement et que $\mathcal{F}_i$
soit un $\mathcal{O}_{U_i}$-module. Les classes d'isomorphisme
de $\mathcal{O}_X$-modules $\mathcal{F}$ tels que
$\mathcal{F}|_{U_i} \cong \mathcal{F}_i$ forment alors un ensemble, car les morphismes
de recollement forment un ensemble. On se ramène donc à montrer que les classes
d'isomorphisme de quotients
$\oplus_{k \in \kappa} \mathcal{O}_X \to \mathcal{F}$
forment un ensemble pour tout espace annelé $X$. C'est clair.
\end{proof}

\noindent
Voici le résultat auquel fait référence le titre de cette section.

\begin{lemma}
\label{lemma-colimit-kappa}
Soit $X$ un schéma. Il existe un cardinal $\kappa$ tel que
tout module quasi-cohérent $\mathcal{F}$ soit la colimite filtrante
de ses sous-modules quasi-cohérents $\kappa$-engendrés.
\end{lemma}

\begin{proof}
Choisissons un recouvrement ouvert affine $X = \bigcup_{i \in I} U_i$. Pour tout couple
$i, j$, choisissons un recouvrement ouvert affine
$U_i \cap U_j = \bigcup_{k \in I_{ij}} U_{ijk}$.
Écrivons $U_i = \Spec(A_i)$ et $U_{ijk} = \Spec(A_{ijk})$.
Soit $\kappa$ un cardinal infini $\geq$ au cardinal
de chacun des ensembles $I$, $I_{ij}$.

\medskip\noindent
Soit $\mathcal{F}$ un faisceau quasi-cohérent. Posons $M_i = \mathcal{F}(U_i)$
et $M_{ijk} = \mathcal{F}(U_{ijk})$. Remarquons que
$$
M_i \otimes_{A_i} A_{ijk} = M_{ijk} = M_j \otimes_{A_j} A_{ijk}.
$$
Voir
Schémas, lemme \ref{schemes-lemma-widetilde-pullback}.
À l'aide de l'axiome du choix, choisissons une application
$$
(i, j, k, m) \mapsto S(i, j, k, m)
$$
qui associe à tous $i, j \in I$, $k \in I_{ij}$ et $m \in M_i$
un sous-ensemble fini $S(i, j, k, m) \subset M_j$ tel que l'on ait
$$
m \otimes 1 = \sum\nolimits_{m' \in S(i, j, k, m)} m' \otimes a_{m'}
$$
dans $M_{ijk}$ pour certains $a_{m'} \in A_{ijk}$. Convenons de plus
que $S(i, i, k, m) = \{m\}$ pour tous $i, j = i, k, m$ comme ci-dessus.
Fixons une telle application.

\medskip\noindent
Étant donnée une famille $\mathcal{S} = (S_i)_{i \in I}$ de sous-ensembles
$S_i \subset M_i$ de cardinal au plus $\kappa$, posons
$\mathcal{S}' = (S'_i)$, où
$$
S'_j = \bigcup\nolimits_{(i, k, m)\text{ tels que }m \in S_i}
S(i, j, k, m)
$$
On a $S_i \subset S'_i$. De plus, $S'_i$ a un cardinal au plus
$\kappa$, car c'est une réunion, indexée par un ensemble de cardinal au plus $\kappa$,
d'ensembles finis. Posons $\mathcal{S}^{(0)} = \mathcal{S}$,
$\mathcal{S}^{(1)} = \mathcal{S}'$ et, par récurrence,
$\mathcal{S}^{(n + 1)} = (\mathcal{S}^{(n)})'$. Posons ensuite
$\mathcal{S}^{(\infty)} = \bigcup_{n \geq 0} \mathcal{S}^{(n)}$.
En écrivant $\mathcal{S}^{(\infty)} = (S^{(\infty)}_i)$, on voit que,
pour tout élément $m \in S^{(\infty)}_i$, l'image de $m$ dans
$M_{ijk}$ peut s'écrire comme une somme finie $\sum m' \otimes a_{m'}$
avec $m' \in S_j^{(\infty)}$. Ainsi, en posant
$$
N_i = A_i\text{-sous-module de }M_i\text{ engendré par }S^{(\infty)}_i
$$
on a
$$
N_i \otimes_{A_i} A_{ijk} = N_j \otimes_{A_j} A_{ijk}.
$$
comme sous-modules de $M_{ijk}$. Il existe donc un sous-faisceau quasi-cohérent
$\mathcal{G} \subset \mathcal{F}$ tel que $\mathcal{G}(U_i) = N_i$.
De plus, par construction, le faisceau $\mathcal{G}$ est $\kappa$-engendré.

\medskip\noindent
Soit $\{\mathcal{G}_t\}_{t \in T}$ l'ensemble des sous-faisceaux
quasi-cohérents $\kappa$-engendrés. Si $t, t' \in T$,
alors $\mathcal{G}_t + \mathcal{G}_{t'}$ est encore un sous-faisceau
quasi-cohérent $\kappa$-engendré, étant l'image du morphisme
$\mathcal{G}_t \oplus \mathcal{G}_{t'} \to \mathcal{F}$.
Le système ordonné par inclusion est donc filtrant.
Les arguments précédents montrent que toute section de $\mathcal{F}$ sur $U_i$
appartient à l'un des $\mathcal{G}_t$ (car on peut partir d'une famille $\mathcal{S}$
telle que la section donnée soit un élément de $S_i$). Ainsi
$\colim_t \mathcal{G}_t \to \mathcal{F}$ est à la fois injectif et surjectif,
comme voulu.
\end{proof}

\begin{proposition}
\label{proposition-coherator}
Soit $X$ un schéma.
\begin{enumerate}
\item La catégorie $\QCoh(\mathcal{O}_X)$ est une catégorie abélienne
de Grothendieck. Par conséquent, $\QCoh(\mathcal{O}_X)$
possède suffisamment d'injectifs et toutes les limites.
\item Le foncteur d'inclusion
$\QCoh(\mathcal{O}_X) \to \textit{Mod}(\mathcal{O}_X)$
possède un adjoint à droite\footnote{Ce foncteur est parfois appelé
le {\it cohérateur}.}
$$
Q : \textit{Mod}(\mathcal{O}_X) \longrightarrow \QCoh(\mathcal{O}_X)
$$
tel que, pour tout faisceau quasi-cohérent $\mathcal{F}$, le morphisme d'adjonction
$Q(\mathcal{F}) \to \mathcal{F}$ soit un isomorphisme.
\end{enumerate}
\end{proposition}

\begin{proof}
L'assertion (1) signifie que $\QCoh(\mathcal{O}_X)$ (a) possède toutes les colimites,
(b) a des colimites filtrantes exactes et (c) possède un générateur, voir
Injectifs, section \ref{injectives-section-grothendieck-conditions}.
D'après Schémas, section \ref{schemes-section-quasi-coherent},
les colimites dans $\QCoh(\mathcal{O}_X)$ existent et coïncident
avec les colimites dans $\textit{Mod}(\mathcal{O}_X)$. D'après
Modules, lemme \ref{modules-lemma-limits-colimits},
les colimites filtrantes sont exactes. Ainsi (a) et (b) sont satisfaites.
Pour construire un générateur $U$, choisissons un cardinal $\kappa$ comme dans le
lemme \ref{lemma-colimit-kappa}. Choisissons une famille
$(\mathcal{F}_t)_{t \in T}$ de faisceaux quasi-cohérents $\kappa$-engendrés comme dans le
lemme \ref{lemma-set-of-iso-classes}. Posons
$U = \bigoplus_{t \in T} \mathcal{F}_t$. Comme tout objet de
$\QCoh(\mathcal{O}_X)$ est une colimite filtrante de modules quasi-cohérents
$\kappa$-engendrés, c'est-à-dire d'objets isomorphes aux $\mathcal{F}_t$,
il est clair que $U$ est un générateur.
Les assertions relatives aux limites et aux injectifs sont vraies dans toute
catégorie abélienne de Grothendieck, voir
Injectifs, théorème
\ref{injectives-theorem-injective-embedding-grothendieck} et
lemme \ref{injectives-lemma-grothendieck-products}.

\medskip\noindent
Démonstration de (2). Pour construire $Q$, on utilise le procédé général suivant.
Étant donné un objet $\mathcal{F}$ de $\textit{Mod}(\mathcal{O}_X)$,
considérons le foncteur
$$
\QCoh(\mathcal{O}_X)^{opp} \longrightarrow \textit{Sets},\quad
\mathcal{G} \longmapsto \Hom_X(\mathcal{G}, \mathcal{F})
$$
Ce foncteur transforme les colimites en limites ;
il est donc représentable, voir
Injectifs, lemme \ref{injectives-lemma-grothendieck-brown}.
Il existe ainsi un faisceau quasi-cohérent $Q(\mathcal{F})$
et un isomorphisme fonctoriel
$\Hom_X(\mathcal{G}, \mathcal{F}) = \Hom_X(\mathcal{G}, Q(\mathcal{F}))$
pour $\mathcal{G}$ dans $\QCoh(\mathcal{O}_X)$. Par le lemme de Yoneda
(Catégories, lemme \ref{categories-lemma-yoneda}),
la construction $\mathcal{F} \leadsto Q(\mathcal{F})$ est
fonctorielle en $\mathcal{F}$. Par construction, $Q$ est adjoint à droite
du foncteur d'inclusion.
Le fait que $Q(\mathcal{F}) \to \mathcal{F}$ soit un isomorphisme
lorsque $\mathcal{F}$ est quasi-cohérent découle formellement du fait
que le foncteur d'inclusion
$\QCoh(\mathcal{O}_X) \to \textit{Mod}(\mathcal{O}_X)$
est pleinement fidèle.
\end{proof}
\section{Sections à support dans un fermé}
\label{section-sections-with-support-in-closed}

\noindent
Étant donnés un espace topologique $X$, un fermé $Z \subset X$ et un
faisceau abélien $\mathcal{F}$, on peut considérer le sous-faisceau des sections dont
le support est contenu dans $Z$. Si $X$ est un schéma, $Z$ un sous-schéma
fermé et $\mathcal{F}$ un module quasi-cohérent, il existe une variante
où l'on considère les sections dont le support est schématiquement
contenu dans $Z$. Il faut toutefois être prudent dans le cadre des schémas, car
le $\mathcal{O}_X$-module obtenu n'est pas nécessairement quasi-cohérent.

\begin{lemma}
\label{lemma-quasi-coherent-finite-type-ideals}
Soit $X$ un schéma quasi-compact et quasi-séparé.
Soit $U \subset X$ un sous-schéma ouvert. Les assertions suivantes sont équivalentes :
\begin{enumerate}
\item $U$ est rétrocompact dans $X$ ;
\item $U$ est quasi-compact ;
\item $U$ est une réunion finie d'ouverts affines ;
\item il existe un faisceau d'idéaux quasi-cohérent de type fini
$\mathcal{I} \subset \mathcal{O}_X$ tel que $X \setminus U = V(\mathcal{I})$
ensemblistement.
\end{enumerate}
\end{lemma}

\begin{proof}
L'équivalence de (1), (2) et (3) résulte du
lemme \ref{lemma-quasi-separated-quasi-compact-open-retrocompact}.
Supposons (1), (2), (3). Soit $T = X \setminus U$. D'après
Schémas, lemme \ref{schemes-lemma-reduced-closed-subscheme}, il existe
un unique faisceau d'idéaux quasi-cohérent $\mathcal{J}$ définissant
la structure de sous-schéma fermé réduit induite sur $T$.
Remarquons que $\mathcal{J}|_U = \mathcal{O}_U$, qui est un
$\mathcal{O}_U$-module de type fini.
Par le lemme \ref{lemma-extend}, il existe un sous-faisceau quasi-cohérent
$\mathcal{I} \subset \mathcal{J}$ de type fini
tel que $\mathcal{I}|_U = \mathcal{J}|_U$.
Alors $X \setminus U = V(\mathcal{I})$, ce qui donne (4). Réciproquement,
si $\mathcal{I}$ est comme dans (4) et si $W = \Spec(R) \subset X$ est un ouvert
affine, alors $\mathcal{I}|_W = \widetilde{I}$ pour un idéal de type fini
$I \subset R$, voir le lemme \ref{lemma-finite-type-module}.
Il en résulte que $U \cap W = \Spec(R) \setminus V(I)$ est quasi-compact,
voir Algèbre, lemme \ref{algebra-lemma-qc-open}. Ainsi $U \subset X$
est rétrocompact par le lemme \ref{lemma-retrocompact}.
\end{proof}

\begin{lemma}
\label{lemma-sections-annihilated-by-ideal}
Soit $X$ un schéma.
Soit $\mathcal{I} \subset \mathcal{O}_X$ un faisceau d'idéaux quasi-cohérent.
Soit $\mathcal{F}$ un $\mathcal{O}_X$-module quasi-cohérent.
Considérons le faisceau de $\mathcal{O}_X$-modules $\mathcal{F}'$
qui associe à tout ouvert $U \subset X$ l'ensemble
$$
\mathcal{F}'(U)
=
\{s \in \mathcal{F}(U) \mid
\mathcal{I}s = 0\}
$$
Supposons $\mathcal{I}$ de type fini. Alors
\begin{enumerate}
\item $\mathcal{F}'$ est un faisceau quasi-cohérent de $\mathcal{O}_X$-modules ;
\item sur tout ouvert affine $U \subset X$, on a
$\mathcal{F}'(U) = \{s \in \mathcal{F}(U) \mid \mathcal{I}(U)s = 0\}$ ;
\item $\mathcal{F}'_x = \{s \in \mathcal{F}_x \mid \mathcal{I}_x s = 0\}$.
\end{enumerate}
\end{lemma}

\begin{proof}
Il est clair que la règle définissant $\mathcal{F}'$ donne un sous-faisceau
de $\mathcal{F}$ (la condition de faisceau se vérifie aisément). On peut donc
travailler localement sur $X$ pour vérifier les autres assertions. Autrement dit,
on peut supposer $X = \Spec(A)$, $\mathcal{F} = \widetilde{M}$
et $\mathcal{I} = \widetilde{I}$. Il est clair que, dans ce cas,
$\mathcal{F}'(U) = \{x \in M \mid Ix = 0\} =: M'$, car $\widetilde{I}$
est engendré par ses sections globales $I$ ; cela démontre (2).
Pour montrer que $\mathcal{F}'$ est quasi-cohérent, il suffit de montrer que,
pour tout $f \in A$, on a
$\{x \in M_f \mid I_f x = 0\} = (M')_f$.
Écrivons $I = (g_1, \ldots, g_t)$, ce qui est possible puisque $\mathcal{I}$
est de type fini, voir le lemme \ref{lemma-finite-type-module}.
Si $x = y/f^n$ et $I_fx = 0$, cela signifie que, pour tout $i$,
il existe un $m \geq 0$ tel que $f^mg_ix = 0$.
On peut choisir un même $m$ convenant à tous les $i$ ; c'est ici que l'on
utilise le fait que $I$ est de type fini. Alors $f^mx \in M'$
et $x/f^n = f^mx/f^{n + m}$ dans $(M')_f$, comme voulu.
La démonstration de (3) est analogue et omise.
\end{proof}

\begin{definition}
\label{definition-subsheaf-sections-annihilated-by-ideal}
Soit $X$ un schéma.
Soit $\mathcal{I} \subset \mathcal{O}_X$ un faisceau d'idéaux quasi-cohérent
de type fini.
Soit $\mathcal{F}$ un $\mathcal{O}_X$-module quasi-cohérent.
Le sous-faisceau $\mathcal{F}' \subset \mathcal{F}$ défini dans le
lemme \ref{lemma-sections-annihilated-by-ideal} ci-dessus est appelé
le {\it sous-faisceau des sections annulées par $\mathcal{I}$}.
\end{definition}

\begin{lemma}
\label{lemma-push-sections-annihilated-by-ideal}
Soit $f : X \to Y$ un morphisme quasi-compact et quasi-séparé
de schémas. Soit $\mathcal{I} \subset \mathcal{O}_Y$ un faisceau d'idéaux
quasi-cohérent de type fini. Soit $\mathcal{F}$ un
$\mathcal{O}_X$-module quasi-cohérent. Soit $\mathcal{F}' \subset \mathcal{F}$
le sous-faisceau des sections annulées par $f^{-1}\mathcal{I}\mathcal{O}_X$.
Alors $f_*\mathcal{F}' \subset f_*\mathcal{F}$ est le sous-faisceau
des sections annulées par $\mathcal{I}$.
\end{lemma}

\begin{proof}
Démonstration omise. (Indication : l'hypothèse que $f$ est quasi-compact et
quasi-séparé entraîne que $f_*\mathcal{F}$ est quasi-cohérent,
si bien que le lemme \ref{lemma-sections-annihilated-by-ideal} s'applique
à $\mathcal{I}$ et à $f_*\mathcal{F}$.)
\end{proof}

\noindent
Pour un faisceau abélien sur un espace topologique, nous avons étudié le sous-faisceau
des sections à support dans un fermé dans
Modules, remarque \ref{modules-remark-sections-support-in-closed}.
Pour les modules quasi-cohérents, ce sous-module n'est pas toujours
quasi-cohérent ; il l'est toutefois si le complémentaire du fermé
est rétrocompact.

\begin{lemma}
\label{lemma-sections-supported-on-closed-subset}
Soient $X$ un schéma et $Z \subset X$ un fermé.
Soit $\mathcal{F}$ un $\mathcal{O}_X$-module quasi-cohérent.
Considérons le faisceau de $\mathcal{O}_X$-modules $\mathcal{F}'$
qui associe à tout ouvert $U \subset X$ l'ensemble
$$
\mathcal{F}'(U)
=
\{s \in \mathcal{F}(U) \mid
\text{le support de }s\text{ est contenu dans }Z \cap U\}
$$
Si $X \setminus Z$ est un ouvert rétrocompact de $X$, alors
\begin{enumerate}
\item pour tout ouvert affine $U \subset X$, il existe un idéal de type fini
$I \subset \mathcal{O}_X(U)$ tel que $Z \cap U = V(I)$ ;
\item pour $U$ et $I$ comme dans (1), on a
$\mathcal{F}'(U) = \{x \in \mathcal{F}(U) \mid
I^nx = 0 \text{ pour un certain } n\}$ ;
\item $\mathcal{F}'$ est un faisceau quasi-cohérent de $\mathcal{O}_X$-modules.
\end{enumerate}
\end{lemma}

\begin{proof}
L'assertion (1) est Algèbre, lemme \ref{algebra-lemma-qc-open}.
Soient $U = \Spec(A)$ et $I$ comme dans (1).
Alors $\mathcal{F}|_U$ est le faisceau quasi-cohérent associé
à un $A$-module $M$. On a
$$
\mathcal{F}'(U) = \{x \in M \mid x = 0\text{ dans }M_\mathfrak p
\text{ pour tout }\mathfrak p \not \in Z\}.
$$
d'après Modules, définition \ref{modules-definition-support}. Ainsi
$x \in \mathcal{F}'(U)$ si et seulement si $V(\text{Ann}(x)) \subset V(I)$, voir
Algèbre, lemme \ref{algebra-lemma-support-element}. Comme $I$ est
de type fini, cela équivaut à $I^n x = 0$ pour un certain $n$.
Cela démontre (2).

\medskip\noindent
Démonstration de (3). Remarquons que, pour tout ouvert $U \subset X$, on a une suite exacte
$$
0 \to \mathcal{F}'(U) \to \mathcal{F}(U) \to \mathcal{F}(U \setminus Z)
$$
Si l'on note $j : X \setminus Z \to X$ le morphisme d'inclusion,
$\mathcal{F}(U \setminus Z)$ est l'ensemble des sections du
module $j_*(\mathcal{F}|_{X \setminus Z})$ sur $U$. On a donc
une suite exacte
$$
0 \to \mathcal{F}' \to \mathcal{F} \to j_*(\mathcal{F}|_{X \setminus Z})
$$
La restriction $\mathcal{F}|_{X \setminus Z}$ est quasi-cohérente.
Donc $j_*(\mathcal{F}|_{X \setminus Z})$ est quasi-cohérent,
d'après Schémas, lemme \ref{schemes-lemma-push-forward-quasi-coherent},
et l'hypothèse que $j$ est quasi-compact (toute immersion ouverte
est séparée). Ainsi $\mathcal{F}'$ est quasi-cohérent,
étant le noyau d'un morphisme de modules quasi-cohérents, voir
Schémas, section \ref{schemes-section-quasi-coherent}.
\end{proof}

\begin{definition}
\label{definition-subsheaf-sections-supported-on-closed}
Soit $X$ un schéma.
Soit $T \subset X$ un fermé dont le complémentaire
est rétrocompact dans $X$.
Soit $\mathcal{F}$ un $\mathcal{O}_X$-module quasi-cohérent.
Le sous-faisceau quasi-cohérent $\mathcal{F}' \subset \mathcal{F}$ défini dans le
lemme \ref{lemma-sections-supported-on-closed-subset} est appelé
le {\it sous-faisceau des sections à support dans $T$}.
\end{definition}

\begin{lemma}
\label{lemma-push-sections-supported-on-closed-subset}
Soit $f : X \to Y$ un morphisme quasi-compact et quasi-séparé
de schémas. Soit $Z \subset Y$ un fermé tel que
$Y \setminus Z$ soit rétrocompact dans $Y$. Soit $\mathcal{F}$ un
$\mathcal{O}_X$-module quasi-cohérent. Soit $\mathcal{F}' \subset \mathcal{F}$
le sous-faisceau des sections à support dans $f^{-1}Z$.
Alors $f_*\mathcal{F}' \subset f_*\mathcal{F}$ est le sous-faisceau
des sections à support dans $Z$.
\end{lemma}

\begin{proof}
Démonstration omise. (Indication : montrer d'abord que $X \setminus f^{-1}Z$ est rétrocompact
dans $X$, puisque $Y \setminus Z$ est rétrocompact dans $Y$. Ainsi,
le lemme \ref{lemma-sections-supported-on-closed-subset}
s'applique à $f^{-1}Z$ et à $\mathcal{F}$. Comme $f$ est quasi-compact et
quasi-séparé, $f_*\mathcal{F}$ est quasi-cohérent.
Le lemme \ref{lemma-sections-supported-on-closed-subset}
s'applique donc à $Z$ et à $f_*\mathcal{F}$. Enfin, identifier directement
les faisceaux.)
\end{proof}

\section{Sections des faisceaux quasi-cohérents}
\label{section-sections-quasi-coherent}

\noindent
Voici un calcul des sections d'un faisceau quasi-cohérent sur un ouvert quasi-compact
d'un spectre affine.

\begin{lemma}
\label{lemma-sections-over-quasi-compact-open-in-affine}
Soit $A$ un anneau.
Soit $I \subset A$ un idéal de type fini.
Soit $M$ un $A$-module.
Il existe alors un homomorphisme canonique
$$
\colim_n \Hom_A(I^n, M)
\longrightarrow
\Gamma(\Spec(A) \setminus V(I), \widetilde{M}).
$$
Cet homomorphisme est toujours injectif.
Si, pour tout $x \in M$, on a $Ix = 0 \Rightarrow x = 0$,
il est un isomorphisme. Dans le cas général, posons
$M_n = \{x \in M \mid I^nx = 0\}$ ; on a alors un
isomorphisme
$$
\colim_n \Hom_A(I^n, M/M_n)
\longrightarrow
\Gamma(\Spec(A) \setminus V(I), \widetilde{M}).
$$
\end{lemma}

\begin{proof}
Comme $I^{n + 1} \subset I^n$ et $M_n \subset M_{n + 1}$, on peut
composer avec ces homomorphismes pour obtenir des homomorphismes canoniques de $A$-modules
$$
\Hom_A(I^n, M)
\longrightarrow
\Hom_A(I^{n + 1}, M)
$$
et
$$
\Hom_A(I^n, M/M_n)
\longrightarrow
\Hom_A(I^{n + 1}, M/M_{n + 1})
$$
que l'on utilisera comme morphismes de transition des systèmes. Étant donné un morphisme de
$A$-modules $\varphi : I^n \to M$, on obtient un morphisme de
faisceaux $\widetilde{\varphi} : \widetilde{I^n} \to \widetilde{M}$
que l'on peut restreindre à l'ouvert $\Spec(A) \setminus V(I)$.
Comme la restriction de $\widetilde{I^n}$ à cet ouvert est le faisceau
structural, on obtient un élément de
$\Gamma(\Spec(A) \setminus V(I), \widetilde{M})$.
Nous omettons de vérifier la compatibilité avec les morphismes de transition
du système $\Hom_A(I^n, M)$. Cela donne la première flèche.
Pour obtenir la seconde, remarquons que
$\widetilde{M}$ et $\widetilde{M/M_n}$ coïncident sur l'ouvert
$\Spec(A) \setminus V(I)$, puisque le faisceau $\widetilde{M_n}$
a évidemment son support dans $V(I)$. On peut donc procéder
de la même manière.

\medskip\noindent
Décrivons maintenant cette flèche en termes algébriques.
Écrivons $I = (f_1, \ldots, f_t)$. Alors
$\Spec(A) \setminus V(I) = \bigcup_{i = 1, \ldots, t} D(f_i)$.
Ainsi la suite
$$
0 \to
\Gamma(\Spec(A) \setminus V(I), \widetilde{M}) \to
\bigoplus\nolimits_i M_{f_i} \to
\bigoplus\nolimits_{i, j} M_{f_if_j}
$$
est exacte. Soit $\varphi : I^n \to M$ un morphisme de $A$-modules.
Considérons le vecteur dont les composantes sont $\varphi(f_i^n)/f_i^n \in M_{f_i}$.
On voit aisément que son image est nulle dans la
seconde somme directe de la suite exacte ci-dessus. Il définit donc un élément
de $\Gamma(\Spec(A) \setminus V(I), \widetilde{M})$.
Nous omettons de vérifier que cette description coïncide avec celle
donnée précédemment.

\medskip\noindent
Montrons à l'aide de cette description que la première flèche est injective.
Si $\varphi$ a une image nulle, alors, pour tout $i$, l'élément
$\varphi(f_i^n)/f_i^n$ est nul dans $M_{f_i}$. Autrement dit,
pour tout $i$, on a $f_i^m\varphi(f_i^n) = 0$ pour un certain $m \geq 0$.
On peut choisir un même $m$ convenant à tous les $i$. Alors
$\varphi(f_i^{n + m}) = 0$ pour tout $i$. On voit aisément que
cela entraîne $\varphi|_{I^{t(n + m - 1) + 1}} = 0$ ; autrement
dit, $\varphi$ a une image nulle dans le terme d'indice $t(n + m - 1) + 1$
de la colimite. L'injectivité en découle.

\medskip\noindent
Remarquons que tous les $M_n = 0$ si l'on a
$Ix = 0 \Rightarrow x = 0$ pour $x \in M$. Pour
achever la démonstration du lemme, il suffit donc de montrer que
la seconde flèche est un isomorphisme.

\medskip\noindent
Cherchons à construire un inverse du second homomorphisme du lemme.
Soit $s \in \Gamma(\Spec(A) \setminus V(I), \widetilde{M})$.
Cette section correspond à un vecteur $x_i/f_i^n$, avec $x_i \in M$, de la
première somme directe de la suite exacte ci-dessus.
Ainsi, pour tous $i, j$, il existe un $m \geq 0$
tel que $f_i^m f_j^m (f_j^n x_i - f_i^n x_j) = 0$ dans $M$.
On peut choisir un même $m$ convenant à tous les couples $i, j$.
En remplaçant $x_i$ par $f_i^mx_i$ et $n$ par $n + m$,
on obtient $f_j^nx_i = f_i^nx_j$ dans $M$ pour tous $i, j$.
Introduisons
$$
K_n = \{x \in M \mid f_1^nx = \ldots = f_t^nx = 0\}
$$
Nous affirmons qu'il existe un morphisme de $A$-modules
$$
\varphi :
I^{t(n - 1) + 1}
\longrightarrow
M/K_n
$$
qui envoie le monôme
$f_1^{e_1} \ldots f_t^{e_t}$, avec $\sum e_i = t(n - 1) + 1$,
sur la classe modulo $K_n$ de l'expression
$f_1^{e_1} \ldots f_i^{e_i - n} \ldots f_t^{e_t}x_i$,
où $i$ est choisi de sorte que $e_i \geq n$ (il existe
au moins un tel $i$).
Pour le vérifier, supposons que
$$
\sum\nolimits_{E = (e_1, \ldots, e_t), |E| = t(n - 1) + 1}
a_E f_1^{e_1} \ldots f_t^{e_t} = 0
$$
soit une relation entre les monômes à coefficients $a_E$ dans $A$.
Cette relation serait envoyée sur
$$
z =
\sum\nolimits_{E = (e_1, \ldots, e_t), |E| = t(n - 1) + 1}
a_E f_1^{e_1} \ldots f_{i(E)}^{e_{i(E)} - n} \ldots f_t^{e_t}x_{i(E)}
$$
où, pour chaque multi-indice $E$, on a choisi un $i(E)$
tel que $e_{i(E)} \geq n$.
Remarquons qu'en multipliant cette expression par $f_j^n$, pour un $j$ quelconque,
on obtient zéro : en effet, les relations $f_j^nx_i = f_i^nx_j$ ci-dessus donnent
\begin{align*}
f_j^nz & = \sum\nolimits_{E = (e_1, \ldots, e_t), |E| = t(n - 1) + 1}
a_E f_1^{e_1} \ldots f_j^{e_j + n}
\ldots f_{i(E)}^{e_{i(E)} - n} \ldots f_t^{e_t}x_{i(E)} \\
& =
\sum\nolimits_{E = (e_1, \ldots, e_t), |E| = t(n - 1) + 1}
a_E f_1^{e_1} \ldots f_t^{e_t}x_j
= 0.
\end{align*}
Ainsi $z \in K_n$, et toute relation est envoyée sur zéro
dans $M/K_n$. L'assertion est démontrée.

\medskip\noindent
Remarquons que $K_n \subset M_{t(n - 1) + 1}$. Le
morphisme $\varphi$ fournit donc en particulier un morphisme de $A$-modules
$I^{t(n - 1) + 1} \to M/M_{t(n - 1) + 1}$.
Cela montre que la seconde flèche du lemme est surjective.
Nous omettons la démonstration de l'injectivité.
\end{proof}

\begin{example}
\label{example-does-not-work-in-general}
Donnons deux exemples où le premier homomorphisme affiché du
lemme \ref{lemma-sections-over-quasi-compact-open-in-affine}
n'est pas un isomorphisme.

\medskip\noindent
Soit $k$ un corps. Considérons l'anneau
$$
A = k[x, y, z_1, z_2, \ldots]/(x^nz_n).
$$
Posons $I = (x)$ et soit $M = A$. L'élément $y/x$ définit alors une section
du faisceau structural de $\Spec(A)$ sur $D(x) = \Spec(A) \setminus V(I)$.
Nous affirmons que $y/x$ n'appartient pas à l'image de l'homomorphisme canonique
$\colim \Hom_A(I^n, A) \to A_x = \mathcal{O}(D(x))$. Sinon,
il proviendrait d'un homomorphisme $\varphi : I^n \to A$ pour un certain $n$.
Posons $a = \varphi(x^n)$. On aurait alors $x^m(xa - x^ny) = 0$ pour un certain
$m > 0$, donc $x^{m + 1}a = x^{m + n}y$, puis
$\varphi(x^{n + m + 1}) = x^{m + n}y$. Cela conduit à une contradiction,
car on aurait
$$
0 = \varphi(0) = \varphi(z_{n + m + 1} x^{n + m + 1}) =
x^{m + n}y z_{n + m + 1}
$$
ce qui est faux dans l'anneau $A$.

\medskip\noindent
Soit $k$ un corps. Considérons l'anneau
$$
A = k[f, g, x, y, \{a_n, b_n\}_{n \geq 1}]/
(fy - gx, \{a_nf^n + b_ng^n\}_{n \geq 1}).
$$
Posons $I = (f, g)$ et soit $M = A$. Alors $x/f \in A_f$ et $y/g \in A_g$
ont la même image dans $A_{fg}$. Ils définissent donc une section $s$
du faisceau structural de $\Spec(A)$ sur
$D(f) \cup D(g) = \Spec(A) \setminus V(I)$. Toutefois, il n'existe aucun
$n \geq 0$ tel que $s$ provienne d'un morphisme de $A$-modules
$\varphi : I^n \to A$ appartenant à la source de la première flèche affichée
du lemme \ref{lemma-sections-over-quasi-compact-open-in-affine}.
En effet, étant donné un tel morphisme de modules, posons
$x_n = \varphi(f^n)$ et $y_n = \varphi(g^n)$.
Alors $f^mx_n = f^{n + m - 1}x$ et
$g^my_n = g^{n + m - 1}y$ pour un certain $m \geq 0$ (voir la démonstration
du lemme). Mais on aurait alors
$0 = \varphi(0) =
\varphi(a_{n + m}f^{n + m} + b_{n + m}g^{n + m}) =
a_{n + m}f^{n + m - 1}x + b_{n + m}g^{n + m - 1}y$, ce qui est faux
dans l'anneau $A$.
\end{example}

\noindent
Nous préciserons le lemme suivant dans le cas noethérien, voir
Cohomologie des schémas, lemme \ref{coherent-lemma-homs-over-open}.

\begin{lemma}
\label{lemma-sections-over-quasi-compact-open}
Soit $X$ un schéma quasi-compact.
Soit $\mathcal{I} \subset \mathcal{O}_X$ un
faisceau d'idéaux quasi-cohérent de type fini.
Soit $Z \subset X$ le sous-schéma fermé
défini par $\mathcal{I}$, et posons $U = X \setminus Z$.
Soit $\mathcal{F}$ un $\mathcal{O}_X$-module quasi-cohérent.
L'homomorphisme canonique
$$
\colim_n \Hom_{\mathcal{O}_X}(\mathcal{I}^n,
\mathcal{F})
\longrightarrow
\Gamma(U, \mathcal{F})
$$
est injectif. Supposons de plus $X$ quasi-séparé.
Soit $\mathcal{F}_n \subset \mathcal{F}$
le sous-faisceau des sections annulées par $\mathcal{I}^n$.
L'homomorphisme canonique
$$
\colim_n \Hom_{\mathcal{O}_X}(\mathcal{I}^n,
\mathcal{F}/\mathcal{F}_n)
\longrightarrow
\Gamma(U, \mathcal{F})
$$
est un isomorphisme.
\end{lemma}

\begin{proof}
Soit $\Spec(A) = W \subset X$ un ouvert affine. Écrivons
$\mathcal{F}|_W = \widetilde{M}$ pour un $A$-module $M$
et $\mathcal{I}|_W = \widetilde{I}$ pour un idéal de type fini
$I \subset A$. En restreignant le premier homomorphisme affiché
du lemme à $W$, on obtient le premier homomorphisme affiché du
lemme \ref{lemma-sections-over-quasi-compact-open-in-affine}.
Comme on peut recouvrir $X$ par un nombre fini d'ouverts affines,
cela montre que le premier homomorphisme affiché du lemme est injectif.

\medskip\noindent
On a $\mathcal{F}_n|_W = \widetilde{M_n}$, où
$M_n \subset M$ est défini comme dans le
lemme \ref{lemma-sections-over-quasi-compact-open-in-affine}
(détails omis). Ce lemme fournit une bijection
$$
\colim_n  \Hom_{\mathcal{O}_W}(
\mathcal{I}^n|_W, (\mathcal{F}/\mathcal{F}_n)|_W)
\longrightarrow
\Gamma(U \cap W, \mathcal{F})
$$
pour tout ouvert affine $W$ comme ci-dessus.

\medskip\noindent
Pour montrer que la seconde flèche affichée du lemme est bijective,
choisissons un recouvrement ouvert affine fini $X = \bigcup_{j = 1, \ldots, m} W_j$.
L'injectivité découle immédiatement de ce qui précède et de la finitude
du recouvrement. Si $X$ est quasi-séparé, pour tout couple
$j, j'$, choisissons un recouvrement ouvert affine fini
$$
W_j \cap W_{j'} = \bigcup\nolimits_{k = 1, \ldots, m_{jj'}} W_{jj'k}.
$$
Soit $s \in \Gamma(U, \mathcal{F})$. Comme on l'a vu, pour tout $j$, il existe
un $n_j$ et un morphisme
$\varphi_j : \mathcal{I}^{n_j}|_{W_j} \to
(\mathcal{F}/\mathcal{F}_{n_j})|_{W_j}$
correspondant à $s|_{U \cap W_j}$.
De même, pour tout triplet $(j, j', k)$, il existe un entier
$n_{jj'k}$ tel que les restrictions de $\varphi_j$ et de $\varphi_{j'}$,
comme morphismes $\mathcal{I}^{n_{jj'k}} \to \mathcal{F}/\mathcal{F}_{n_{jj'k}}$,
coïncident sur $W_{jj'k}$. Posons $n = \max\{n_j, n_{jj'k}\}$ ; alors
les $\varphi_j$ se recollent en un morphisme
$\mathcal{I}^n \to \mathcal{F}/\mathcal{F}_n$ sur $X$.
Cela démontre la surjectivité de l'homomorphisme.
\end{proof}

\section{Faisceaux inversibles amples}
\label{section-ample}

\noindent
Rappelons d'après Modules, lemme \ref{modules-lemma-s-open},
qu'étant donnés un faisceau inversible $\mathcal{L}$ sur un espace
localement annelé $X$ et une section globale $s$ de $\mathcal{L}$,
l'ensemble $X_s = \{x \in X \mid s \not \in \mathfrak m_x\mathcal{L}_x\}$
est ouvert. On remarque en général que
$X_s \cap X_{s'} = X_{ss'}$, où $ss'$ désigne
la section $s \otimes s' \in \Gamma(X, \mathcal{L} \otimes \mathcal{L}')$.

\begin{definition}
\label{definition-ample}
\begin{reference}
\cite[II Definition 4.5.3]{EGA}
\end{reference}
Soit $X$ un schéma.
Soit $\mathcal{L}$ un $\mathcal{O}_X$-module inversible.
On dit que $\mathcal{L}$ est {\it ample} si
\begin{enumerate}
\item $X$ est quasi-compact, et
\item pour tout $x \in X$, il existe un entier $n \geq 1$
et $s \in \Gamma(X, \mathcal{L}^{\otimes n})$ tels
que $x \in X_s$ et que $X_s$ soit affine.
\end{enumerate}
\end{definition}

\begin{lemma}
\label{lemma-ample-power-ample}
\begin{reference}
\cite[II Proposition 4.5.6(i)]{EGA}
\end{reference}
Soit $X$ un schéma. Soit $\mathcal{L}$ un $\mathcal{O}_X$-module inversible.
Soit $n \geq 1$. Alors $\mathcal{L}$ est ample si et seulement si
$\mathcal{L}^{\otimes n}$ est ample.
\end{lemma}

\begin{proof}
Cela résulte de l'égalité $X_{s^n} = X_s$.
\end{proof}

\begin{lemma}
\label{lemma-ample-on-closed}
Soit $X$ un schéma.
Soit $\mathcal{L}$ un $\mathcal{O}_X$-module inversible ample.
Pour tout sous-schéma fermé $Z \subset X$, la restriction de
$\mathcal{L}$ à $Z$ est ample.
\end{lemma}

\begin{proof}
C'est clair, car une partie fermée d'un espace quasi-compact est quasi-compacte
et un sous-schéma fermé d'un schéma affine est affine (voir
Schémas, lemme \ref{schemes-lemma-closed-immersion-affine-case}).
\end{proof}

\begin{lemma}
\label{lemma-affine-cap-s-open}
Soit $X$ un schéma. Soit $\mathcal{L}$ un $\mathcal{O}_X$-module inversible.
Soit $s \in \Gamma(X, \mathcal{L})$. Pour tout ouvert affine $U \subset X$,
l'intersection $U \cap X_s$ est affine.
\end{lemma}

\begin{proof}
Cela se traduit par le problème d'algèbre suivant.
Soit $R$ un anneau. Soit $N$ un $R$-module inversible
(c'est-à-dire localement libre de rang 1). Soit $s \in N$.
Alors $V = \{\mathfrak p \mid s \not \in \mathfrak p N\}$ est
un ouvert affine de $\Spec(R)$.

\medskip\noindent
Soit $A = \bigoplus_{n \geq 0} A_n$ l'algèbre symétrique de $N$
(qui est commutative) et considérons $s$ comme un élément de $A_1$.
Posons $B = A/(s - 1)A$. C'est une $R$-algèbre dont la construction
commute à tout changement de base $R \to R'$. Ainsi $B' = B \otimes_R R'$
est l'anneau nul si $s$ a pour image zéro dans $N' = N \otimes_R R'$.
Il s'ensuit que, si $x \in \Spec(R) \setminus V$, alors $B \otimes_R \kappa(x) = 0$.
On en conclut que $\Spec(B) \to \Spec(R)$ se factorise par $V$, car les
fibres aux points $x \not \in V$ sont vides. D'autre part, si
$\Spec(R') \subset V$ est un ouvert affine, alors $s$ a pour image un élément
de base de $N'$ et l'on voit que $B' = R'[s]/(s - 1) \cong R'$.
Il en résulte que $\Spec(B) \to V$ est un isomorphisme et que $V$ est
bien affine.
\end{proof}

\begin{lemma}
\label{lemma-ample-tensor-globally-generated}
\begin{reference}
\cite[II Proposition 4.5.6(ii)]{EGA}
\end{reference}
Soit $X$ un schéma. Soient $\mathcal{L}$ et $\mathcal{M}$
des $\mathcal{O}_X$-modules inversibles. Si
\begin{enumerate}
\item $\mathcal{L}$ est ample, et
\item les ouverts $X_t$, où $t \in \Gamma(X, \mathcal{M}^{\otimes m})$
pour $m > 0$, recouvrent $X$,
\end{enumerate}
alors $\mathcal{L} \otimes \mathcal{M}$ est ample.
\end{lemma}

\begin{proof}
Vérifions les conditions de la définition \ref{definition-ample}.
Comme $\mathcal{L}$ est ample, $X$ est quasi-compact.
Soit $x \in X$. Choisissons $n \geq 1$, $m \geq 1$,
$s \in \Gamma(X, \mathcal{L}^{\otimes n})$ et
$t \in \Gamma(X, \mathcal{M}^{\otimes m})$
tels que $x \in X_s$, $x \in X_t$ et que $X_s$ soit affine.
Alors $s^mt^n \in \Gamma(X, (\mathcal{L} \otimes \mathcal{M})^{\otimes nm})$,
$x \in X_{s^mt^n}$, et $X_{s^mt^n}$ est affine d'après
le lemme \ref{lemma-affine-cap-s-open}.
\end{proof}

\begin{lemma}
\label{lemma-affine-s-opens}
Soit $X$ un schéma. Soit $\mathcal{L}$ un $\mathcal{O}_X$-module inversible.
Supposons que les ouverts $X_s$, où $s \in \Gamma(X, \mathcal{L}^{\otimes n})$
et $n \geq 1$, forment une base de la topologie de $X$.
Alors, parmi ces ouverts, ceux des $X_s$ qui sont affines
forment une base de la topologie de $X$.
\end{lemma}

\begin{proof}
Soit $x \in X$. Choisissons un voisinage ouvert affine
$\Spec(R) = U \subset X$ de $x$.
Par hypothèse, il existe
un entier $n \geq 1$ et une section $s \in \Gamma(X, \mathcal{L}^{\otimes n})$
tels que $X_s \subset U$. D'après le lemme \ref{lemma-affine-cap-s-open} ci-dessus,
l'intersection $X_s = U \cap X_s$ est affine.
Comme on peut choisir $U$ arbitrairement petit, cela donne la conclusion.
\end{proof}

\begin{lemma}
\label{lemma-affine-s-opens-cover-quasi-separated}
Soient $X$ un schéma et $\mathcal{L}$ un $\mathcal{O}_X$-module inversible.
Supposons que, pour tout point $x$ de $X$, il existe $n \geq 1$ et
$s \in \Gamma(X, \mathcal{L}^{\otimes n})$ tels que
$x \in X_s$ et que $X_s$ soit affine. Alors $X$ est séparé.
\end{lemma}

\begin{proof}
Montrons d'abord que $X$ est quasi-séparé. Par hypothèse, on
peut trouver un recouvrement de $X$ par des ouverts affines de la forme
$X_s$. D'après le lemme \ref{lemma-affine-cap-s-open},
l'intersection de deux tels ouverts est affine ; Schémas, lemme
\ref{schemes-lemma-characterize-quasi-separated}, entraîne donc
que $X$ est quasi-séparé.

\medskip\noindent
Pour montrer que $X$ est séparé, on peut utiliser le critère
valuatif, Schémas, lemme
\ref{schemes-lemma-valuative-criterion-separatedness}.
Soit donc $A$ un anneau de valuation de corps des fractions $K$,
et considérons deux morphismes $f, g : \Spec(A) \to X$ tels que
les deux composés $\Spec(K) \to \Spec(A) \to X$ coïncident.
Comme $A$ est local, il existe $p, q \ge 1$, $s \in \Gamma(X,
\mathcal{L}^{\otimes p})$ et $t \in \Gamma(X,
\mathcal{L}^{\otimes q})$ tels que $X_s$ et $X_t$ soient
affines, $f(\Spec A) \subseteq X_s$ et $g(\Spec A)
\subseteq X_t$. Remplaçons maintenant $s$ par $s^q$, $t$ par $t^p$
et $\mathcal{L}$ par $\mathcal{L}^{\otimes pq}$. Cela ne change
rien, puisque $X_s = X_{s^q}$ et $X_t = X_{t^p}$, et désormais $s$
et $t$ sont des sections du même faisceau $\mathcal{L}$.

\medskip\noindent
Le module quasi-cohérent $f^*\mathcal{L}$ correspond à un $A$-module $M$ et
$g^*\mathcal{L}$ correspond à un $A$-module $N$, d'après notre
classification des modules quasi-cohérents sur les schémas affines
(Schémas, lemme \ref{schemes-lemma-quasi-coherent-affine}).
Les $A$-modules $M$ et $N$ sont localement libres de rang
$1$ (lemme \ref{lemma-locally-free-module}) et, puisque $A$ est
local, ils sont libres (Algèbre, lemme
\ref{algebra-lemma-K0-local}). On peut donc identifier
$M$ et $N$ à des $A$-sous-modules de $M \otimes_A K$ et $N
\otimes_A K$. L'égalité $f|_{\Spec(K)} = g|_{\Spec(K)}$
détermine un isomorphisme $\phi \colon M \otimes_A K \to N
\otimes_A K$.

\medskip\noindent
Soient $x \in M$ et $y \in N$ les éléments correspondant aux
images réciproques de $s$ par $f$ et $g$, respectivement. Ils
vérifient $\phi(x \otimes 1) = y \otimes 1$. L'image de $f$
est contenue dans $X_s$, donc $x \not\in \mathfrak{m}_A M$,
autrement dit $x$ engendre $M$. Ainsi $\phi$ détermine un
isomorphisme de $M$ sur le sous-module de $N$ engendré par
$y$. En raisonnant symétriquement à l'aide de $t$, on voit que $\phi^{-1}$
détermine un isomorphisme de $N$ sur un sous-module de $M$.
Par conséquent, $\phi$ se restreint à un isomorphisme de $M$ sur
$N$. Puisque $x$ engendre $M$, son image $y$ engendre
$N$, d'où $y \not\in \mathfrak{m}_A N$. On a donc
$g(\Spec(A)) \subseteq X_s$. Comme $X_s$ est affine, il est
séparé d'après Schémas, lemme \ref{schemes-lemma-affine-separated},
et l'on en conclut que $f = g$.
\end{proof}

\begin{lemma}
\label{lemma-ample-separated}
Soit $X$ un schéma. S'il existe un faisceau inversible ample sur $X$,
alors $X$ est séparé.
\end{lemma}

\begin{proof}
Cela résulte immédiatement du
lemme \ref{lemma-affine-s-opens-cover-quasi-separated} et de la
définition \ref{definition-ample}.
\end{proof}

\begin{lemma}
\label{lemma-map-into-proj}
Soit $X$ un schéma.
Soit $\mathcal{L}$ un $\mathcal{O}_X$-module inversible.
Posons $S = \Gamma_*(X, \mathcal{L})$, considéré comme anneau gradué.
Si tout point de $X$ appartient à l'un des
ouverts $X_s$, pour un élément homogène $s \in S_{+}$, alors
il existe un morphisme canonique de schémas
$$
f : X \longrightarrow Y = \text{Proj}(S),
$$
à valeurs dans le spectre homogène de $S$ (voir
Constructions, section \ref{constructions-section-proj}).
Ce morphisme possède les propriétés suivantes :
\begin{enumerate}
\item $f^{-1}(D_{+}(s)) = X_s$ pour tout élément homogène $s \in S_{+}$,
\item il existe des homomorphismes de $\mathcal{O}_X$-modules
$f^*\mathcal{O}_Y(n) \to \mathcal{L}^{\otimes n}$
compatibles aux applications de multiplication, voir
Constructions, équation (\ref{constructions-equation-multiply}),
\item le composé
$S_n \to \Gamma(Y, \mathcal{O}_Y(n)) \to \Gamma(X, \mathcal{L}^{\otimes n})$
est l'application identique, et
\item pour tout $x \in X$, il existe un entier $d \geq 1$
et un voisinage ouvert $U \subset X$ de $x$
tels que $f^*\mathcal{O}_Y(dn)|_U \to \mathcal{L}^{\otimes dn}|_U$
soit un isomorphisme pour tout $n \in \mathbf{Z}$.
\end{enumerate}
\end{lemma}

\begin{proof}
Notons $\psi : S \to \Gamma_*(X, \mathcal{L})$ l'application identique.
Nous allons utiliser le triplet
$(U(\psi), r_{\mathcal{L}, \psi}, \theta)$ de
Constructions, lemme \ref{constructions-lemma-invertible-map-into-proj}.
Par hypothèse, le sous-schéma ouvert $U(\psi)$ est égal à $X$. Ainsi
$r_{\mathcal{L}, \psi} : U(\psi) \to Y$ est défini sur tout $X$.
Posons $f = r_{\mathcal{L}, \psi}$.
Les applications de (2) sont les composantes de $\theta$.
L'assertion (3) résulte de la condition (2) du lemme cité ci-dessus.
L'assertion (1) résulte de (3), combiné à la condition (1) du lemme
cité ci-dessus. L'assertion (4) résulte de la dernière assertion de
Constructions, lemme \ref{constructions-lemma-invertible-map-into-proj},
puisque l'application $\alpha$ qui y figure est un isomorphisme.
\end{proof}

\begin{lemma}
\label{lemma-map-into-proj-quasi-compact}
Soit $X$ un schéma. Soit $\mathcal{L}$ un $\mathcal{O}_X$-module inversible.
Posons $S = \Gamma_*(X, \mathcal{L})$.
Supposons (a) que tout point de $X$ appartienne à l'un des
ouverts $X_s$, pour un élément homogène $s \in S_{+}$,
et (b) que $X$ soit quasi-compact. Alors le morphisme canonique de schémas
$f : X \longrightarrow \text{Proj}(S)$ du lemme \ref{lemma-map-into-proj}
ci-dessus est quasi-compact et d'image dense.
\end{lemma}

\begin{proof}
Pour montrer que $f$ est quasi-compact, il suffit de montrer que $f^{-1}(D_{+}(s))$
est quasi-compact pour tout élément homogène $s \in S_{+}$. Écrivons
$X = \bigcup_{i = 1, \ldots, n} X_i$ comme réunion finie
d'ouverts affines. D'après le lemme \ref{lemma-affine-cap-s-open}, chaque intersection
$X_s \cap X_i$ est affine. Ainsi $X_s = \bigcup_{i = 1, \ldots, n} X_s \cap X_i$
est quasi-compact. Supposons que l'image de $f$ ne soit pas dense et cherchons
une contradiction. Puisque les ouverts $D_+(s)$, avec $s \in S_+$ homogène,
forment une base de la topologie de $\text{Proj}(S)$, on peut trouver
un tel $s$ avec $D_+(s) \not = \emptyset$ et $f(X) \cap D_+(s) = \emptyset$.
D'après le lemme \ref{lemma-map-into-proj},
cela signifie que $X_s = \emptyset$. D'après le lemme \ref{lemma-invert-s-sections},
une puissance $s^n$ est alors la section nulle de
$\mathcal{L}^{\otimes n\deg(s)}$.
Cela entraîne à son tour que $D_+(s) = \emptyset$, contradiction
cherchée.
\end{proof}

\begin{lemma}
\label{lemma-ample-immersion-into-proj}
Soit $X$ un schéma. Soit $\mathcal{L}$ un $\mathcal{O}_X$-module inversible.
Posons $S = \Gamma_*(X, \mathcal{L})$.
Supposons $\mathcal{L}$ ample. Alors le morphisme canonique de schémas
$f : X \longrightarrow \text{Proj}(S)$ du lemme \ref{lemma-map-into-proj}
est une immersion ouverte d'image dense.
\end{lemma}

\begin{proof}
D'après le lemme \ref{lemma-affine-s-opens-cover-quasi-separated},
$X$ est quasi-séparé. Choisissons un nombre fini d'éléments
homogènes $s_1, \ldots, s_n \in S_{+}$
tels que les $X_{s_i}$ soient affines et que $X = \bigcup X_{s_i}$.
Supposons $s_i$ de degré $d_i$. L'image réciproque de
$D_{+}(s_i)$ par $f$ est $X_{s_i}$, voir le lemme \ref{lemma-map-into-proj}.
D'après le lemme \ref{lemma-invert-s-sections}, l'homomorphisme d'anneaux
$$
(S^{(d_i)})_{(s_i)} = \Gamma(D_{+}(s_i), \mathcal{O}_{\text{Proj}(S)})
\longrightarrow
\Gamma(X_{s_i}, \mathcal{O}_X)
$$
est un isomorphisme. Ainsi $f$ induit un isomorphisme
$X_{s_i} \to D_{+}(s_i)$. Donc $f$ est un isomorphisme de $X$ sur le sous-schéma
ouvert $\bigcup_{i = 1, \ldots, n} D_{+}(s_i)$ de $\text{Proj}(S)$.
Son image est dense d'après le lemme \ref{lemma-map-into-proj-quasi-compact}.
\end{proof}


\begin{lemma}
\label{lemma-open-in-proj-ample}
Soit $X$ un schéma.
Soit $S$ un anneau gradué. Supposons $X$ quasi-compact
et supposons qu'il existe une immersion ouverte
$$
j : X \longrightarrow Y = \text{Proj}(S).
$$
Alors $j^*\mathcal{O}_Y(d)$ est un faisceau inversible ample
pour un entier $d > 0$.
\end{lemma}

\begin{proof}
C'est Constructions, lemme \ref{constructions-lemma-ample-on-proj}.
\end{proof}

\begin{proposition}
\label{proposition-characterize-ample}
Soit $X$ un schéma quasi-compact.
Soit $\mathcal{L}$ un faisceau inversible sur $X$.
Posons $S = \Gamma_*(X, \mathcal{L})$.
Les conditions suivantes sont équivalentes :
\begin{enumerate}
\item
\label{item-ample}
$\mathcal{L}$ est ample,
\item
\label{item-immersion}
les ouverts $X_s$, avec $s \in S_{+}$ homogène,
recouvrent $X$ et le morphisme associé $X \to \text{Proj}(S)$
est une immersion ouverte,
\item
\label{item-s-basis}
les ouverts $X_s$, avec $s \in S_{+}$ homogène,
forment une base de la topologie de $X$,
\item
\label{item-s-affine-basis}
les ouverts $X_s$, avec $s \in S_{+}$ homogène,
qui sont affines forment une base de la topologie de $X$,
\item
\label{item-qc-gg}
pour tout faisceau quasi-cohérent $\mathcal{F}$ sur $X$,
la somme des images des applications canoniques
$$
\Gamma(X, \mathcal{F} \otimes_{\mathcal{O}_X} \mathcal{L}^{\otimes n})
\otimes_{\mathbf{Z}} \mathcal{L}^{\otimes -n}
\longrightarrow
\mathcal{F}
$$
pour $n \geq 1$ est égale à $\mathcal{F}$,
\item
\label{item-qc-i-gg}
même propriété qu'en (\ref{item-qc-gg}), où $\mathcal{F}$
parcourt tous les faisceaux quasi-cohérents d'idéaux,
\item
\label{item-c-gg}
$X$ est quasi-séparé et,
pour tout faisceau quasi-cohérent $\mathcal{F}$ de type fini sur $X$,
il existe un entier $n_0$ tel que
$\mathcal{F} \otimes_{\mathcal{O}_X} \mathcal{L}^{\otimes n}$
soit engendré par ses sections globales pour tout $n \geq n_0$,
\item
\label{item-c-q}
$X$ est quasi-séparé et,
pour tout faisceau quasi-cohérent $\mathcal{F}$ de type fini sur $X$,
il existe des entiers $n > 0$, $k \geq 0$ tels que
$\mathcal{F}$ soit quotient d'une somme directe de $k$ exemplaires de
$\mathcal{L}^{\otimes - n}$, et
\item
\label{item-c-i-q}
même propriété qu'en (\ref{item-c-q}), où $\mathcal{F}$ parcourt tous
les faisceaux d'idéaux de type fini sur $X$.
\end{enumerate}
\end{proposition}

\begin{proof}
Le lemme \ref{lemma-ample-immersion-into-proj} donne
(\ref{item-ample}) $\Rightarrow$ (\ref{item-immersion}).
Les lemmes \ref{lemma-ample-power-ample} et \ref{lemma-open-in-proj-ample}
donnent l'implication
(\ref{item-ample}) $\Leftarrow$ (\ref{item-immersion}).
Les implications (\ref{item-immersion}) $\Rightarrow$
(\ref{item-s-affine-basis}) $\Rightarrow$ (\ref{item-s-basis})
résultent clairement de Constructions, section \ref{constructions-section-proj}.
Le lemme \ref{lemma-affine-s-opens} donne
(\ref{item-s-basis}) $\Rightarrow$ (\ref{item-ample}).
Les quatre premières conditions sont donc équivalentes.

\medskip\noindent
Supposons vérifiées les conditions équivalentes (1) -- (4).
Notons qu'en particulier $X$ est séparé (comme sous-schéma
ouvert du schéma séparé $\text{Proj}(S)$).
Soit $\mathcal{F}$ un faisceau quasi-cohérent sur $X$.
Choisissons $s \in S_{+}$ homogène tel que $X_s$ soit affine.
Nous affirmons que toute section $m \in \Gamma(X_s, \mathcal{F})$
appartient à l'image de l'une des applications figurant dans
(\ref{item-qc-gg}) ci-dessus. Cela entraînera (\ref{item-qc-gg}),
puisque ces ouverts affines $X_s$ recouvrent $X$.
En effet, d'après le lemme \ref{lemma-invert-s-sections}, on peut écrire
$m$ comme image de $m' \otimes s^{-n}$ pour un certain
$n \geq 1$ et un certain
$m' \in \Gamma(X, \mathcal{F} \otimes \mathcal{L}^{\otimes n})$.
Cela prouve l'affirmation.

\medskip\noindent
Clairement, (\ref{item-qc-gg}) $\Rightarrow$ (\ref{item-qc-i-gg}).
Supposons (\ref{item-qc-i-gg}) et montrons que $\mathcal{L}$ est
ample. Choisissons $x \in X$. Soit $U \subset X$ un ouvert affine
contenant $x$. Posons $Z = X \setminus U$. On peut considérer
$Z$ comme un sous-schéma fermé réduit, voir
Schémas, section \ref{schemes-section-reduced}.
Soit $\mathcal{I} \subset \mathcal{O}_X$ le faisceau quasi-cohérent
d'idéaux correspondant au sous-schéma fermé $Z$.
D'après l'hypothèse (\ref{item-qc-i-gg}), il existe un entier $n \geq 1$ et une section
$s \in \Gamma(X, \mathcal{I} \otimes \mathcal{L}^{\otimes n})$
tels que $s$ ne s'annule pas en $x$ (plus précisément, tels que
$s \not \in \mathfrak m_x \mathcal{I}_x \otimes \mathcal{L}_x^{\otimes n}$).
On peut considérer $s$ comme une section de $\mathcal{L}^{\otimes n}$.
Puisqu'elle s'annule clairement le long de $Z$, on a
$X_s \subset U$. Ainsi $X_s$ est affine, voir
le lemme \ref{lemma-affine-cap-s-open}.
Cela prouve que $\mathcal{L}$ est ample.
Nous avons donc montré à ce stade que (1) -- (6) sont équivalentes.

\medskip\noindent
Supposons vérifiées les conditions équivalentes (1) -- (6). Dans la suite,
nous utiliserons sans autre mention le fait que le produit tensoriel de deux
faisceaux de modules engendrés par leurs sections globales est lui-même
engendré par ses sections globales (voir
Modules, lemme \ref{modules-lemma-tensor-product-globally-generated}).
D'après (1), on peut trouver des éléments $s_i \in S_{d_i}$ avec $d_i \geq 1$
tels que $X = \bigcup_{i = 1, \ldots, n} X_{s_i}$.
Posons $d = d_1\ldots d_n$. Il en résulte que $\mathcal{L}^{\otimes d}$
est engendré par les sections globales
$$
s_1^{d/d_1}, \ldots, s_n^{d/d_n}.
$$
Ainsi, si $\mathcal{L}^{\otimes j}$ est engendré par ses sections globales,
il en est de même de $\mathcal{L}^{\otimes j + dn}$ pour tout $n \geq 0$.
Fixons $j \in \{0, \ldots, d - 1\}$. Pour tout point $x \in X$, il
existe un entier $n \geq 1$ et une section globale $s$ de $\mathcal{L}^{j + dn}$
qui ne s'annule pas en $x$, comme il résulte de (\ref{item-qc-gg}) appliqué
à $\mathcal{F} = \mathcal{L}^{\otimes j}$ et au faisceau inversible ample
$\mathcal{L}^{\otimes d}$. Comme $X$ est quasi-compact,
on peut trouver une famille finie d'entiers $n_i$ et de sections globales
$s_i$ de $\mathcal{L}^{\otimes j + dn_i}$ qui n'ont aucun zéro commun
sur $X$. Puisque $\mathcal{L}^{\otimes d}$ est engendré par ses sections globales,
$\mathcal{L}^{\otimes j + dn}$ est engendré par ses sections globales pour $n = \max\{n_i\}$.
Cela étant établi pour chaque classe de congruence modulo $d$,
on en conclut qu'il existe un entier $n_0 = n_0(\mathcal{L})$ tel que
$\mathcal{L}^{\otimes n}$ soit engendré par ses sections globales pour tout $n \geq n_0$.
On voit dès lors que, si $\mathcal{F}$ est engendré par ses sections globales,
il en est de même de $\mathcal{F} \otimes \mathcal{L}^{\otimes n}$ pour tout
$n \geq n_0$.

\medskip\noindent
Supposons toujours vérifiées les conditions équivalentes (1) -- (6).
Soit $\mathcal{F}$ un faisceau quasi-cohérent
de $\mathcal{O}_X$-modules de type fini.
Notons $\mathcal{F}_n \subset \mathcal{F}$ l'image de l'application
canonique de (\ref{item-qc-gg}). Par construction,
$\mathcal{F}_n \otimes \mathcal{L}^{\otimes n}$ est
engendré par ses sections globales. D'après (\ref{item-qc-gg}),
$\mathcal{F}$ est la somme des sous-faisceaux $\mathcal{F}_n$,
$n \geq 1$. D'après
Modules, lemme \ref{modules-lemma-finite-type-quasi-compact-colimit},
on a $\mathcal{F} = \sum_{n = 1, \ldots, N} \mathcal{F}_n$
pour un entier $N \geq 1$. Il en résulte que
$\mathcal{F} \otimes \mathcal{L}^{\otimes n}$ est engendré par
ses sections globales dès que $n \geq N + n_0(\mathcal{L})$, où $n_0(\mathcal{L})$
est défini ci-dessus. Ainsi (1) -- (6) entraînent (\ref{item-c-gg}).

\medskip\noindent
Supposons (\ref{item-c-gg}). Soit $\mathcal{F}$ un faisceau quasi-cohérent
de $\mathcal{O}_X$-modules de type fini.
D'après (\ref{item-c-gg}), il existe un entier $n \geq 1$ tel que
l'application canonique
$$
\Gamma(X, \mathcal{F} \otimes_{\mathcal{O}_X} \mathcal{L}^{\otimes n})
\otimes_{\mathbf{Z}} \mathcal{L}^{\otimes -n}
\longrightarrow
\mathcal{F}
$$
soit surjective. Soit $I$ l'ensemble des parties finies de
$\Gamma(X, \mathcal{F} \otimes_{\mathcal{O}_X} \mathcal{L}^{\otimes n})$,
ordonné par inclusion. Alors $I$ est un ensemble ordonné filtrant.
Pour $i = \{s_1, \ldots, s_{r(i)}\}$, soit $\mathcal{F}_i \subset \mathcal{F}$
l'image de l'application
$$
\bigoplus\nolimits_{j = 1, \ldots, r(i)} \mathcal{L}^{\otimes -n}
\longrightarrow
\mathcal{F}
$$
qui est la multiplication par $s_j$ sur le $j$-ième facteur. La surjectivité ci-dessus
entraîne que $\mathcal{F} = \colim_{i \in I} \mathcal{F}_i$.
On peut donc appliquer Modules, lemme \ref{modules-lemma-finite-type-quasi-compact-colimit},
et l'on en conclut que
$\mathcal{F} = \mathcal{F}_i$ pour un certain $i$.
Nous avons donc établi (\ref{item-c-q}). Autrement dit,
(\ref{item-c-gg}) $\Rightarrow$ (\ref{item-c-q}).

\medskip\noindent
L'implication (\ref{item-c-q}) $\Rightarrow$ (\ref{item-c-i-q}) est immédiate.

\medskip\noindent
Supposons enfin (\ref{item-c-i-q}).
Soit $\mathcal{I} \subset \mathcal{O}_X$ un faisceau quasi-cohérent
d'idéaux. D'après le lemme \ref{lemma-quasi-coherent-colimit-finite-type}
(c'est ici qu'intervient l'hypothèse que $X$ est quasi-séparé),
on a $\mathcal{I} = \colim_\alpha I_\alpha$, où
chaque $I_\alpha$ est quasi-cohérent de type fini. Puisque, par hypothèse,
chaque $I_\alpha$ est quotient d'une somme de puissances tensorielles négatives de
$\mathcal{L}$, il en est de même de $\mathcal{I}$ (mais, bien entendu,
sans finitude ni borne sur les puissances). On en conclut donc
que (\ref{item-c-i-q}) entraîne (\ref{item-qc-i-gg}).
Cela achève la démonstration de la proposition.
\end{proof}

\begin{lemma}
\label{lemma-ample-on-locally-closed}
Soit $X$ un schéma. Soit $\mathcal{L}$ un
$\mathcal{O}_X$-module inversible ample. Soit $i : X' \to X$ un morphisme de schémas.
Supposons vérifiée au moins l'une des conditions suivantes :
\begin{enumerate}
\item $i$ est une immersion quasi-compacte,
\item $X'$ est quasi-compact et $i$ est une immersion,
\item $i$ est quasi-compact et induit un homéomorphisme
entre $X'$ et $i(X')$,
\item $X'$ est quasi-compact et $i$ induit un homéomorphisme
entre $X'$ et $i(X')$.
\end{enumerate}
Alors $i^*\mathcal{L}$ est ample sur $X'$.
\end{lemma}

\begin{proof}
Observons que, dans les cas (1) et (3), le schéma $X'$ est quasi-compact,
puisque $X$ est quasi-compact d'après la définition \ref{definition-ample}.
Il suffit donc de traiter les cas (2) et (4). Comme (2) est un cas particulier
de (4), il suffit de traiter (4).

\medskip\noindent
Supposons vérifiée la condition (4). Pour $s \in \Gamma(X, \mathcal{L}^{\otimes d})$,
notons $s' = i^*s$ l'image réciproque de $s$ sur $X'$. Alors
$s'$ est une section de $(i^*\mathcal{L})^{\otimes d}$. D'après la
proposition \ref{proposition-characterize-ample},
les ouverts $X_s$, pour $s \in \Gamma(X, \mathcal{L}^{\otimes d})$,
forment une base de la topologie de $X$. Comme $X'_{s'} = i^{-1}(X_s)$ d'après
Modules, remarque \ref{modules-remark-pullback-s-open},
et comme $X' \to i(X')$ est un homéomorphisme, on en conclut
que les ouverts $X'_{s'}$ forment une base de la topologie de $X'$. Ainsi
$i^*\mathcal{L}$ est ample d'après la
proposition \ref{proposition-characterize-ample}.
\end{proof}

\begin{lemma}
\label{lemma-ample-on-product}
Soit $S$ un schéma quasi-séparé. Soient $X$, $Y$ des schémas sur $S$.
Soit $\mathcal{L}$ un $\mathcal{O}_X$-module inversible ample
et soit $\mathcal{N}$ un $\mathcal{O}_Y$-module inversible ample.
Alors $\mathcal{M} = \text{pr}_1^*\mathcal{L}
\otimes_{\mathcal{O}_{X \times_S Y}} \text{pr}_2^*\mathcal{N}$
est un faisceau inversible ample sur $X \times_S Y$.
\end{lemma}

\begin{proof}
Le morphisme $i : X \times_S Y \to X \times Y$ est une immersion
quasi-compacte, voir Schémas, lemme \ref{schemes-lemma-fibre-product-after-map}.
D'autre part, $\mathcal{M}$ est l'image réciproque par
$i$ du module inversible correspondant sur $X \times Y$.
D'après le lemme \ref{lemma-ample-on-locally-closed}, il suffit de démontrer le
lemme pour $X \times Y$. Vérifions (1) et (2) de la
définition \ref{definition-ample} pour $\mathcal{M}$ sur $X \times Y$.

\medskip\noindent
Puisque $X$ et $Y$ sont quasi-compacts, il en est de même de $X \times Y$.
Soit $z \in X \times Y$ un point. Soient $x \in X$ et $y \in Y$
ses projections. Choisissons $n > 0$ et
$s \in \Gamma(X, \mathcal{L}^{\otimes n})$
tels que $X_s$ soit un voisinage ouvert affine de $x$.
Choisissons $m > 0$ et
$t \in \Gamma(Y, \mathcal{N}^{\otimes m})$
tels que $Y_t$ soit un voisinage ouvert affine de $y$.
Alors $r = \text{pr}_1^*s \otimes \text{pr}_2^*t$ est une section
de $\mathcal{M}$ avec $(X \times Y)_r = X_s \times Y_t$.
C'est un voisinage ouvert affine de $z$, ce qui achève la démonstration.
\end{proof}


\section{Schémas affines et quasi-affines}
\label{section-affine-quasi-affine}

\begin{lemma}
\label{lemma-quasi-affine-O-ample}
Soit $X$ un schéma.
Alors $X$ est quasi-affine si et seulement si $\mathcal{O}_X$ est ample.
\end{lemma}

\begin{proof}
Supposons $X$ quasi-affine. Posons $A = \Gamma(X, \mathcal{O}_X)$.
Considérons l'immersion ouverte
$$
j : X \longrightarrow \Spec(A)
$$
du lemme \ref{lemma-quasi-affine}. Notons que
$\Spec(A) = \text{Proj}(A[T])$, voir
Constructions, exemple \ref{constructions-example-trivial-proj}.
On peut donc appliquer le lemme \ref{lemma-open-in-proj-ample}
pour en déduire que $\mathcal{O}_X$ est ample.

\medskip\noindent
Supposons $\mathcal{O}_X$ ample.
Notons que $\Gamma_*(X, \mathcal{O}_X) \cong A[T]$
comme anneaux gradués. Le résultat découle donc des lemmes
\ref{lemma-ample-immersion-into-proj} et \ref{lemma-quasi-affine},
compte tenu de l'égalité
$\Spec(A) = \text{Proj}(A[T])$ pour tout anneau $A$,
établie ci-dessus.
\end{proof}

\begin{lemma}
\label{lemma-quasi-affine-locally-closed}
Soit $X$ un schéma quasi-affine. Pour toute immersion quasi-compacte
$i : X' \to X$, le schéma $X'$ est quasi-affine.
\end{lemma}

\begin{proof}
On peut le démontrer directement, sans faire appel aux résultats sur
les faisceaux inversibles amples ; nous invitons le lecteur à en esquisser une preuve.
Comme $X$ est quasi-affine, $\mathcal{O}_X$ est ample d'après le
lemme \ref{lemma-quasi-affine-O-ample}.
Alors $\mathcal{O}_{X'}$ est ample d'après le
lemme \ref{lemma-ample-on-locally-closed}. Donc $X'$ est quasi-affine d'après le
lemme \ref{lemma-quasi-affine-O-ample}.
\end{proof}

\begin{lemma}
\label{lemma-characterize-affine}
Soit $X$ un schéma. Supposons qu'il existe un nombre fini d'éléments
$f_1, \ldots, f_n \in \Gamma(X, \mathcal{O}_X)$ tels que
\begin{enumerate}
\item chaque $X_{f_i}$ soit un ouvert affine de $X$, et
\item l'idéal engendré par $f_1, \ldots, f_n$ dans
$\Gamma(X, \mathcal{O}_X)$ soit l'idéal unité.
\end{enumerate}
Alors $X$ est affine.
\end{lemma}

\begin{proof}
Supposons donnés $f_1, \ldots, f_n$ comme dans l'énoncé.
On peut écrire $1 = \sum g_i f_i$ avec $g_j \in \Gamma(X, \mathcal{O}_X)$,
d'où clairement $X = \bigcup X_{f_i}$. (Les $f_i$ ne peuvent
pas tous s'annuler en un même point.) Puisque chaque $X_{f_i}$
est quasi-compact (étant affine), $X$ est quasi-compact.
Il en résulte que $X$ est quasi-affine d'après le
lemme \ref{lemma-quasi-affine-O-ample} ci-dessus.
Considérons l'immersion ouverte
$$
j : X \to \Spec(\Gamma(X, \mathcal{O}_X)),
$$
voir le lemme \ref{lemma-quasi-affine}. L'image réciproque de l'ouvert standard
$D(f_i)$ du membre de droite est $X_{f_i}$ dans le
membre de gauche, et le morphisme $j$ induit un isomorphisme
$X_{f_i} \cong D(f_i)$, voir le
lemme \ref{lemma-invert-f-affine}. Puisque les $f_i$ engendrent l'idéal unité,
on a $\Spec(\Gamma(X, \mathcal{O}_X))
= \bigcup_{i = 1, \ldots, n} D(f_i)$. Ainsi $j$ est un isomorphisme.
\end{proof}


\section{Faisceaux quasi-cohérents et faisceaux inversibles amples}
\label{section-ample-quasi-coherent}

\noindent
Thème de cette section : en présence d'un faisceau inversible ample,
tout faisceau quasi-cohérent provient d'un module gradué.

\begin{situation}
\label{situation-ample}
Soit $X$ un schéma.
Soit $\mathcal{L}$ un faisceau inversible ample sur $X$.
Posons $S = \Gamma_*(X, \mathcal{L})$, considéré comme anneau gradué.
Posons $Y = \text{Proj}(S)$.
Soit $f : X \to Y$ le morphisme canonique du lemme \ref{lemma-map-into-proj}.
Il est muni d'un homomorphisme $\mathbf{Z}$-gradué de $\mathcal{O}_X$-algèbres
$\bigoplus f^*\mathcal{O}_Y(n) \to \bigoplus \mathcal{L}^{\otimes n}$.
\end{situation}

\noindent
Le lemme suivant est en fait un cas particulier de celui qui le suit,
mais il paraît utile de l'établir d'abord.

\begin{lemma}
\label{lemma-ample-gcd-is-one}
Plaçons-nous dans la situation \ref{situation-ample}.
Le morphisme canonique $f : X \to Y$
envoie $X$ dans le sous-schéma ouvert $W = W_1 \subset Y$
sur lequel $\mathcal{O}_Y(1)$ est inversible et toutes
les applications de multiplication
$\mathcal{O}_Y(n) \otimes_{\mathcal{O}_Y} \mathcal{O}_Y(m) \to
\mathcal{O}_Y(n + m)$
sont des isomorphismes (voir
Constructions, lemme \ref{constructions-lemma-where-invertible}).
De plus, les applications $f^*\mathcal{O}_Y(n) \to \mathcal{L}^{\otimes n}$
sont toutes des isomorphismes.
\end{lemma}

\begin{proof}
D'après la proposition \ref{proposition-characterize-ample}, il existe un entier
$n_0$ tel que $\mathcal{L}^{\otimes n}$ soit engendré par ses sections globales pour tout
$n \geq n_0$. Soit $x \in X$. D'après ce qui précède, on peut trouver
$a \in S_{n_0}$ et $b \in S_{n_0 + 1}$ tels que
$a$ et $b$ ne s'annulent pas en $x$. Ainsi
$f(x) \in D_{+}(a) \cap D_{+}(b) = D_{+}(ab)$. D'après
Constructions, lemme \ref{constructions-lemma-where-invertible},
on a $f(x) \in W_1$, comme voulu. D'après
Constructions, lemme \ref{constructions-lemma-invertible-map-into-proj},
utilisé dans la construction de l'application $f$,
les applications
$f^*\mathcal{O}_Y(n_0) \to \mathcal{L}^{\otimes n_0}$ et
$f^*\mathcal{O}_Y(n_0 + 1) \to \mathcal{L}^{\otimes n_0 + 1}$
sont des isomorphismes dans un voisinage de $x$. Par compatibilité avec
la structure d'algèbre et du fait que $f$ prend ses valeurs dans $W$,
on en conclut que toutes les applications
$f^*\mathcal{O}_Y(n) \to \mathcal{L}^{\otimes n}$ sont des isomorphismes
dans un voisinage de $x$, d'où la conclusion.
\end{proof}

\noindent
Rappelons d'après Modules, définition \ref{modules-definition-gamma-star},
qu'étant donnés un espace localement annelé $X$, un faisceau inversible $\mathcal{L}$
et un $\mathcal{O}_X$-module $\mathcal{F}$, on dispose du
$\Gamma_*(X, \mathcal{L})$-module gradué
$$
\Gamma_*(X, \mathcal{L}, \mathcal{F}) =
\bigoplus\nolimits_{n \in \mathbf{Z}}
\Gamma(X, \mathcal{F} \otimes_{\mathcal{O}_X} \mathcal{L}^{\otimes n}).
$$
Le lemme suivant affirme que, dans la situation \ref{situation-ample},
on peut retrouver un $\mathcal{O}_X$-module quasi-cohérent $\mathcal{F}$
à partir de ce module gradué. Voir aussi
Constructions, lemme \ref{constructions-lemma-quasi-coherent-projective-space},
où ce lemme est démontré dans le cas particulier $X = \mathbf{P}^n_R$.

\begin{lemma}
\label{lemma-ample-quasi-coherent}
Plaçons-nous dans la situation \ref{situation-ample}.
Soit $\mathcal{F}$ un faisceau quasi-cohérent sur $X$.
Posons $M = \Gamma_*(X, \mathcal{L}, \mathcal{F})$, considéré comme $S$-module gradué.
Il existe des isomorphismes
$$
f^*\widetilde{M} \longrightarrow \mathcal{F}
$$
fonctoriels en $\mathcal{F}$, tels que
$M_0 \to \Gamma(\text{Proj}(S), \widetilde{M}) \to \Gamma(X, \mathcal{F})$
soit l'application identique.
\end{lemma}

\begin{proof}
Soit $s \in S_{+}$ homogène tel que $X_s$ soit un ouvert affine de $X$.
Rappelons que $\widetilde{M}|_{D_{+}(s)}$ correspond au
$S_{(s)}$-module $M_{(s)}$, voir
Constructions, lemme \ref{constructions-lemma-proj-sheaves}.
Rappelons que $f^{-1}(D_{+}(s)) = X_s$.
Puisque $X$ possède un faisceau inversible ample, il est quasi-compact et
quasi-séparé, voir la section \ref{section-ample}.
D'après le lemme \ref{lemma-invert-s-sections}, il existe un isomorphisme canonique
$M_{(s)} = \Gamma_*(X, \mathcal{L}, \mathcal{F})_{(s)} \to
\Gamma(X_s, \mathcal{F})$.
Puisque $\mathcal{F}$ est quasi-cohérent, cela donne
un isomorphisme canonique
$$
f^*\widetilde{M}|_{X_s} \to \mathcal{F}|_{X_s}
$$
Comme $\mathcal{L}$ est ample sur $X$, on sait que $X$ est recouvert
par les ouverts affines de la forme $X_s$. Il suffit donc de montrer
que les applications ci-dessus se recollent sur les intersections. Nous en omettons
la démonstration.
\end{proof}

\begin{remark}
\label{remark-neurotic}
Sous les hypothèses et avec les notations du lemme \ref{lemma-ample-quasi-coherent},
notons $\theta_\mathcal{F}$ l'application qui y figure.
Notons que l'isomorphisme $f^*\mathcal{O}_Y(n) \to \mathcal{L}^{\otimes n}$
du lemme \ref{lemma-ample-gcd-is-one} n'est autre que
$\theta_{\mathcal{L}^{\otimes n}}$.
Considérons les applications de multiplication
$$
\widetilde{M} \otimes_{\mathcal{O}_Y} \mathcal{O}_Y(n)
\longrightarrow
\widetilde{M(n)}
$$
voir
Constructions, équation (\ref{constructions-equation-multiply-more-generally}).
Prenons-en l'image réciproque sur $X$ et considérons
$$
\xymatrix{
f^*\widetilde{M} \otimes_{\mathcal{O}_X} f^*\mathcal{O}_Y(n)
\ar[r]
\ar[d]_{\theta_\mathcal{F} \otimes \theta_{\mathcal{L}^{\otimes n}}}
&
f^*\widetilde{M(n)}
\ar[d]^{\theta_{\mathcal{F} \otimes \mathcal{L}^{\otimes n}}}
\\
\mathcal{F} \otimes \mathcal{L}^{\otimes n} \ar[r]^{\text{id}} &
\mathcal{F} \otimes \mathcal{L}^{\otimes n}
}
$$
On a utilisé ici l'identification évidente
$M(n) = \Gamma_*(X, \mathcal{L}, \mathcal{F} \otimes \mathcal{L}^{\otimes n})$.
Ce diagramme est commutatif. Démonstration omise.
\end{remark}

\noindent
On devrait pouvoir déduire le lemme suivant du
lemme \ref{lemma-ample-quasi-coherent} (ou réciproquement), mais il paraît
plus simple d'en reprendre la démonstration.

\begin{lemma}
\label{lemma-proj-quasi-coherent}
Soit $S$ un anneau gradué tel que $X = \text{Proj}(S)$ soit quasi-compact.
Soit $\mathcal{F}$ un $\mathcal{O}_X$-module quasi-cohérent. Posons
$M = \bigoplus_{n \in \mathbf{Z}} \Gamma(X, \mathcal{F}(n))$,
considéré comme $S$-module gradué, voir
Constructions, section \ref{constructions-section-invertible-on-proj}.
L'application
$$
\widetilde{M} \longrightarrow \mathcal{F}
$$
de Constructions, lemme
\ref{constructions-lemma-comparison-proj-quasi-coherent},
est un isomorphisme.
Si $X$ est recouvert par des ouverts standards $D_+(f)$, où $f$ est de degré $1$,
alors les applications induites
$M_n \to \Gamma(X, \mathcal{F}(n))$ sont les applications identiques.
\end{lemma}

\begin{proof}
Comme $X$ est quasi-compact, on peut trouver des éléments homogènes
$f_1, \ldots, f_n \in S$ de degrés strictement positifs tels que
$X = D_+(f_1) \cup \ldots \cup D_+(f_n)$. Soit $d$ le
plus petit commun multiple des degrés de $f_1, \ldots, f_n$.
En remplaçant chaque $f_i$ par une puissance, on peut supposer que chaque
$f_i$ est de degré $d$. Alors $\mathcal{L} = \mathcal{O}_X(d)$ est
inversible, les applications de multiplication
$\mathcal{O}_X(ad) \otimes \mathcal{O}_X(bd) \to \mathcal{O}_X((a + b)d)$
sont des isomorphismes, et chaque $f_i$ détermine une section globale $s_i$
de $\mathcal{L}$ telle que $X_{s_i} = D_+(f_i)$, voir
Constructions, lemmes \ref{constructions-lemma-where-invertible} et
\ref{constructions-lemma-principal-open}.
Ainsi $\Gamma(X, \mathcal{F}(ad)) =
\Gamma(X, \mathcal{F} \otimes \mathcal{L}^{\otimes a})$.
Rappelons que $\widetilde{M}|_{D_{+}(f_i)}$ correspond au
$S_{(f_i)}$-module $M_{(f_i)}$, voir
Constructions, lemme \ref{constructions-lemma-proj-sheaves}.
Comme le degré de $f_i$ est $d$, la classe d'isomorphisme de
$M_{(f_i)}$ ne dépend que des composantes homogènes de $M$ de
degré divisible par $d$. Plus précisément, la classe d'isomorphisme de
$M_{(f_i)}$ ne dépend que du $\Gamma_*(X, \mathcal{L})$-module gradué
$\Gamma_*(X, \mathcal{L}, \mathcal{F})$
et de l'image $s_i$ de $f_i$ dans $\Gamma_*(X, \mathcal{L})$.
Le schéma $X$ est quasi-compact par hypothèse et
séparé d'après Constructions, lemme \ref{constructions-lemma-proj-separated}.
D'après le lemme \ref{lemma-invert-s-sections}, il existe un isomorphisme canonique
$$
M_{(f_i)} = \Gamma_*(X, \mathcal{L}, \mathcal{F})_{(s_i)} \to
\Gamma(X_{s_i}, \mathcal{F}).
$$
La construction de l'application de Constructions, lemme
\ref{constructions-lemma-comparison-proj-quasi-coherent},
montre alors qu'elle est un isomorphisme sur $D_+(f_i)$,
donc un isomorphisme puisque ces ouverts recouvrent $X$.
Nous omettons la démonstration de la dernière assertion.
\end{proof}

\section{Recherche d'ouverts affines convenables}
\label{section-finding-affine-opens}

\noindent
Dans cette section, nous rassemblons quelques résultats sur l'existence
d'ouverts affines dans des situations plus ou moins générales.

\begin{lemma}
\label{lemma-maximal-points-affine}
Soit $X$ un schéma quasi-séparé.
Soient $Z_1, \ldots, Z_n$ des composantes irréductibles deux à deux distinctes de $X$,
voir Topologie, section \ref{topology-section-irreducible-components}.
Soient $\eta_i \in Z_i$ leurs points génériques, voir
Schémas, lemme \ref{schemes-lemma-scheme-sober}.
Il existe des voisinages ouverts affines $\eta_i \in U_i$
tels que $U_i \cap U_j = \emptyset$ pour tous $i \not = j$.
En particulier, $U = U_1 \cup \ldots \cup U_n$ est un ouvert
affine contenant tous les points $\eta_1, \ldots, \eta_n$.
\end{lemma}

\begin{proof}
Soit $V_i$ un ouvert affine quelconque contenant $\eta_i$
et disjoint du fermé $Z_1 \cup \ldots \hat Z_i \ldots \cup Z_n$.
Comme $X$ est quasi-séparé, pour chaque $i$, la réunion
$W_i = \bigcup_{j, j \not = i} V_i \cap V_j$ est un ouvert quasi-compact
de $V_i$ ne contenant pas $\eta_i$.
On peut trouver des voisinages ouverts $U_i \subset V_i$
contenant $\eta_i$ et disjoints de $W_i$ d'après
Algèbre, lemme \ref{algebra-lemma-standard-open-containing-maximal-point}.
Enfin, $U$ est affine, car c'est le spectre
de l'anneau $R_1 \times \ldots \times R_n$, où $R_i = \mathcal{O}_X(U_i)$,
voir Schémas, lemme \ref{schemes-lemma-disjoint-union-affines}.
\end{proof}

\begin{remark}
\label{remark-maximal-points-affine}
Le lemme \ref{lemma-maximal-points-affine} ci-dessus est faux si $X$
n'est pas quasi-séparé. En voici un exemple. Prenons
$R = \mathbf{Q}[x, y_1, y_2, \ldots]/((x-i)y_i)$.
Considérons l'idéal premier minimal $\mathfrak p = (y_1, y_2, \ldots)$
de $R$. Recollons deux copies de $\Spec(R)$ le long de
l'ouvert (non quasi-compact) $\Spec(R) \setminus V(\mathfrak p)$
pour obtenir un schéma $X$ (par le recollement décrit dans
Schémas, exemple \ref{schemes-example-affine-space-zero-doubled}).
Alors les deux points maximaux de $X$ correspondant à $\mathfrak p$
ne sont contenus dans aucun ouvert affine commun. En effet,
tout ouvert de $\Spec(R)$ contenant $\mathfrak p$
contient une infinité des « droites » $x = i$, $y_j = 0$,
$j \not = i$, de paramètre $y_i$. Détails omis.
\end{remark}

\noindent
Malgré l'exemple ci-dessus, pour « la plupart » des familles finies de
fermés irréductibles, on peut appliquer le lemme \ref{lemma-maximal-points-affine} ci-dessus,
du moins si $X$ est quasi-compact. Cela tient au fait que $X$ contient un
ouvert dense qui est séparé.

\begin{lemma}
\label{lemma-quasi-compact-dense-open-separated}
Soit $X$ un schéma quasi-compact.
Il existe un ouvert dense $V \subset X$ qui est séparé.
\end{lemma}

\begin{proof}
Écrivons $X = \bigcup_{i = 1, \ldots, n} U_i$ comme réunion de $n$ sous-schémas
ouverts affines. Nous allons démontrer le lemme par récurrence sur $n$. Il est immédiat pour
$n = 1$. Soit $V' \subset \bigcup_{i = 1, \ldots, n - 1} U_i$ un sous-schéma ouvert
dense séparé, dont l'existence résulte de l'hypothèse de récurrence. Considérons
$$
V = V' \amalg (U_n \setminus \overline{V'}).
$$
Il est clair que $V$ est séparé et est un sous-schéma ouvert dense de $X$.
\end{proof}

\noindent
Même si $X$ est à la fois quasi-séparé et quasi-compact,
il n'existe pas nécessairement d'ouvert dense séparé et quasi-compact, voir
Exemples, lemme
\ref{examples-lemma-no-dense-separated-quasi-compact-open-in-qcqs}.
Voici une légère précision du lemme \ref{lemma-maximal-points-affine} ci-dessus.

\begin{lemma}
\label{lemma-point-and-maximal-points-affine}
Soit $X$ un schéma quasi-séparé. Soient $Z_1, \ldots, Z_n$ des composantes
irréductibles deux à deux distinctes de $X$. Soient $\eta_i \in Z_i$ leurs
points génériques. Soit $x \in X$ un point quelconque.
Il existe un ouvert affine $U \subset X$ contenant
$x$ et tous les $\eta_i$.
\end{lemma}

\begin{proof}
Supposons que $x \in Z_1 \cap \ldots \cap Z_r$ et que
$x \not \in Z_{r + 1}, \ldots, Z_n$. On peut alors choisir
un ouvert affine $W \subset X$ tel que
$x \in W$ et $W \cap Z_i = \emptyset$ pour
$i = r + 1, \ldots, n$. Il est clair que $\eta_i \in W$
pour $i = 1, \ldots, r$. D'après le lemme \ref{lemma-maximal-points-affine},
on peut choisir des ouverts affines $U_i \subset X$ deux à deux
disjoints tels que $\eta_i \in U_i$ pour $i = r + 1, \ldots, n$.
Comme $X$ est quasi-séparé, les ouverts $W \cap U_i$
sont quasi-compacts et ne contiennent pas $\eta_i$ pour
$i = r + 1, \ldots, n$. Ainsi, d'après
Algèbre, lemme \ref{algebra-lemma-standard-open-containing-maximal-point},
on peut rétrécir $U_i$ de sorte que $W \cap U_i = \emptyset$
pour $i = r + 1, \ldots, n$. Alors la réunion
$U = W \cup \bigcup_{i = r + 1, \ldots, n} U_i$ est disjointe, donc
(d'après Schémas, lemme \ref{schemes-lemma-disjoint-union-affines})
c'est un ouvert affine convenable.
\end{proof}

\begin{lemma}
\label{lemma-ample-finite-set-in-affine}
Soit $X$ un schéma. Supposons vérifiée l'une des conditions suivantes :
\begin{enumerate}
\item Le schéma $X$ est quasi-affine.
\item Le schéma $X$ est isomorphe à un sous-schéma localement fermé
d'un schéma affine.
\item Il existe un faisceau inversible ample sur $X$.
\item Le schéma $X$ est isomorphe à un sous-schéma localement fermé
de $\text{Proj}(S)$ pour un anneau gradué $S$.
\end{enumerate}
Alors, pour toute partie finie $E \subset X$, il existe un
ouvert affine $U \subset X$ tel que $E \subset U$.
\end{lemma}

\begin{proof}
D'après Propriétés, définition \ref{definition-quasi-affine},
un schéma quasi-affine est un sous-schéma ouvert quasi-compact
d'un schéma affine. Tout schéma affine $\Spec(R)$ est isomorphe à
$\text{Proj}(R[X])$, où $R[X]$ est gradué en posant $\deg(X) = 1$.
D'après la proposition \ref{proposition-characterize-ample},
si $X$ possède un faisceau inversible ample, alors $X$ est isomorphe à un
sous-schéma ouvert de $\text{Proj}(S)$ pour un anneau gradué $S$.
Il suffit donc de démontrer le lemme dans le cas (4).
(Nous invitons le lecteur à démontrer directement le cas (2).)

\medskip\noindent
Supposons donc que $X \subset \text{Proj}(S)$ soit un sous-schéma localement fermé,
où $S$ est un anneau gradué. Soit $T = \overline{X} \setminus X$.
Rappelons que les ouverts standards $D_{+}(f)$ forment une base de la
topologie de $\text{Proj}(S)$. Comme $E$ est fini, on peut choisir un nombre fini
d'éléments homogènes $f_i \in S_{+}$ tels que
$$
E \subset
D_{+}(f_1) \cup \ldots \cup D_{+}(f_n) \subset
\text{Proj}(S) \setminus T
$$
Supposons que $E = \{\mathfrak p_1, \ldots, \mathfrak p_m\}$,
considéré comme partie de $\text{Proj}(S)$.
Considérons l'idéal $I = (f_1, \ldots, f_n) \subset S$.
Puisque $I \not \subset \mathfrak p_j$ pour tous $j = 1, \ldots, m$,
Algèbre, lemme \ref{algebra-lemma-graded-silly}, montre qu'il
existe un élément homogène $f \in I$, avec $f \not \in \mathfrak p_j$
pour tous $j = 1, \ldots, m$. Alors $E \subset D_{+}(f) \subset
D_{+}(f_1) \cup \ldots \cup D_{+}(f_n)$. Comme $D_{+}(f)$ ne
rencontre pas $T$, on voit que $X \cap D_{+}(f)$ est un sous-schéma fermé du
schéma affine $D_{+}(f)$, donc un ouvert affine de $X$, comme voulu.
\end{proof}

\begin{lemma}
\label{lemma-ample-finite-set-in-principal-affine}
Soit $X$ un schéma. Soit $\mathcal{L}$ un faisceau inversible ample sur $X$.
Soient
$$
E \subset W \subset X
$$
avec $E$ fini et $W$ ouvert dans $X$. Alors il existe un entier $n > 0$
et une section $s \in \Gamma(X, \mathcal{L}^{\otimes n})$ tels que
$X_s$ soit affine et que $E \subset X_s \subset W$.
\end{lemma}

\begin{proof}
Le lecteur peut adapter la démonstration du lemme \ref{lemma-ample-finite-set-in-affine}
pour établir celui-ci ; nous allons plutôt l'en déduire.
D'après le lemme \ref{lemma-ample-finite-set-in-affine}, on peut choisir un
ouvert affine $U \subset W$ tel que $E \subset U$.
Considérons l'anneau gradué $S = \Gamma_*(X, \mathcal{L}) =
\bigoplus_{n \geq 0} \Gamma(X, \mathcal{L}^{\otimes n})$.
Pour chaque $x \in E$, soit $\mathfrak p_x \subset S$ l'idéal gradué
des sections qui s'annulent en $x$. Il est clair que $\mathfrak p_x$ est
un idéal premier et, puisqu'une puissance de $\mathcal{L}$ est engendrée par
ses sections globales, il est clair que $S_{+} \not \subset \mathfrak p_x$.
Soit $I \subset S$ l'idéal gradué des sections qui s'annulent en tous
les points de $X \setminus U$. Comme les ensembles $X_s$ forment une base
de la topologie, on a $I \not \subset \mathfrak p_x$ pour
tous $x \in E$.
Par le lemme d'évitement des idéaux premiers dans le cas gradué (Algèbre, lemme \ref{algebra-lemma-graded-silly}),
on peut trouver $s \in I$ homogène
tel que $s \not \in \mathfrak p_x$ pour tous $x \in E$.
Alors $E \subset X_s \subset U$ et $X_s$ est affine d'après le
lemme \ref{lemma-affine-cap-s-open}.
\end{proof}

\begin{lemma}
\label{lemma-quasi-affine-invertible-nonvanishing-section}
Soit $X$ un schéma quasi-affine. Soit $\mathcal{L}$ un
$\mathcal{O}_X$-module inversible. Soient $E \subset W \subset X$, avec $E$ fini
et $W$ ouvert. Alors il existe $s \in \Gamma(X, \mathcal{L})$
tel que $X_s$ soit affine et que $E \subset X_s \subset W$.
\end{lemma}

\begin{proof}
La démonstration de ce lemme présente de nombreux points communs avec celle donnée dans
Algèbre, lemme \ref{algebra-lemma-silly}.
Écrivons $E = \{x_1, \ldots, x_n\}$. Si $E = W = \emptyset$, alors $s = 0$
convient. Si $W \not = \emptyset$, on peut supposer $E \not = \emptyset$
en ajoutant un point au besoin. On peut donc supposer $n \geq 1$.
Nous allons démontrer le lemme par récurrence sur $n$.

\medskip\noindent
Cas initial : $n = 1$. En remplaçant $W$ par un voisinage ouvert affine
de $x_1$ dans $W$, on peut supposer $W$ affine. En combinant
le lemme \ref{lemma-quasi-affine-O-ample} et la
proposition \ref{proposition-characterize-ample},
on voit que tout
$\mathcal{O}_X$-module quasi-cohérent est engendré par ses sections globales.
Il existe donc
une section globale $s$ de $\mathcal{L}$ qui ne s'annule pas en $x_1$.
D'autre part, soit $Z \subset X$ le
sous-schéma fermé réduit induit sur $X \setminus W$.
En appliquant la propriété d'engendrement par les sections globales au faisceau quasi-cohérent d'idéaux
$\mathcal{I}$ de $Z$, on trouve une section globale $f$ de $\mathcal{I}$
qui ne s'annule pas en $x_1$. Alors $s' = fs$ est une section globale
de $\mathcal{L}$ qui ne s'annule pas en $x_1$ et telle que
$X_{s'} \subset W$. Alors $X_{s'}$ est affine d'après le
lemme \ref{lemma-affine-cap-s-open}.

\medskip\noindent
Étape de récurrence pour $n > 1$. S'il existe une spécialisation
$x_i \leadsto x_j$ pour $i \not = j$, il suffit
de démontrer le lemme pour $\{x_1, \ldots, x_n\} \setminus \{x_i\}$,
et l'on conclut par récurrence. On peut donc supposer qu'il n'existe aucune
spécialisation entre les $x_i$.
D'après le lemme \ref{lemma-ample-finite-set-in-affine} ou le
lemme \ref{lemma-ample-finite-set-in-principal-affine},
on peut supposer $W$ affine.
Par récurrence, on peut trouver une section globale
$s$ de $\mathcal{L}$ telle que $X_s \subset W$ soit affine et contienne
$x_1, \ldots, x_{n - 1}$. Si $x_n \in X_s$, la démonstration est terminée.
Supposons que $s$ s'annule en $x_n$. D'après le cas $n = 1$, on peut trouver
une section globale $s'$ de $\mathcal{L}$ telle que
$\{x_n\} \subset X_{s'} \subset
W \setminus \overline{\{x_1, \ldots, x_{n - 1}\}}$.
On utilise ici le fait que $x_n$ n'est une spécialisation d'aucun des $x_1, \ldots, x_{n - 1}$.
Alors $s + s'$
est une section globale de $\mathcal{L}$ qui ne s'annule pas
en $x_1, \ldots, x_n$, avec $X_{s + s'} \subset W$,
et l'on conclut comme précédemment.
\end{proof}

\begin{lemma}
\label{lemma-ring-affine-open-injective-local-ring}
Soient $X$ un schéma et $x \in X$ un point. Il existe un voisinage ouvert
affine $U \subset X$ de $x$ tel que l'application canonique
$\mathcal{O}_X(U) \to \mathcal{O}_{X, x}$ soit injective dans chacun
des cas suivants :
\begin{enumerate}
\item $X$ est intègre,
\item $X$ est localement noethérien,
\item $X$ est réduit et possède un nombre fini de composantes irréductibles.
\end{enumerate}
\end{lemma}

\begin{proof}
Après traduction en termes d'algèbre, on applique
Algèbre, lemme \ref{algebra-lemma-subring-of-local-ring}.
\end{proof}

\begin{lemma}
\label{lemma-standard-open-two-affines}
Soient $U$, $V$ des schémas affines et soient $W \to U$ et $W \to V$
des immersions ouvertes. Pour tout $w \in W$, il existe un voisinage ouvert
affine $W' \subset W$ de $w$ tel que $W'$
ait pour image un ouvert standard de $U$ et de $V$.
\end{lemma}

\begin{proof}
Soit $X$ le schéma obtenu en recollant $U$ et $V$ le long de $W$ ;
appliquons Schémas, lemme \ref{schemes-lemma-standard-open-two-affines}.
\end{proof}






\begin{multicols}{2}[\section{Autres chapitres}]
\noindent
Préliminaires
\begin{enumerate}
\item \hyperref[introduction-section-phantom]{Introduction}
\item \hyperref[conventions-section-phantom]{Conventions}
\item \hyperref[sets-section-phantom]{Théorie des ensembles}
\item \hyperref[categories-section-phantom]{Catégories}
\item \hyperref[topology-section-phantom]{Topologie}
\item \hyperref[sheaves-section-phantom]{Faisceaux sur les espaces}
\item \hyperref[sites-section-phantom]{Sites et faisceaux}
\item \hyperref[stacks-section-phantom]{Champs}
\item \hyperref[fields-section-phantom]{Corps}
\item \hyperref[algebra-section-phantom]{Algèbre commutative}
\item \hyperref[brauer-section-phantom]{Groupes de Brauer}
\item \hyperref[homology-section-phantom]{Algèbre homologique}
\item \hyperref[derived-section-phantom]{Catégories dérivées}
\item \hyperref[simplicial-section-phantom]{Méthodes simpliciales}
\item \hyperref[more-algebra-section-phantom]{Compléments d'algèbre}
\item \hyperref[smoothing-section-phantom]{Lissage des morphismes d'anneaux}
\item \hyperref[modules-section-phantom]{Faisceaux de modules}
\item \hyperref[sites-modules-section-phantom]{Modules sur les sites}
\item \hyperref[injectives-section-phantom]{Objets injectifs}
\item \hyperref[cohomology-section-phantom]{Cohomologie des faisceaux}
\item \hyperref[sites-cohomology-section-phantom]{Cohomologie sur les sites}
\item \hyperref[dga-section-phantom]{Algèbre différentielle graduée}
\item \hyperref[dpa-section-phantom]{Algèbre à puissances divisées}
\item \hyperref[sdga-section-phantom]{Faisceaux différentiels gradués}
\item \hyperref[hypercovering-section-phantom]{Hyperrecouvrements}
\end{enumerate}
Schémas
\begin{enumerate}
\setcounter{enumi}{25}
\item \hyperref[schemes-section-phantom]{Schémas}
\item \hyperref[constructions-section-phantom]{Constructions de schémas}
\item \hyperref[properties-section-phantom]{Propriétés des schémas}
\item \hyperref[morphisms-section-phantom]{Morphismes de schémas}
\item \hyperref[coherent-section-phantom]{Cohomologie des schémas}
\item \hyperref[divisors-section-phantom]{Diviseurs}
\item \hyperref[limits-section-phantom]{Limites de schémas}
\item \hyperref[varieties-section-phantom]{Variétés}
\item \hyperref[topologies-section-phantom]{Topologies sur les schémas}
\item \hyperref[descent-section-phantom]{Descente}
\item \hyperref[perfect-section-phantom]{Catégories dérivées de schémas}
\item \hyperref[more-morphisms-section-phantom]{Compléments sur les morphismes}
\item \hyperref[flat-section-phantom]{Compléments sur la platitude}
\item \hyperref[groupoids-section-phantom]{Schémas en groupoïdes}
\item \hyperref[more-groupoids-section-phantom]{Compléments sur les schémas en groupoïdes}
\item \hyperref[etale-section-phantom]{Morphismes étales de schémas}
\end{enumerate}
Thèmes de théorie des schémas
\begin{enumerate}
\setcounter{enumi}{41}
\item \hyperref[chow-section-phantom]{Homologie de Chow}
\item \hyperref[intersection-section-phantom]{Théorie de l'intersection}
\item \hyperref[pic-section-phantom]{Schémas de Picard des courbes}
\item \hyperref[weil-section-phantom]{Théories cohomologiques de Weil}
\item \hyperref[adequate-section-phantom]{Modules adéquats}
\item \hyperref[dualizing-section-phantom]{Complexes dualisants}
\item \hyperref[duality-section-phantom]{Dualité pour les schémas}
\item \hyperref[discriminant-section-phantom]{Discriminants et différentes}
\item \hyperref[derham-section-phantom]{Cohomologie de de Rham}
\item \hyperref[local-cohomology-section-phantom]{Cohomologie locale}
\item \hyperref[algebraization-section-phantom]{Géométrie algébrique et formelle}
\item \hyperref[curves-section-phantom]{Courbes algébriques}
\item \hyperref[resolve-section-phantom]{Résolution des surfaces}
\item \hyperref[models-section-phantom]{Réduction semi-stable}
\item \hyperref[functors-section-phantom]{Foncteurs et morphismes}
\item \hyperref[equiv-section-phantom]{Catégories dérivées de variétés}
\item \hyperref[pione-section-phantom]{Groupes fondamentaux de schémas}
\item \hyperref[etale-cohomology-section-phantom]{Cohomologie étale}
\item \hyperref[crystalline-section-phantom]{Cohomologie cristalline}
\item \hyperref[proetale-section-phantom]{Cohomologie pro-étale}
\item \hyperref[relative-cycles-section-phantom]{Cycles relatifs}
\item \hyperref[more-etale-section-phantom]{Compléments de cohomologie étale}
\item \hyperref[trace-section-phantom]{La formule des traces}
\end{enumerate}
Espaces algébriques
\begin{enumerate}
\setcounter{enumi}{64}
\item \hyperref[spaces-section-phantom]{Espaces algébriques}
\item \hyperref[spaces-properties-section-phantom]{Propriétés des espaces algébriques}
\item \hyperref[spaces-morphisms-section-phantom]{Morphismes d'espaces algébriques}
\item \hyperref[decent-spaces-section-phantom]{Espaces algébriques décents}
\item \hyperref[spaces-cohomology-section-phantom]{Cohomologie des espaces algébriques}
\item \hyperref[spaces-limits-section-phantom]{Limites d'espaces algébriques}
\item \hyperref[spaces-divisors-section-phantom]{Diviseurs sur les espaces algébriques}
\item \hyperref[spaces-over-fields-section-phantom]{Espaces algébriques sur des corps}
\item \hyperref[spaces-topologies-section-phantom]{Topologies sur les espaces algébriques}
\item \hyperref[spaces-descent-section-phantom]{Descente et espaces algébriques}
\item \hyperref[spaces-perfect-section-phantom]{Catégories dérivées d'espaces}
\item \hyperref[spaces-more-morphisms-section-phantom]{Compléments sur les morphismes d'espaces}
\item \hyperref[spaces-flat-section-phantom]{Platitude sur les espaces algébriques}
\item \hyperref[spaces-groupoids-section-phantom]{Groupoïdes dans les espaces algébriques}
\item \hyperref[spaces-more-groupoids-section-phantom]{Compléments sur les groupoïdes dans les espaces}
\item \hyperref[bootstrap-section-phantom]{Amorçage}
\item \hyperref[spaces-pushouts-section-phantom]{Sommes amalgamées d'espaces algébriques}
\end{enumerate}
Thèmes de géométrie
\begin{enumerate}
\setcounter{enumi}{81}
\item \hyperref[spaces-chow-section-phantom]{Groupes de Chow des espaces}
\item \hyperref[groupoids-quotients-section-phantom]{Quotients de groupoïdes}
\item \hyperref[spaces-more-cohomology-section-phantom]{Compléments sur la cohomologie des espaces}
\item \hyperref[spaces-simplicial-section-phantom]{Espaces simpliciaux}
\item \hyperref[spaces-duality-section-phantom]{Dualité pour les espaces}
\item \hyperref[formal-spaces-section-phantom]{Espaces algébriques formels}
\item \hyperref[restricted-section-phantom]{Algébrisation des espaces formels}
\item \hyperref[spaces-resolve-section-phantom]{Résolution des surfaces revisitée}
\end{enumerate}
Théorie des déformations
\begin{enumerate}
\setcounter{enumi}{89}
\item \hyperref[formal-defos-section-phantom]{Théorie formelle des déformations}
\item \hyperref[defos-section-phantom]{Théorie des déformations}
\item \hyperref[cotangent-section-phantom]{Le complexe cotangent}
\item \hyperref[examples-defos-section-phantom]{Problèmes de déformation}
\end{enumerate}
Champs algébriques
\begin{enumerate}
\setcounter{enumi}{93}
\item \hyperref[algebraic-section-phantom]{Champs algébriques}
\item \hyperref[examples-stacks-section-phantom]{Exemples de champs}
\item \hyperref[stacks-sheaves-section-phantom]{Faisceaux sur les champs algébriques}
\item \hyperref[criteria-section-phantom]{Critères de représentabilité}
\item \hyperref[artin-section-phantom]{Axiomes d'Artin}
\item \hyperref[quot-section-phantom]{Espaces de Quot et de Hilbert}
\item \hyperref[stacks-properties-section-phantom]{Propriétés des champs algébriques}
\item \hyperref[stacks-morphisms-section-phantom]{Morphismes de champs algébriques}
\item \hyperref[stacks-limits-section-phantom]{Limites de champs algébriques}
\item \hyperref[stacks-cohomology-section-phantom]{Cohomologie des champs algébriques}
\item \hyperref[stacks-perfect-section-phantom]{Catégories dérivées de champs}
\item \hyperref[stacks-introduction-section-phantom]{Introduction aux champs algébriques}
\item \hyperref[stacks-more-morphisms-section-phantom]{Compléments sur les morphismes de champs}
\item \hyperref[stacks-geometry-section-phantom]{La géométrie des champs}
\end{enumerate}
Thèmes de théorie des espaces de modules
\begin{enumerate}
\setcounter{enumi}{107}
\item \hyperref[moduli-section-phantom]{Champs de modules}
\item \hyperref[moduli-curves-section-phantom]{Espaces de modules de courbes}
\end{enumerate}
Divers
\begin{enumerate}
\setcounter{enumi}{109}
\item \hyperref[examples-section-phantom]{Exemples}
\item \hyperref[exercises-section-phantom]{Exercices}
\item \hyperref[guide-section-phantom]{Guide bibliographique}
\item \hyperref[desirables-section-phantom]{Desiderata}
\item \hyperref[coding-section-phantom]{Style de programmation}
\item \hyperref[obsolete-section-phantom]{Obsolète}
\item \hyperref[fdl-section-phantom]{Licence de documentation libre GNU}
\item \hyperref[index-section-phantom]{Index généré automatiquement}
\end{enumerate}
\end{multicols}


\bibliography{my}
\bibliographystyle{amsalpha}

\end{document}
